\input{preamble}

% OK, start here.
%
\begin{document}

\title{Exercices}


\maketitle

\phantomsection
\label{section-phantom}

\tableofcontents


\section{Algèbre}
\label{section-algebra}

\noindent
Cette première section ne contient que quelques questions diverses.

\begin{exercise}
\label{exercise-isomorphism-localizations}
Soient $A$ un anneau et ${\mathfrak m}$ un idéal maximal. Dans $A[X]$,
posons $\tilde {\mathfrak m}_1 = ({\mathfrak m}, X)$ et
$\tilde {\mathfrak m}_2 = ({\mathfrak m}, X-1)$. Montrer
que
$$
A[X]_{\tilde {\mathfrak m}_1} \cong A[X]_{\tilde {\mathfrak m}_2}.
$$
\end{exercise}

\begin{exercise}
\label{exercise-coherent}
Trouver un exemple d'anneau non noethérien $R$ tel que tout idéal
de type fini de $R$ soit de présentation finie comme $R$-module.
(On dit qu'un anneau est {\it cohérent} si cette dernière propriété est vérifiée.)
\end{exercise}

\begin{exercise}
\label{exercise-flat-ideals-pid}
Supposons que $(A, {\mathfrak m}, k)$ soit un anneau local noethérien. Pour tout
$A$-module de type fini $M$, notons $r(M)$ le nombre minimal
de générateurs de $M$ comme $A$-module. D'après le lemme de Nakayama, ce nombre vaut
$\dim_k M/{\mathfrak m} M = \dim_k M \otimes_A k$.
\begin{enumerate}
\item Montrer que $r(M \otimes_A N) = r(M)r(N)$.
\item Soit $I\subset A $ un idéal tel que $r(I) > 1$. Montrer que
$r(I^2) < r(I)^2$.
\item En déduire que, si tout idéal de $A$ est un module plat, alors
$A$ est un anneau principal (ou un corps).
\end{enumerate}
\end{exercise}

\begin{exercise}
\label{exercise-not-isomorphic}
Soit $k$ un corps. Montrer que les paires suivantes de
$k$-algèbres ne sont pas isomorphes :
\begin{enumerate}
\item $k[x_1, \ldots, x_n]$ et $k[x_1, \ldots, x_{n + 1}]$ pour tout
$n\geq 1$.
\item $k[a, b, c, d, e, f]/(ab + cd + ef)$ et $k[x_1, \ldots, x_n]$
pour $n = 5$.
\item $k[a, b, c, d, e, f]/(ab + cd + ef)$ et $k[x_1, \ldots, x_n]$
pour $n = 6$.
\end{enumerate}
\end{exercise}

\begin{remark}
\label{remark-simple-geometric}
L'idée de cet exercice est bien entendu de trouver, dans chaque cas,
un argument simple plutôt que d'appliquer un « gros » théorème.
Il est néanmoins utile de se laisser guider par des principes généraux.
\end{remark}

\begin{exercise}
\label{exercise-silly}
Algèbre. (Exercice un peu futile, qui devrait être facile.)
\begin{enumerate}
\item Donner un exemple d'un anneau $A$ et d'une
suite exacte courte non scindée de $A$-modules
$$
0 \to M_1 \to M_2 \to M_3 \to 0.
$$
\item Donner un exemple d'une suite non scindée de $A$-modules
comme ci-dessus et d'un homomorphisme fidèlement plat $A \to B$ tel que
$$
0 \to M_1\otimes_A B \to M_2\otimes_A B \to M_3\otimes_A B \to 0.
$$
soit scindée comme suite de $B$-modules.
\end{enumerate}
\end{exercise}

\begin{exercise}
\label{exercise-field-kummer}
Supposons que le corps $k$ contienne une racine primitive $n$-ième
de l'unité $\zeta$. Cela signifie que $\zeta^n = 1$, mais que $\zeta^m\not = 1$ pour
$0 < m < n$.
\begin{enumerate}
\item Montrer que la caractéristique de $k$ est première à $n$.
\item Supposons que $a \in k$ soit un élément de $k$ qui ne soit une
puissance $d$-ième dans $k$ pour aucun diviseur $d$ de $n$ tel que $n \geq d > 1$. Montrer que
$k[x]/(x^n-a)$ est un corps. (Indication : considérer un corps de décomposition de
$x^n-a$ et utiliser la théorie de Galois.)
\end{enumerate}
\end{exercise}

\begin{exercise}
\label{exercise-valuation}
Soit $\nu : k[x]\setminus \{0\}  \to {\mathbf Z}$ une application
ayant les propriétés suivantes : $\nu(fg) = \nu(f) + \nu(g)$ lorsque
$f$, $g$ sont non nuls, et $\nu(f + g) \geq min(\nu(f), \nu(g))$ lorsque
$f$, $g$ et $f + g$ sont non nuls, et $\nu(c) = 0$ pour tout $c\in k^*$.
\begin{enumerate}
\item Montrer que, si $f$, $g$ et $f + g$ sont non nuls et si
$\nu(f) \not = \nu(g)$, alors $\nu(f + g) = min(\nu(f), \nu(g))$.
\item Montrer que, si $f = \sum a_i x^i$, $f\not = 0$, alors
$\nu(f) \geq min(\{i\nu(x)\}_{a_i\not = 0})$. Quand a-t-on égalité ?
\item Montrer que, si $\nu$ prend une valeur négative, alors
$\nu(f) = -n \deg(f)$ pour un certain $n\in {\mathbf N}$.
\item Supposons que $\nu(x) \geq 0$. Montrer que
$\{f \mid f = 0, \ \text{ou}\ \nu(f) > 0\}$ est un idéal premier de $k[x]$.
\item Décrire toutes les applications $\nu$ possibles.
\end{enumerate}
\end{exercise}


\noindent
Soit $A$ un anneau. Un {\it idempotent} est un élément $e \in A$
tel que $e^2 = e$. Les éléments $1$ et $0$ sont toujours idempotents.
Un {\it idempotent non trivial} est un idempotent distinct de zéro et de l'unité.
Deux idempotents $e, e' \in A$ sont dits {\it orthogonaux}
si $ee' = 0$.

\begin{exercise}
\label{exercise-product}
Soit $A$ un anneau. Montrer que $A$ est un produit de deux anneaux non nuls si
et seulement si $A$ possède un idempotent non trivial.
\end{exercise}

\begin{exercise}
\label{exercise-lift-idempotents}
Soient $A$ un anneau et $I \subset A$ un idéal localement nilpotent.
Montrer que l'application $A \to A/I$ induit une bijection sur les idempotents.
(Indication : il peut être plus facile de commencer par le cas où $I$ est nilpotent.
Utiliser ensuite la « réduction noethérienne absolue » pour se ramener à ce cas.)
\end{exercise}







% Fin de la tranche source louée.
\section{Limites inductives}
\label{section-colimits}


\begin{definition}
\label{definition-directed-poset}
Un {\it ensemble préordonné filtrant} est un ensemble non vide $I$ muni d'un préordre $\leq$
tel que, pour toute paire $i, j \in I$, il existe un $k \in I$ tel que
$i \leq k$ et $j \leq k$. Un {\it système inductif d'anneaux} indexé par $I$ est la donnée d'un
anneau $A_i$ pour chaque $i \in I$ et d'un homomorphisme d'anneaux $\varphi_{ij} : A_i \to A_j$
dès que $i \leq j$, de sorte que le composé $A_i \to A_j \to A_k$ soit égal à
$A_i \to A_k$ dès que $i \leq j \leq k$.
\end{definition}

\noindent
On définit de même les systèmes inductifs de groupes, de modules sur un anneau fixé,
d'espaces vectoriels sur un corps, etc.

\begin{exercise}
\label{exercise-directed-colimit}
Soit $I$ un ensemble préordonné filtrant et soit
$(A_i, \varphi_{ij})$ un système inductif d'anneaux indexé par $I$.
Montrer qu'il existe un anneau $A$ et des homomorphismes $\varphi_i : A_i \to A$
tels que $\varphi_j \circ \varphi_{ij} = \varphi_i$ pour tous $i \leq j$
et possédant la propriété universelle suivante : étant donnés un anneau $B$
et des homomorphismes $\psi_i : A_i \to B$ tels que
$\psi_j \circ \varphi_{ij} = \psi_i$ pour tous $i \leq j$, il
existe un unique homomorphisme d'anneaux $\psi : A \to B$ tel que
$\psi_i = \psi \circ \varphi_i$.
\end{exercise}

\begin{definition}
\label{definition-colimit}
L'anneau $A$ construit dans l'exercice \ref{exercise-directed-colimit}
est appelé la {\it limite inductive} du système. Notation : $\colim A_i$.
\end{definition}

\begin{exercise}
\label{exercise-prime-in-colimit}
Soit $(I, \geq)$ un ensemble préordonné filtrant et soit
$(A_i, \varphi_{ij})$ un système inductif d'anneaux indexé par $I$, de limite inductive
$A$. Montrer qu'il existe une bijection
$$
\Spec(A) = \{(\mathfrak p_i)_{i \in I} \mid
\mathfrak p_i \subset A_i \text{ et }
\mathfrak p_i = \varphi_{ij}^{-1}(\mathfrak p_j)\ \forall i \leq j\}
\subset \prod\nolimits_{i \in I} \Spec(A_i)
$$
L'ensemble figurant au membre de droite de l'égalité est la limite projective des ensembles
$\Spec(A_i)$. Notation : $\lim \Spec(A_i)$.
\end{exercise}

\begin{exercise}
\label{exercise-colimit-surjective}
Soit $(I, \geq)$ un ensemble préordonné filtrant et soit
$(A_i, \varphi_{ij})$ un système inductif d'anneaux indexé par $I$, de limite inductive
$A$. Supposons que l'application $\Spec(A_j) \to \Spec(A_i)$ soit
surjective pour tous $i \leq j$. Montrer que
$\Spec(A) \to \Spec(A_i)$ est surjective pour tout $i$.
(Indication : on peut essayer d'utiliser le théorème de Tychonoff, mais il existe aussi une
démonstration algébrique directe essentiellement triviale fondée sur
Algèbre, lemme \ref{algebra-lemma-in-image}.)
\end{exercise}

\begin{exercise}
\label{exercise-integral-colimit-finite}
Soit une extension entière d'anneaux $A \subset B$. Montrer que
$\Spec(B) \to \Spec(A)$ est surjective.
Utiliser les exercices ci-dessus, le fait que l'assertion est vraie pour une extension finie
d'anneaux (ce qui a été démontré dans les cours), et montrer que
$B = \colim B_i$ est une limite inductive filtrante d'extensions finies $A \subset B_i$.
\end{exercise}

\begin{exercise}
\label{exercise-colimit-tensor}
Soit $(I, \geq)$ un ensemble préordonné filtrant.
Soit $A$ un anneau et soit $(N_i, \varphi_{i, i'})$ un système inductif de
$A$-modules indexé par $I$. Soit en outre $M$ un $A$-module. Montrer
que
$$
\colim_{i\in I} M \otimes_A N_i\cong
M \otimes_A \Big( \colim_{i\in I} N_i\Big).
$$
\end{exercise}

\begin{definition}
\label{definition-finite-presentation}
Un module $M$ sur $R$ est dit de {\it présentation finie} sur
$R$ s'il est isomorphe au conoyau d'un homomorphisme de modules libres de type fini
$ R^{\oplus n} \to R^{\oplus m}$.
\end{definition}

\begin{exercise}
\label{exercise-colimit-modules}
Montrer que tout module sur un anneau quelconque est
\begin{enumerate}
\item la limite inductive de ses sous-modules de type fini, et
\item, d'une certaine manière, une limite inductive de modules de présentation finie.
\end{enumerate}
\end{exercise}





\section{Catégories additives et abéliennes}
\label{section-additive}

\begin{exercise}
\label{exercise-filtered-vector-spaces}
Soit $k$ un corps. Soit $\mathcal{C}$ la catégorie des espaces vectoriels filtrés
sur $k$ ; voir
Homologie, définition \ref{homology-definition-filtered}
pour la définition d'un objet filtré d'une catégorie quelconque.
\begin{enumerate}
\item Montrer que c'est une catégorie additive (expliquer soigneusement ce qu'est la
somme directe de deux objets).
\item Soit $f : (V, F) \to (W, F)$ un morphisme de $\mathcal{C}$.
Montrer que $f$ possède un noyau et un conoyau (expliquer précisément
ce que sont le noyau et le conoyau de $f$).
\item Donner un exemple d'un morphisme de $\mathcal{C}$ tel que
le morphisme canonique $\Coim(f) \to \Im(f)$ ne soit pas un isomorphisme.
\end{enumerate}
\end{exercise}

\begin{exercise}
\label{exercise-torsion-free}
Soit $R$ un anneau intègre noethérien. Soit $\mathcal{C}$ la catégorie des
$R$-modules sans torsion de type fini.
\begin{enumerate}
\item Montrer que c'est une catégorie additive.
\item Soit $f : N \to M$ un morphisme de $\mathcal{C}$.
Montrer que $f$ possède un noyau et un conoyau (veiller à définir précisément
ce que sont le noyau et le conoyau de $f$).
\item Donner un exemple d'un anneau intègre noethérien $R$ et d'un morphisme de
$\mathcal{C}$ tel que le morphisme canonique $\Coim(f) \to \Im(f)$
ne soit pas un isomorphisme.
\end{enumerate}
\end{exercise}

\begin{exercise}
\label{exercise-other}
Donner un exemple d'une catégorie additive qui possède des noyaux
et des conoyaux, mais qui soit autre que celles des
exercices \ref{exercise-filtered-vector-spaces} et
\ref{exercise-torsion-free}.
\end{exercise}





\section{Produit tensoriel}
\label{section-tensor-product}

\noindent
Les produits tensoriels sont introduits dans Algèbre, section
\ref{algebra-section-tensor-product}.
Soit $R$ un anneau. Soit $\text{Mod}_R$ la catégorie des $R$-modules.
Nous dirons qu'un foncteur $F : \text{Mod}_R \to \text{Mod}_R$
\begin{enumerate}
\item est additif si
$F : \Hom_R(M, N) \to \Hom_R(F(M), F(N))$
est un homomorphisme de groupes abéliens
pour tous $R$-modules $M, N$ ; voir
Homologie, définition \ref{homology-definition-preadditive} ;
\item est $R$-linéaire si $F : \Hom_R(M, N) \to \Hom_R(F(M), F(N))$ est $R$-linéaire
pour tous $R$-modules $M, N$ ;
\item est exact à droite si, pour toute suite exacte courte
$0 \to M_1 \to M_2 \to M_3 \to 0$, la suite
$F(M_1) \to F(M_2) \to F(M_3) \to 0$ est exacte ;
\item est exact à gauche si, pour toute suite exacte courte
$0 \to M_1 \to M_2 \to M_3 \to 0$, la suite
$0 \to F(M_1) \to F(M_2) \to F(M_3)$ est exacte ;
\item commute aux sommes directes si, étant donnés un ensemble
$I$ et des $R$-modules $M_i$, les homomorphismes
$F(M_i) \to F(\bigoplus M_i)$ induisent un isomorphisme
$\bigoplus F(M_i) = F(\bigoplus M_i)$.
\end{enumerate}

\begin{exercise}
\label{exercise-characterize-tensor-functor}
Soit $R$ un anneau. On conserve les notations ci-dessus.
\begin{enumerate}
\item Donner un exemple d'un anneau $R$ et d'un foncteur additif
$F : \text{Mod}_R \to \text{Mod}_R$ qui ne soit pas $R$-linéaire.
\item Soit $N$ un $R$-module. Montrer que le foncteur
$F(M) = M \otimes_R N$ est $R$-linéaire,
exact à droite et commute aux sommes directes.
\item Réciproquement, montrer que tout foncteur $F : \text{Mod}_R \to \text{Mod}_R$
qui est $R$-linéaire,
exact à droite et commute aux sommes directes est de la
forme $F(M) = M \otimes_R N$ pour un certain $R$-module $N$.
\item Montrer que, si dans (3) on supprime l'hypothèse
selon laquelle $F$ commute aux sommes directes, la conclusion n'est
plus vraie.
\end{enumerate}
\end{exercise}





\section{Homomorphismes plats d'anneaux}
\label{section-flat}

\begin{exercise}
\label{exercise-localization-flat}
Soit $S$ une partie multiplicative de l'anneau $A$.
\begin{enumerate}
\item Pour un $A$-module $M$, montrer que $S^{-1}M = S^{-1}A \otimes_A M$.
\item Montrer que $S^{-1}A$ est plat sur $A$.
\end{enumerate}
\end{exercise}

\begin{exercise}
\label{exercise-examples-not-flat}
Trouver une injection $M_1 \to M_2$ de $A$-modules telle que
$M_1\otimes N \to M_2 \otimes N$ ne soit pas injective dans les
cas suivants :
\begin{enumerate}
\item $A = k[x, y]$ et $N = (x, y) \subset A$. (Ici et ci-dessous, $k$ est un corps.)
\item $A = k[x, y]$ et $N = A/(x, y)$.
\end{enumerate}
\end{exercise}

\begin{exercise}
\label{exercise-flat-not-projective}
Donner un exemple d'un anneau $A$ et d'un $A$-module de type fini $M$
qui soit plat mais non projectif comme $A$-module.
\end{exercise}

\begin{remark}
\label{remark-flat-not-projective}
Si $M$ est de présentation finie et plat sur $A$,
alors $M$ est projectif sur $A$. L'exemple devra donc faire
intervenir un anneau $A$ qui ne soit pas noethérien. Je connais un exemple
où $A$ est l'anneau des fonctions de classe ${\mathcal C}^\infty$ sur ${\mathbf R}$.
\end{remark}

\begin{exercise}
\label{exercise-flat-not-free-dvr}
Trouver un module plat mais non libre sur ${\mathbf Z}_{(2)}$.
\end{exercise}

\begin{exercise}
\label{exercise-flat-deformations}
Déformations plates.
\begin{enumerate}
\item Supposons que $k$ soit un corps et que $k[\epsilon]$ soit l'algèbre des
nombres duaux $k[\epsilon] = k[x]/(x^2)$, avec $\epsilon = \bar x$. Montrer que, pour
toute $k$-algèbre $A$, il existe une $k[\epsilon]$-algèbre plate $B$ telle que
$A$ soit isomorphe à $B/\epsilon B$.
\item Supposons que $k = {\mathbf F}_p = {\mathbf Z}/p{\mathbf Z}$ et
$$
A = k[x_1, x_2, x_3, x_4, x_5, x_6]/
(x_1^p, x_2^p, x_3^p, x_4^p, x_5^p, x_6^p).
$$
Montrer qu'il existe une ${\mathbf Z}/p^2{\mathbf Z}$-algèbre plate $B$ telle
que $B/pB$ soit isomorphe à $A$. (Ainsi, $p$ joue ici le rôle de $\epsilon$.)
\item Posons maintenant $p = 2$ et considérons la même question pour
$k = {\mathbf F}_2 = {\mathbf Z}/2{\mathbf Z}$ et
$$
A = k[x_1, x_2, x_3, x_4, x_5, x_6]/
(x_1^2, x_2^2, x_3^2, x_4^2, x_5^2, x_6^2, x_1x_2 + x_3x_4 + x_5x_6).
$$
Cependant, montrer que dans ce cas il n'existe {\it pas} de
${\mathbf Z}/4{\mathbf Z}$-algèbre plate $B$ telle que $B/2B$ soit isomorphe à
$A$. (Trouver l'astuce ! Le même exemple fonctionne en toute caractéristique
$p > 0$, mais le calcul est plus difficile.)
\end{enumerate}
\end{exercise}

\begin{exercise}
\label{exercise-flat-given-residue-field-extension}
Soit $(A, {\mathfrak m}, k)$ un anneau local et soit $k'/k$
une extension finie de corps. Montrer qu'il existe un homomorphisme local et plat
d'anneaux locaux $A \to B$ tel que ${\mathfrak m}_B = {\mathfrak m} B$ et que
$B/{\mathfrak m} B$ soit
isomorphe à $k'$ comme $k$-algèbre. (Indication : commencer par le cas où
l'extension $k \subset k'$ est engendrée par un seul élément.)
\end{exercise}

\begin{remark}
\label{remark-flat-given-residue-field-extension-general}
Le même résultat vaut pour toute extension de corps $K/k$.
\end{remark}





% Fin de la tranche source assignée.




\section{Le spectre d'un anneau}
\label{section-spectrum-ring}

\begin{exercise}
\label{exercise-spec-Z}
Déterminer $\Spec(\mathbf{Z})$ en tant qu'ensemble et décrire sa topologie.
\end{exercise}

\begin{exercise}
\label{exercise-basis-opens-standard}
Soit $A$ un anneau quelconque. Pour $f\in A$, on définit
$D(f):= \{\mathfrak p \subset A \mid f \not \in \mathfrak p\}$.
Montrer que les ouverts $D(f)$ forment une base de la topologie de
$\Spec(A)$.
\end{exercise}

\begin{exercise}
\label{exercise-radical-ideals-closed}
Montrer que l'application $I\mapsto V(I)$
définit une bijection naturelle
$$
\{I\subset A\text{ tel que }I = \sqrt{I}\}
\longrightarrow
\{T\subset \Spec(A)\text{ fermée}\}
$$
\end{exercise}

\begin{definition}
\label{definition-quasi-compact}
Un espace topologique $X$ est dit {\it quasi-compact}
si, pour tout recouvrement ouvert $X = \bigcup_{i\in I} U_i$, il existe une
partie finie $\{i_1, \ldots, i_n\}\subset I$ telle que $X = U_{i_1}\cup\ldots
U_{i_n}$.
\end{definition}

\begin{exercise}
\label{exercise-spec-quasi-compact}
Montrer que $\Spec(A)$ est quasi-compact pour tout anneau $A$.
\end{exercise}

\begin{definition}
\label{definition-Hausdorff}
On dit qu'un espace topologique $X$ vérifie l'axiome de séparation $T_0$
si, pour tout couple de points $x, y\in X$ tels que $x\not = y$, il existe un
ouvert de $X$ qui contient l'un mais pas l'autre.
On dit que $X$ est {\it séparé} si, pour tout couple $x, y\in X$ tel que $x\not = y$,
il existe des ouverts disjoints $U, V$ tels que $x\in U$
et $y\in V$.
\end{definition}

\begin{exercise}
\label{exercise-not-hausdorff}
Montrer que $\Spec(A)$ n'est {\bf pas} séparé en général.
Montrer que $\Spec(A)$ est $T_0$. Donner un exemple d'espace topologique
$X$ qui n'est pas $T_0$.
\end{exercise}

\begin{remark}
\label{remark-not-hausdorff}
En général, on réserve le mot compact aux espaces quasi-compacts et
séparés.
\end{remark}

\begin{definition}
\label{definition-irreducible}
Un espace topologique $X$ est dit {\it irréductible} si $X$ est non vide
et si $X = Z_1\cup Z_2$, avec $Z_1, Z_2\subset X$ fermés, alors soit
$Z_1 = X$, soit $Z_2 = X$. Une partie $T\subset X$ d'un espace topologique
est dite {\it irréductible} si elle est un espace topologique
irréductible pour la topologie induite par $X$.
Cette définition implique que $T$ est irréductible si et seulement
si l'adhérence $\bar T$ de $T$ dans $X$ est irréductible.
\end{definition}

\begin{exercise}
\label{exercise-irreducible-spec}
Montrer que $\Spec(A)$ est irréductible si et seulement si
$Nil(A)$ est un idéal premier et que, dans ce cas, c'est l'unique
idéal premier minimal de $A$.
\end{exercise}

\begin{exercise}
\label{exercise-irreducible-prime}
Montrer qu'une partie fermée $T\subset \Spec(A)$
est irréductible si et seulement si elle est de la forme $T = V({\mathfrak p})$ pour
un idéal premier ${\mathfrak p}\subset A$.
\end{exercise}

\begin{definition}
\label{definition-generic-point}
Un point $x$ d'un espace topologique irréductible $X$ est appelé
un {\it point générique} de $X$ si $X$ est égal à l'adhérence de
la partie $\{x\}$.
\end{definition}

\begin{exercise}
\label{exercise-irreducible-T0-at-most-one-generic}
Montrer que, dans un espace $T_0$ $X$, toute partie fermée
irréductible admet au plus un point générique.
\end{exercise}

\begin{exercise}
\label{exercise-spec-sober}
Montrer que, dans $\Spec(A)$, toute
partie fermée irréductible {\it admet bien} un point générique.
Montrer en fait que l'application
${\mathfrak p} \mapsto \overline{\{{\mathfrak p}\}}$ est
une bijection de $\Spec(A)$ sur l'ensemble des parties fermées
irréductibles de $X$.
\end{exercise}

\begin{exercise}
\label{exercise-irreducible-subset-not-generic}
Donner un exemple montrant qu'une partie irréductible de
$\Spec(\mathbf{Z})$ ne possède pas nécessairement de point générique.
\end{exercise}

\begin{definition}
\label{definition-Noetherian-space}
Un espace topologique $X$ est dit {\it noethérien} si toute
suite décroissante $Z_1\supset Z_2 \supset Z_3\supset \ldots$
de parties fermées de $X$ est stationnaire.
(Il est dit {\it artinien} si toute suite croissante de parties
fermées est stationnaire.)
\end{definition}

\begin{exercise}
\label{exercise-Noetherian-spec}
Montrer que, si l'anneau $A$ est noethérien, alors
l'espace topologique $\Spec(A)$ est noethérien. Donner un
exemple montrant que la réciproque est fausse. (Même question dans le cas
artinien, si vous le souhaitez.)
\end{exercise}

\begin{definition}
\label{definition-irreducible-component}
Une partie irréductible maximale $T\subset X$ est appelée une
{\it composante irréductible} de l'espace $X$. Une telle composante
irréductible de $X$ est automatiquement une partie fermée de $X$.
\end{definition}

\begin{exercise}
\label{exercise-irreducible-in-irreducible}
Montrer que toute partie irréductible
de $X$ est contenue dans une composante irréductible de $X$.
\end{exercise}

\begin{exercise}
\label{exercise-Noetherian-finite-nr-irreducible}
Montrer qu'un espace topologique noethérien $X$
n'a qu'un nombre fini de composantes irréductibles, disons $X_1, \ldots, X_n$,
et que $X = X_1\cup X_2\cup\ldots\cup X_n$. (Noter que
tout $X$ est toujours la réunion de ses composantes irréductibles, mais que,
si $X = {\mathbf R}$ est par exemple muni de sa topologie usuelle, les composantes irréductibles
de $X$ sont les parties réduites à un point. Ce n'est guère
intéressant.)
\end{exercise}

\begin{exercise}
\label{exercise-irreducible-components-minimal-primes}
Montrer que les composantes irréductibles de $\Spec(A)$
correspondent aux idéaux premiers minimaux de $A$.
\end{exercise}

\begin{definition}
\label{definition-closed}
Un point $x\in X$ est dit {\it fermé} si $\overline{\{x\}} = \{ x\}$.
Soient $x, y$ des points de $X$. On dit que $x$ est une {\it spécialisation}
de $y$, ou que $y$ est une {\it générisation} de $x$, si
$x\in \overline{\{y\}}$.
\end{definition}

\begin{exercise}
\label{exercise-closed-maximal}
Montrer que les points fermés de $\Spec(A)$
correspondent aux idéaux maximaux de $A$.
\end{exercise}

\begin{exercise}
\label{exercise-generalization}
Montrer que ${\mathfrak p}$ est une générisation de ${\mathfrak q}$
dans $\Spec(A)$ si et seulement si ${\mathfrak p}\subset {\mathfrak q}$.
Caractériser les points fermés,
les idéaux maximaux, les points génériques et les idéaux premiers minimaux en termes de
générisation et de spécialisation. (On utilise ici la terminologie selon laquelle un point
d'un espace topologique éventuellement réductible $X$ est appelé point générique
s'il est un point générique de l'une des composantes irréductibles de $X$.)
\end{exercise}

\begin{exercise}
\label{exercise-disjoint-closed-spec}
Soient $I$ et $J$ des idéaux de $A$.
À quelle condition $V(I)$ et $V(J)$ sont-ils disjoints ?
\end{exercise}

\begin{definition}
\label{definition-connected-component}
Un espace topologique $X$ est dit {\it connexe} s'il est non vide et n'est pas la
réunion de deux ouverts non vides et disjoints. Une {\it composante connexe}
de $X$ est une partie connexe maximale. Tout point de $X$ est contenu
dans une composante connexe de $X$, et toute composante connexe de $X$ est
fermée dans $X$. (Mais, en général, une composante connexe n'est pas nécessairement ouverte dans $X$.)
\end{definition}

\begin{exercise}
\label{exercise-disconnected-spec}
Soit $A$ un anneau non nul.
Montrer que $\Spec(A)$ est non connexe
si et seulement si $A\cong B \times C$ pour certains anneaux non nuls $B, C$.
\end{exercise}

\begin{exercise}
\label{exercise-connected-component-stable-generalization}
Soit $T$ une composante connexe
de $\Spec(A)$. Montrer que $T$ est stable par générisation.
Montrer que $T$ est une partie ouverte de $\Spec(A)$ si $A$ est noethérien.
(Remarque : cette assertion est fausse lorsque $A$ est, par exemple, un produit infini de copies de
${\mathbf F}_2$. Le spectre de cet anneau est constitué d'une infinité
de points fermés.)
\end{exercise}

\begin{exercise}
\label{exercise-primes-kx}
Déterminer $\Spec(k[x])$, c'est-à-dire décrire
les idéaux premiers de cet anneau, décrire les spécialisations possibles et
décrire la topologie. (Traiter le cas où $k$ est algébriquement clos, mais
aussi celui où $k$ ne l'est pas.)
\end{exercise}

\begin{exercise}
\label{exercise-primes-kxy}
 Déterminer $\Spec(k[x, y])$, où $k$ est algébriquement
clos.
[Indication : utiliser le morphisme
$\varphi : \Spec(k[x, y]) \to \Spec(k[x])$ ; si
$\varphi({\mathfrak p}) = (0)$, alors localiser suivant
$S = \{f\in k[x] \mid f \not = 0\}$ et utiliser le résultat du cours
sur la localisation et $\Spec$.]
(Pourquoi, à votre avis, les géomètres algébristes appellent-ils cet espace l'espace affine de dimension 2 ?)
\end{exercise}

\begin{exercise}
\label{exercise-primes-Zy}
Déterminer $\Spec(\mathbf{Z}[y])$.
[Indication : procéder comme ci-dessus.] (Espace affine de dimension 1 sur $\mathbf{Z}$.)
\end{exercise}






\section{Localisation}
\label{section-localization}


\begin{exercise}
\label{exercise-submodule-localization}
Soit $A$ un anneau. Soit $S \subset A$ une partie multiplicative.
Soit $M$ un $A$-module. Soit $N \subset S^{-1}M$ un $S^{-1}A$-sous-module.
Montrer qu'il existe un $A$-sous-module $N' \subset M$ tel que
$N = S^{-1}N'$. (Ce résultat utile s'applique en particulier aux idéaux
de $S^{-1}A$.)
\end{exercise}

\begin{exercise}
\label{exercise-localize-zero}
Soit $A$ un anneau. Soit $M$ un $A$-module. Soit $m \in M$.
\begin{enumerate}
\item Montrer que $I = \{a \in A \mid am = 0\}$ est un idéal de $A$.
\item Pour un idéal premier $\mathfrak p$ de $A$, montrer que l'image de $m$
dans $M_\mathfrak p$ est nulle si et seulement si $I \not \subset \mathfrak p$.
\item Montrer que $m$ est nul si et seulement si l'image de $m$ est nulle
dans $M_\mathfrak p$ pour tout idéal premier $\mathfrak p$ de $A$.
\item Montrer que $m$ est nul si et seulement si l'image de $m$ est nulle
dans $M_\mathfrak m$ pour tout idéal maximal $\mathfrak m$ de $A$.
\item Montrer que $M = 0$ si et seulement si $M_{\mathfrak m}$ est nul
pour tout idéal maximal $\mathfrak m$.
\end{enumerate}
\end{exercise}

\begin{exercise}
\label{exercise-localization-is-field}
Trouver un couple $(A, f)$ où $A$ est un anneau intègre possédant au moins trois idéaux
premiers deux à deux distincts et où $f \in A$ est tel que le localisé
principal $A_f = \{1, f, f^2, \ldots \}^{-1}A$ soit un corps.
\end{exercise}

\begin{exercise}
\label{exercise-localize-finite-module-zero}
Soit $A$ un anneau. Soit $M$ un $A$-module de type fini. Soit $S \subset A$
une partie multiplicative. Supposons que $S^{-1}M = 0$. Montrer qu'il
existe $f \in S$ tel que le localisé principal
$M_f = \{1, f, f^2, \ldots \}^{-1}M$ soit nul.
\end{exercise}
% Borne de la tranche : la ligne source 750 est vide.
\begin{exercise}
\label{exercise-localization-is-quotient}
Donner un exemple de triplet $(A, I, S)$ où $A$ est un anneau,
$0 \not = I \not = A$ est un idéal propre non nul et $S \subset A$
est une partie multiplicative
telle que $A/I \cong S^{-1}A$ comme $A$-algèbres.
\end{exercise}




\section{Lemme de Nakayama}
\label{section-nakayama}

\begin{exercise}
\label{exercise-nakayama}
Soit $A$ un anneau.
Soit $I$ un idéal de $A$.
Soit $M$ un $A$-module.
Soient $x_1, \ldots, x_n \in M$.
Supposons satisfaites les conditions suivantes :
\begin{enumerate}
\item $M/IM$ est engendré par $x_1, \ldots, x_n$,
\item $M$ est un $A$-module de type fini,
\item $I$ est contenu dans tout idéal maximal de $A$.
\end{enumerate}
Montrer que $x_1, \ldots, x_n$ engendrent $M$. (Solution suggérée :
se ramener, par localisation, au cas d'un idéal maximal de $A$ à l'aide de
l'exercice \ref{exercise-localize-zero} et de l'exactitude de la localisation.
Se ramener ensuite à l'énoncé du lemme de Nakayama donné dans le cours
en considérant le quotient de $M$ par le sous-module engendré par
$x_1, \ldots, x_n$.)
\end{exercise}






\section{Longueur}
\label{section-length}

\begin{definition}
\label{definition-length}
Soient $A$ un anneau et $M$ un $A$-module. La
{\it longueur} de $M$ comme $A$-module est
$$
\text{longueur}_A(M)
=
\sup
\{
n
\mid
\exists\ 0 = M_0 \subset M_1 \subset \ldots \subset M_n = M,
\text{ }M_i \not = M_{i + 1}
\}.
$$
Autrement dit, c'est la borne supérieure des longueurs des chaînes de sous-modules.
\end{definition}

\begin{exercise}
\label{exercise-length-is-one}
Montrer qu'un module $M$ sur un anneau $A$ est de longueur $1$ si et seulement s'il
est isomorphe à $A/\mathfrak m$ pour un idéal maximal $\mathfrak m$
de $A$.
\end{exercise}

\begin{exercise}
\label{exercise-length-easy}
Calculer la longueur des modules suivants sur les anneaux indiqués.
Expliquer brièvement (!) chaque réponse. (On pourra utiliser librement l'additivité de
la fonction longueur dans les suites exactes courtes ; voir
Algèbre, lemme \ref{algebra-lemma-length-additive}.)
\begin{enumerate}
\item La longueur de $\mathbf{Z}/120\mathbf{Z}$ sur $\mathbf{Z}$.
\item La longueur de $\mathbf{C}[x]/(x^{100} + x + 1)$ sur $\mathbf{C}[x]$.
\item La longueur de $\mathbf{R}[x]/(x^4 + 2x^2 + 1)$ sur $\mathbf{R}[x]$.
\end{enumerate}
\end{exercise}

\begin{exercise}
\label{exercise-compute-length}
Soit $A = k[x, y]_{(x, y)}$ l'anneau local du plan affine en
l'origine. On pourra imposer au corps $k$ toute hypothèse que l'on souhaite. Supposons
que $f = x^3 + x^2y^2 + y^{100}$ et $g = y^3 - x^{999}$. Quelle est la longueur
de $A/(f, g)$ comme $A$-module ? (Une démarche possible consiste à considérer
l'idéal engendré par $f$ et $g$ dans les quotients de la forme $A/{\mathfrak m}_A^n=
k[x, y]/(x, y)^n$, où $n$ varie. Essayer de trouver $n$ tel que
$A/(f, g)+{\mathfrak m}_A^n \cong A/(f, g)+{\mathfrak m}_A^{n + 1}$
et utiliser NAK.)
\end{exercise}






\section{Idéaux premiers associés}
\label{section-ass}

\noindent
Les idéaux premiers associés sont étudiés dans
Algèbre, section \ref{algebra-section-ass}.

\begin{exercise}
\label{exercise-compute-ass}
Calculer l'ensemble des idéaux premiers associés à chacun des modules suivants.
\begin{enumerate}
\item $R = k[x, y]$ et $M = R/(xy(x + y))$,
\item $R = \mathbf{Z}[x]$ et $M = R/(300x + 75)$, et
\item $R = k[x, y, z]$ et $M = R/(x^3, x^2y, xz)$.
\end{enumerate}
Comme d'habitude, $k$ désigne un corps.
\end{exercise}

\begin{exercise}
\label{exercise-prime-power-not-primary}
Donner un exemple d'anneau noethérien $R$ et d'idéal premier $\mathfrak p$
tel que $\mathfrak p$ ne soit pas le seul idéal premier associé à $R/\mathfrak p^2$.
\end{exercise}

\begin{exercise}
\label{exercise-product-primes-not-only-associated}
Soit $R$ un anneau noethérien possédant deux idéaux premiers incomparables
$\mathfrak p$, $\mathfrak q$, c'est-à-dire $\mathfrak p \not \subset \mathfrak q$
et $\mathfrak q \not \subset \mathfrak p$.
\begin{enumerate}
\item Montrer que, pour $N = R/(\mathfrak p \cap \mathfrak q)$,
on a $\text{Ass}(N) = \{\mathfrak p, \mathfrak q\}$.
\item Montrer par un exemple que le module
$M = R/\mathfrak p \mathfrak q$ peut
avoir un idéal premier associé distinct de $\mathfrak p$ et de $\mathfrak q$.
\end{enumerate}
\end{exercise}





\section{Groupes d'Ext}
\label{section-ext}

\noindent
Les groupes d'Ext sont définis dans Algèbre, section \ref{algebra-section-ext}.

\begin{exercise}
\label{exercise-compute-ext-abelian-groups}
Calculer tous les groupes d'Ext $\Ext^i(M, N)$ des modules donnés
dans la catégorie des $\mathbf{Z}$-modules
(appelée aussi catégorie des groupes abéliens).
\begin{enumerate}
\item $M = \mathbf{Z}$ et $N = \mathbf{Z}$,
\item $M = \mathbf{Z}/4\mathbf{Z}$ et $N = \mathbf{Z}/8\mathbf{Z}$,
\item $M = \mathbf{Q}$ et $N = \mathbf{Z}/2\mathbf{Z}$, et
\item $M = \mathbf{Z}/2\mathbf{Z}$ et $N = \mathbf{Q}/\mathbf{Z}$.
\end{enumerate}
\end{exercise}

\begin{exercise}
\label{exercise-compute-in-regular}
Soit $R = k[x, y]$, où $k$ est un corps.
\begin{enumerate}
\item Montrer directement que le complexe de Koszul
$$
0 \to R
\xrightarrow{
\left(
\begin{matrix}
y \\
-x
\end{matrix}
\right)
}
R^{\oplus 2} \xrightarrow{(x, y)} R \xrightarrow{f \mapsto f(0, 0)} k \to 0
$$
est exact.
\item Calculer $\Ext^i_R(k, k)$, où $k = R/(x, y)$ est considéré comme un $R$-module.
\end{enumerate}
\end{exercise}

\begin{exercise}
\label{exercise-infinitely-many-nonzero-ext}
Donner un exemple d'anneau noethérien $R$ et de modules de type fini $M$, $N$
tels que $\Ext^i_R(M, N)$ soit non nul pour tout $i \geq 0$.
\end{exercise}

\begin{exercise}
\label{exercise-infinite-ext}
Donner un exemple d'anneau $R$ et d'idéal $I$ tels que
$\Ext^1_R(R/I, R/I)$ ne soit pas un $R$-module de type fini.
(On sait que cela ne peut pas se produire si $R$ est noethérien, d'après
Algèbre, lemme \ref{algebra-lemma-ext-noetherian}.)
\end{exercise}







\section{Profondeur}
\label{section-depth}

\noindent
La profondeur est définie dans Algèbre, section \ref{algebra-section-depth},
et étudiée plus avant dans
Complexes dualisants, section \ref{dualizing-section-depth}.

\begin{exercise}
\label{exercise-compute-depth}
Soient $R$ un anneau, $I \subset R$ un idéal et $M$ un $R$-module.
Calculer $\text{prof}_I(M)$ dans les cas suivants.
\begin{enumerate}
\item $R = \mathbf{Z}$, $I = (30)$, $M = \mathbf{Z}$,
\item $R = \mathbf{Z}$, $I = (30)$, $M = \mathbf{Z}/(300)$,
\item $R = \mathbf{Z}$, $I = (30)$, $M = \mathbf{Z}/(7)$,
\item $R = k[x, y, z]/(x^2 + y^2 + z^2)$, $I = (x, y, z)$, $M = R$,
\item $R = k[x, y, z, w]/(xz, xw, yz, yw)$, $I = (x, y, z, w)$,
$M = R$.
\end{enumerate}
Ici, $k$ est un corps. Dans les deux derniers cas, on pourra localiser
en l'idéal maximal $I$.
\end{exercise}

\begin{exercise}
\label{exercise-depth-not-inherited-localization}
Donner un exemple d'anneau local noethérien $(R, \mathfrak m, \kappa)$
de profondeur $\geq 1$ et d'idéal premier $\mathfrak p$ tels que
\begin{enumerate}
\item $\text{prof}_\mathfrak m(R) \geq 1$,
\item $\text{prof}_\mathfrak p(R_\mathfrak p) = 0$, et
\item $\dim(R_\mathfrak p) \geq 1$.
\end{enumerate}
Si l'on ne demande pas (3), l'exercice est trop facile. Pourquoi ?
\end{exercise}

\begin{exercise}
\label{exercise-depth-torsion-free}
Soit $(R, \mathfrak m)$ un anneau local noethérien intègre.
Soit $M$ un $R$-module de type fini.
\begin{enumerate}
\item Si $M$ est sans torsion, montrer que $M$ est de profondeur au moins $1$ sur $R$.
\item Donner un exemple où la profondeur vaut $1$.
\end{enumerate}
\end{exercise}

\begin{exercise}
\label{exercise-depth-examples}
Pour tous $m \geq n \geq 0$, donner un exemple d'anneau local
noethérien $R$ tel que $\dim(R) = m$ et $\text{prof}(R) = n$.
\end{exercise}

\begin{exercise}
\label{exercise-make-depth-1}
Soit $(R, \mathfrak m)$ un anneau local noethérien.
Soit $M$ un $R$-module de type fini.
Montrer qu'il existe une suite exacte courte canonique
$$
0 \to K \to M \to Q \to 0
$$
ayant les propriétés suivantes :
\begin{enumerate}
\item $\text{prof}(Q) \geq 1$,
\item $K$ est nul ou $\text{Supp}(K) = \{\mathfrak m\}$, et
\item $\text{longueur}_R(K) < \infty$.
\end{enumerate}
Indication : en utilisant la propriété noethérienne, montrer qu'il
existe un sous-module $K$ maximal parmi ceux qui vérifient (2),
puis que $Q = M/K$ vérifie (1) et que $K$ vérifie (3).
\end{exercise}

\begin{exercise}
\label{exercise-make-depth-2}
Soit $(R, \mathfrak m)$ un anneau local noethérien.
Soit $M$ un $R$-module de type fini et de profondeur $\geq 2$.
Soit $N \subset M$ un sous-module non nul.
\begin{enumerate}
\item Montrer que $\text{prof}(N) \geq 1$.
\item Montrer que $\text{prof}(N) = 1$ si et seulement si le
module quotient $M/N$ vérifie $\text{prof}(M/N) = 0$.
\item Montrer qu'il existe un sous-module $N' \subset M$ tel que
$N \subset N'$ et que le quotient soit de longueur finie, c'est-à-dire que
$\text{longueur}_R(N'/N) < \infty$, et tel que $N'$ soit
de profondeur $\geq 2$. Indication : appliquer l'exercice \ref{exercise-make-depth-1} à $M/N$
et prendre pour $N'$ l'image inverse de $K$.
\end{enumerate}
\end{exercise}

\begin{exercise}
\label{exercise-Hartshorne-reduced}
Soit $(R, \mathfrak m)$ un anneau local noethérien.
Supposons $R$ réduit, c'est-à-dire que $R$ ne possède aucun élément
nilpotent non nul. Supposons en outre que $R$ possède deux
idéaux premiers minimaux distincts $\mathfrak p$ et $\mathfrak q$.
\begin{enumerate}
\item Montrer que la suite de $R$-modules
$$
0 \to R \to R/\mathfrak p \oplus R/\mathfrak q \to
R/\mathfrak p + \mathfrak q \to 0
$$
est exacte (le vérifier en chaque terme). Les applications sont
$x \mapsto (x \bmod \mathfrak p, x \bmod \mathfrak q)$
et $(y \bmod \mathfrak p, z \bmod \mathfrak q) \mapsto
(y - z \bmod \mathfrak p + \mathfrak q)$.
\item Montrer que si $\text{prof}(R) \geq 2$, alors
$\dim(R/\mathfrak p + \mathfrak q) \geq 1$.
\item Montrer que si $\text{prof}(R) \geq 2$, alors
$U = \Spec(R) \setminus \{\mathfrak m\}$ est un espace topologique
connexe.
\end{enumerate}
Cela démontre un cas très particulier du théorème de connexité de Hartshorne,
selon lequel le spectre épointé $U$ d'un anneau local
noethérien de $\text{prof} \geq 2$ est connexe.
\end{exercise}

\begin{exercise}
\label{exercise-depth-2}
Soit $(R, \mathfrak m)$ un anneau local noethérien.
Supposons que $x, y \in \mathfrak m$ forment une suite régulière de longueur $2$.
Pour tout $n \geq 2$, montrer qu'il n'existe pas de $a, b \in R$ tels que
$$
x^{n - 1}y^{n - 1} = a x^n + b y^n
$$
Indication : commencer par traiter le cas $n = 2$ afin de voir comment procéder.
Remarque : ce résultat admet une vaste généralisation
appelée la conjecture des monômes.
\end{exercise}







\section{Modules et anneaux de Cohen-Macaulay}
\label{section-CM}

\noindent
Les modules de Cohen-Macaulay sont étudiés dans
Algèbre, section \ref{algebra-section-CM},
et les anneaux de Cohen-Macaulay sont étudiés dans
Algèbre, section \ref{algebra-section-CM-ring}.

\begin{exercise}
\label{exercise-examples-CM}
Dans chacun des cas suivants, répondre par oui ou par non.
Aucune explication ni démonstration n'est nécessaire.
\begin{enumerate}
\item Soit $p$ un nombre premier.
L'anneau local $\mathbf{Z}_{(p)}$ est-il un anneau local de Cohen-Macaulay ?
\item Soit $p$ un nombre premier.
L'anneau local $\mathbf{Z}_{(p)}$ est-il un anneau local régulier ?
\item Soit $k$ un corps.
L'anneau local $k[x]_{(x)}$ est-il un anneau local de Cohen-Macaulay ?
\item Soit $k$ un corps.
L'anneau local $k[x]_{(x)}$ est-il un anneau local régulier ?
\item Soit $k$ un corps.
L'anneau local $(k[x, y]/(y^2 - x^3))_{(x, y)} =
k[x, y]_{(x, y)}/(y^2 - x^3)$ est-il un anneau local de Cohen-Macaulay ?
\item Soit $k$ un corps.
L'anneau local $(k[x, y]/(y^2, xy))_{(x, y)} =
k[x, y]_{(x, y)}/(y^2, xy)$ est-il un anneau local de Cohen-Macaulay ?
\end{enumerate}
\end{exercise}













\section{Singularités}
\label{section-singularities}

\begin{exercise}
\label{exercise-singularities}
Soit $k$ un corps quelconque. Supposons que $A = k[[x, y]]/(f)$ et
$B = k[[u, v]]/(g)$, où $f = xy$ et $g = uv + \delta$ avec
$\delta \in (u, v)^3$. Montrer que $A$ et $B$ sont des anneaux isomorphes.
\end{exercise}

\begin{remark}
\label{remark-singularities}
Une singularité d'une courbe sur un corps $k$ est appelée
point double ordinaire si l'anneau local complété de la courbe au
point est de la forme $k'[[x, y]]/(f)$, où (a) $k'$ est une extension finie
séparable de $k$, (b) le terme initial de $f$ est de degré deux, c'est-à-dire qu'il
est de la forme $q = ax^2 + bxy + cy^2$ pour des $a, b, c\in k'$ non tous nuls, et
(c) $q$ est une forme quadratique non dégénérée sur $k'$ (en caractéristique 2,
cela signifie que $b$ est non nul). En général, il y a une classe d'isomorphisme
de tels anneaux pour chaque classe d'isomorphisme de couples $(k', q)$.
\end{remark}

\begin{exercise}
\label{exercise-periodic-resolution}
Soit $R$ un anneau. Soit $n \geq 1$. Soient $A$, $B$ des matrices $n \times n$ à
coefficients dans $R$ telles que $AB = f 1_{n \times n}$
pour un non-diviseur de zéro $f$ de $R$. Posons $S = R/(f)$. Montrer que
$$
\ldots \to
S^{\oplus n} \xrightarrow{B}
S^{\oplus n} \xrightarrow{A}
S^{\oplus n} \xrightarrow{B}
S^{\oplus n} \to \ldots
$$
est exacte.
\end{exercise}







\section{Ensembles constructibles}
\label{section-constructible}

\noindent
Soit $k$ un corps algébriquement clos, par exemple le corps $\mathbf{C}$
des nombres complexes. Soit $n \geq 0$. Un polynôme $f \in k[x_1, \ldots, x_n]$
définit par évaluation une fonction $f : k^n \to k$. Un sous-ensemble $Z \subset k^n$
est appelé un {\it ensemble algébrique} s'il est l'ensemble des zéros communs d'une
famille de polynômes.

\begin{exercise}
\label{exercise-finite-nr-equations}
Démontrer qu'un ensemble algébrique peut toujours s'écrire comme l'ensemble des zéros
d'un nombre fini de polynômes.
\end{exercise}

\noindent
Avec les notations ci-dessus, un sous-ensemble $E \subset k^n$ est dit
{\it constructible} s'il est une réunion finie d'ensembles de la forme
$Z \cap \{f \not = 0\}$, où $f$ est un polynôme.

\begin{exercise}
\label{exercise-constructible-classical}
Montrer les assertions suivantes :
\begin{enumerate}
\item le complémentaire d'un ensemble constructible est constructible,
\item une réunion finie d'ensembles constructibles est constructible,
\item une intersection finie d'ensembles constructibles est constructible, et
\item tout ensemble constructible $E$ peut s'écrire comme une réunion disjointe finie
$E = \coprod E_i$, chaque $E_i$ étant de la forme $Z \cap \{f \not = 0\}$,
où $Z$ est un ensemble algébrique et $f$ un polynôme.
\end{enumerate}
\end{exercise}

\begin{exercise}
\label{exercise-division-with-remainder}
Soit $R$ un anneau. Soit
$f = a_d x^d + a_{d - 1} x^{d - 1} + \ldots + a_0 \in R[x]$.
(Comme d'habitude, cette notation signifie que $a_0, \ldots, a_d \in R$.) Soit $g \in R[x]$.
Démontrer qu'il existe $N \geq 0$ et $r, q \in R[x]$ tels que
$$
a_d^N g = q f + r
$$
avec $\deg(r) < d$, c'est-à-dire que, pour certains $c_i \in R$, on a
$r = c_0 + c_1 x + \ldots + c_{d - 1}x^{d - 1}$.
\end{exercise}







\section{Nullstellensatz de Hilbert}
\label{section-Hilbert-Nullstellensatz}


\begin{exercise}
\label{exercise-uncountable}
{\it Un argument saugrenu utilisant les nombres complexes !}
Soit ${\mathbf C}$ le corps des nombres complexes. Soit $V$ un espace vectoriel
sur ${\mathbf C}$. Le spectre d'un opérateur linéaire
$T : V \to V$ est l'ensemble des nombres complexes $\lambda \in {\mathbf C}$
tels que l'opérateur $T - \lambda \text{id}_V$ ne soit pas inversible.
\begin{enumerate}
\item Montrer que $\mathbf{C}(X)$
est de dimension non dénombrable sur ${\mathbf C}$.
\item Montrer que tout opérateur linéaire sur $V$ possède un
spectre non vide si la dimension de $V$ est finie ou
dénombrable.
\item Montrer que si une ${\mathbf C}$-algèbre de type fini
$R$ est un corps, alors l'application ${\mathbf C}\to R$ est un isomorphisme.
\item Montrer que tout idéal maximal ${\mathfrak m}$ de
${\mathbf C}[x_1, \ldots, x_n]$ est de la forme
$(x_1-\alpha_1, \ldots, x_n-\alpha_n)$ pour certains $\alpha_i \in {\mathbf C}$.
\end{enumerate}
\end{exercise}

\begin{remark}
\label{remark-HNSS}
Soit $k$ un corps. Alors, pour tout entier $n\in {\mathbf N}$ et
tout idéal maximal ${\mathfrak m} \subset k[x_1, \ldots, x_n]$,
le quotient $k[x_1, \ldots, x_n]/{\mathfrak m}$ est une extension finie
de $k$. Ce résultat sera démontré plus loin dans le cours. Bien entendu
(le vérifier), il entraîne un énoncé analogue pour les idéaux maximaux
des $k$-algèbres de type fini. L'exercice précédent le démontre
dans le cas $k = {\mathbf C}$.
\end{remark}

\begin{exercise}
\label{exercise-Hilbert-Nullstellensatz}
Soit $k$ un corps. Utiliser la remarque \ref{remark-HNSS}.
\begin{enumerate}
\item Soit $R$ une $k$-algèbre. Supposons que $\dim_k R < \infty$
et que $R$ soit intègre. Montrer que $R$ est un corps.
\item Supposons que $R$ soit une $k$-algèbre de type fini et que
$f\in R$ ne soit pas nilpotent. Montrer qu'il existe un idéal maximal
${\mathfrak m} \subset R$ tel que $f\not\in {\mathfrak m}$.
\item Montrer par un exemple que cette assertion est fausse lorsque $R$
n'est pas de type fini sur un corps.
\item Montrer que tout idéal radical $I \subset {\mathbf C}[x_1, \ldots, x_n]$
est l'intersection des idéaux maximaux qui le contiennent.
\end{enumerate}
\end{exercise}

\begin{remark}
\label{remark-Hilbert-Nullstellensatz}
C'est le Nullstellensatz de Hilbert. Il affirme en effet
que les parties fermées de $\Spec(k[x_1, \ldots, x_n])$
(qui correspondent aux idéaux radicaux d'après un exercice précédent)
sont déterminées par les points fermés qu'elles contiennent.
\end{remark}

\begin{exercise}
\label{exercise-product-matrices-ring}
Soit $A =
{\mathbf C}[x_{11}, x_{12}, x_{21}, x_{22}, y_{11}, y_{12}, y_{21}, y_{22}]$.
Soit $I$ l'idéal de $A$ engendré par les coefficients de la
matrice $XY$, où
$$
X = \left(
\begin{matrix}
x_{11} & x_{12}\\
x_{21} & x_{22}
\end{matrix}
\right)
\quad\text{et}\quad
Y = \left(
\begin{matrix}
y_{11} & y_{12}\\
y_{21} & y_{22}
\end{matrix}
\right).
$$
Déterminer les composantes irréductibles du sous-ensemble fermé $V(I)$ de
$\Spec(A)$.
(Il s'agit de les décrire et de donner des équations pour chacune. Il n'est pas
nécessaire de démontrer que les équations écrites définissent des idéaux premiers.) Indications :
\begin{enumerate}
\item On peut utiliser le Nullstellensatz de Hilbert, et il suffit de déterminer des
sous-ensembles localement fermés irréductibles qui recouvrent l'ensemble des points fermés de
$V(I)$.
\item Il y a deux composantes faciles à trouver.
\item L'image d'un ensemble irréductible par une application continue est
irréductible.
\end{enumerate}
\end{exercise}




\section{Dimension}
\label{section-dimension}

\begin{exercise}
\label{exercise-dimension-bigger-one-finite-nr-primes}
Construire un anneau $A$ de dimension $> 1$ qui n'a qu'un nombre fini d'idéaux premiers.
\end{exercise}

\begin{exercise}
\label{exercise-hypersurface-in-A2-dimension-one}
Soit $f \in \mathbf{C}[x, y]$ un polynôme non constant.
Montrer que $\mathbf{C}[x, y]/(f)$ est de dimension $1$.
\end{exercise}

\begin{exercise}
\label{exercise-dimension-polynomial-ring}
Soit $(R, \mathfrak m)$ un anneau local noethérien.
Soit $n \geq 1$. Soit $\mathfrak m' = (\mathfrak m, x_1, \ldots, x_n)$
dans l'anneau de polynômes $R[x_1, \ldots, x_n]$.
Montrer que
$$
\dim(R[x_1, \ldots, x_n]_{\mathfrak m'}) = \dim(R) + n.
$$
\end{exercise}





\section{Anneaux caténaires}
\label{section-catenary}

\begin{definition}
\label{definition-catenary}
On dit qu'un anneau noethérien $A$ est {\it caténaire}
si, pour tout triplet d'idéaux premiers
${\mathfrak p}_1 \subset {\mathfrak p}_2 \subset {\mathfrak p}_3$,
on a
$$
ht({\mathfrak p}_3 / {\mathfrak p}_1) = ht({\mathfrak p}_3/{\mathfrak p}_2) +
ht({\mathfrak p}_2/{\mathfrak p}_1).
$$
Ici, $ht(\mathfrak p/\mathfrak q)$ désigne la hauteur de
$\mathfrak p/\mathfrak q$ dans l'anneau $A/\mathfrak q$.
Autrement dit,
$$
ht(\mathfrak p/\mathfrak q) =
\dim(A_\mathfrak p/\mathfrak qA_\mathfrak p) =
\dim((A/\mathfrak q)_\mathfrak p) =
\dim((A/\mathfrak q)_{\mathfrak p/\mathfrak q})
$$
Un espace topologique $X$ est {\it caténaire} si, étant donnés $T \subset T' \subset X$
avec $T$ et $T'$ fermés et irréductibles, il existe une chaîne maximale
de parties fermées irréductibles
$$
T = T_0 \subset T_1 \subset \ldots \subset T_n = T'
$$
et toutes ces chaînes ont la même longueur (finie).
\end{definition}

\begin{exercise}
\label{exercise-catenary-the-same}
Montrer que la notion d'anneau caténaire définie dans
Algèbre, définition \ref{algebra-definition-catenary}
coïncide avec celle de la définition \ref{definition-catenary}
pour les anneaux noethériens.
\end{exercise}

\begin{exercise}
\label{exercise-Noetherian-local-domain-dim-2-catenary}
Montrer qu'un anneau local noethérien intègre de dimension $2$ est caténaire.
\end{exercise}

\begin{exercise}
\label{exercise-finite-type-over-field-catenary}
Soit $k$ un corps.
Montrer qu'une $k$-algèbre de type fini est caténaire.
\end{exercise}

\begin{exercise}
\label{exercise-example-no-dim-function}
Donner un exemple d'espace topologique fini, sobre et caténaire
$X$ qui n'admette aucune fonction de dimension $\delta : X \to \mathbf{Z}$.
Ici, $\delta : X \to \mathbf{Z}$ est une fonction de dimension si,
pour $x, y \in X$, on a
\begin{enumerate}
\item $x \leadsto y$ et $x \not = y$ implique $\delta(x) > \delta(y)$,
\item $x \leadsto y$ et $\delta(x) \geq \delta(y) + 2$ implique
qu'il existe $z \in X$ tel que $x \leadsto z \leadsto y$ et
$\delta(x) > \delta(z) > \delta(y)$.
\end{enumerate}
Décrire clairement cet espace et expliquer brièvement pourquoi il
ne peut admettre de fonction de dimension.
\end{exercise}




\section{Corps des fractions}
\label{section-fraction-fields}

\begin{exercise}
\label{exercise-find-fraction-field}
Considérer l'anneau intègre
$$
{\mathbf Q}[r, s, t]/(s^2-(r-1)(r-2)(r-3), t^2-(r + 1)(r + 2)(r + 3)).
$$
Trouver un anneau intègre de la forme ${\mathbf Q}[x, y]/(f)$ dont le
corps des fractions soit isomorphe à celui de l'anneau ci-dessus.
\end{exercise}



\section{Degré de transcendance}
\label{section-transcendence}

\begin{exercise}
\label{exercise-algebraic-extension}
Soient $K'/K/k$ des extensions de corps telles que $K'$ soit
algébrique sur $K$. Montrer que $\text{deg.tr.}_k(K) = \text{deg.tr.}_k(K')$.
(Indication : montrer que si $x_1, \ldots, x_d \in K$ sont algébriquement indépendants
sur $k$ et si $d < \text{deg.tr.}_k(K')$, alors l'extension $k(x_1, \ldots, x_d) \subset K$
ne peut être algébrique.)
\end{exercise}

\begin{exercise}
\label{exercise-growth-powers-subvector-space}
Soient $k$ un corps et $K/k$ une extension de type fini
de degré de transcendance $d$.
Si $V, W \subset K$ sont des sous-espaces vectoriels sur $k$ de dimension finie,
posons
$$
VW = \{f \in K \mid f = \sum\nolimits_{i = 1, \ldots, n} v_i w_i
\text{ pour un certain }n\text{ et }v_i \in V, w_i \in W\}
$$
C'est un sous-espace vectoriel sur $k$ de dimension finie.
Posons $V^2 = VV$, $V^3 = V V^2$, etc.
\begin{enumerate}
\item Montrer qu'on peut trouver $V \subset K$ et $\epsilon > 0$ tels que
$\dim V^n \geq \epsilon n^d$ pour tout $n \geq 1$.
\item Réciproquement, montrer que, pour tout sous-espace $V \subset K$ de dimension finie,
il existe $C > 0$ tel que $\dim V^n \leq C n^d$ pour tout $n \geq 1$.
(Une démarche possible consiste à procéder comme suit :
établir d'abord le résultat pour les sous-espaces vectoriels de $k[x_1, \ldots, x_d]$,
puis pour ceux de $k(x_1, \ldots, x_d)$.
Enfin, si $K/k(x_1, \ldots, x_d)$ est une extension finie,
choisir une base de $K$ sur $k(x_1, \ldots, x_d)$ et
raisonner en développant les éléments sur cette base.)
\item En déduire que l'on peut redéfinir le degré de transcendance en fonction
de la croissance des puissances des sous-espaces vectoriels de dimension finie de $K$.
\end{enumerate}
Cela est lié à la dimension de Gelfand--Kirillov des algèbres
(non commutatives) sur $k$.
\end{exercise}








\section{Dimension des fibres}
\label{section-dimension-fibres}

\noindent
Quelques questions liées à la formule de dimension ; voir
Algèbre, section \ref{algebra-section-dimension-formula}.

\begin{exercise}
\label{exercise-nr-components-fibre}
Soit $k$ un corps algébriquement clos quelconque.
Dans la suite, $k[x]$ et $k[x, y]$ désignent les anneaux de polynômes.
\begin{enumerate}
\item Pour tout entier $n \geq 0$, trouver une extension de type fini
$k[x] \subset A$ d'anneaux intègres telle que le spectre de $A/xA$
ait exactement $n$ composantes irréductibles.
\item Construire un exemple d'extension de type fini $k[x] \subset A$
d'anneaux intègres telle que le spectre de $A/(x - \alpha)A$
soit non vide et réductible pour tout $\alpha \in k$.
\item Construire un exemple d'extension de type fini $k[x, y] \subset A$
d'anneaux intègres telle que le spectre de $A/(x - \alpha, y - \beta)A$
soit irréductible\footnote{Rappelons qu'irréductible implique non vide.}
pour tout $(\alpha, \beta) \in k^2 \setminus \{(0, 0)\}$
et que le spectre de $A/(x, y)A$ soit non vide et réductible.
\end{enumerate}
\end{exercise}

\begin{exercise}
\label{exercise-codim-1}
Soit $k$ un corps algébriquement clos quelconque.
Soit $n \geq 1$.
Soit $k[x_1, \ldots, x_n]$ l'anneau de polynômes.
Posons $\mathfrak m = (x_1, \ldots, x_n)$.
Soit $k[x_1, \ldots, x_n] \subset A$ une extension de type fini
d'anneaux intègres. Posons $d = \dim(A)$.
\begin{enumerate}
\item Montrer que $d - 1 \geq \dim(A/\mathfrak m A) \geq d - n$
si $A/\mathfrak mA \not = 0$.
\item Montrer par des exemples que toutes les valeurs peuvent se produire.
\item Montrer par un exemple que $\Spec(A/\mathfrak m A)$ peut
avoir des composantes irréductibles de dimensions différentes.
\end{enumerate}
\end{exercise}




\section{Modules localement libres de type fini}
\label{section-finite-locally-free}

\begin{definition}
\label{definition-finite-locally-free}
Soit $A$ un anneau. Rappelons qu'un $A$-module $M$ est
{\it localement libre de type fini} si, pour tout ${\mathfrak p} \in \Spec(A)$,
il existe un
$f\in A$, $f \not \in {\mathfrak p}$, tel que $M_f$ soit un $A_f$-module libre
de type fini. On dit que $M$ est un {\it module inversible} si
$M$ est localement libre de type fini et de rang $1$, c'est-à-dire si, pour tout
${\mathfrak p} \in \Spec(A)$, il existe un
$f\in A$, $f \not \in \mathfrak p$, tel que $M_f \cong A_f$
comme $A_f$-module.
\end{definition}

\begin{exercise}
\label{exercise-tensor-finite-locally-free}
Montrer que le produit tensoriel de modules localement libres de type fini
est localement libre de type fini. Montrer que le produit tensoriel de deux
modules inversibles est inversible.
\end{exercise}

\begin{definition}
\label{definition-class-group}
Soit $A$ un anneau. Le {\it groupe des classes de $A$}, parfois appelé
{\it groupe de Picard de $A$}, est l'ensemble $\Pic(A)$
des classes d'isomorphie de $A$-modules inversibles, muni
de la loi de groupe définie par le produit tensoriel (voir
l'exercice \ref{exercise-tensor-finite-locally-free}).
\end{definition}

\noindent
Notons que le groupe des classes de $A$ est trivial si et seulement si tout module
inversible est isomorphe à un module libre de rang 1.

\begin{exercise}
\label{exercise-class-group-trivial}
Montrer que les groupes des classes des anneaux suivants sont triviaux :
\begin{enumerate}
\item un anneau de polynômes $A = k[x]$, où $k$ est un corps,
\item l'anneau des entiers $A = \mathbf{Z}$,
\item un anneau de polynômes $A = k[x, y]$, où $k$ est un corps, et
\item l'anneau quotient $k[x, y]/(xy)$, où $k$ est un corps.
\end{enumerate}
\end{exercise}

\begin{exercise}
\label{exercise-class-group-not-trivial}
Montrer que le groupe des classes de l'anneau
$A = k[x, y]/(y^2 - f(x))$, où $k$ est un corps de caractéristique différente de $2$
et où $f(x) = (x - t_1) \ldots (x - t_n)$ avec $t_1, \ldots, t_n \in k$
distincts et $n \geq 3$ impair, n'est pas trivial. (Indication : montrer que
l'idéal $(y, x - t_1)$ définit un élément non trivial de $\Pic(A)$.)
\end{exercise}

\begin{exercise}
\label{exercise-trace-det}
Soit $A$ un anneau.
\begin{enumerate}
\item Supposons que $M$ soit un $A$-module localement libre de type fini et
que $\varphi : M \to M$ soit un endomorphisme. Définir/construire
la {\it trace} et le {\it déterminant} de $\varphi$, et montrer que cette
construction est « fonctorielle par rapport au triplet $(A, M, \varphi)$ ».
\item Montrer que, si $M, N$ sont des $A$-modules localement libres de type fini,
et si $\varphi : M \to N$ et $\psi : N \to M$, alors
$\text{Trace}(\varphi \circ \psi) = \text{Trace}(\psi \circ \varphi)$ et
$\det(\varphi \circ \psi) = \det(\psi \circ \varphi)$.
\item Lorsque $M$ est localement libre de type fini, montrer que
$\text{Trace}$ définit une application $A$-linéaire $\text{End}_A(M) \to A$ et que
$\det$ définit une application multiplicative $\text{End}_A(M) \to A$.
\end{enumerate}
\end{exercise}

\begin{exercise}
\label{exercise-trace-det-rings}
Supposons maintenant que $B$ soit une $A$-algèbre localement libre
de type fini comme $A$-module ; autrement dit, $B$ est une $A$-algèbre
localement libre de type fini.
\begin{enumerate}
\item Définir $\text{Trace}_{B/A}$ et $\text{Norme}_{B/A}$ au moyen de
$\text{Trace}$ et de $\det$ dans l'exercice \ref{exercise-trace-det}.
\item Soit $b\in B$ et soit $\pi : \Spec(B) \to \Spec(A)$
le morphisme induit. Montrer que $\pi(V(b)) = V(\text{Norme}_{B/A}(b))$.
(Rappelons que $V(f) = \{ {\mathfrak p} \mid f \in {\mathfrak p}\}$.)
\item (Changement de base.) Supposons que $i : A \to A'$ soit un homomorphisme d'anneaux. Posons
$B' = B \otimes_A A'$. Indiquer pourquoi $i(\text{Norme}_{B/A}(b))$ est égal à
$\text{Norme}_{B'/A'}(b \otimes 1)$.
\item Calculer $\text{Norme}_{B/A}(b)$ lorsque
$B = A \times A \times A \times \ldots \times A$
et $b = (a_1, \ldots, a_n)$.
\item Calculer la norme de $y-y^3$ pour l'homomorphisme fini et plat
${\mathbf Q}[x] \to {\mathbf Q}[y]$, $x \to y^n$. (Indication : utiliser
le « changement de base »
$A = {\mathbf Q}[x] \subset A' = {\mathbf Q}(\zeta_n)(x^{1/n})$.)
\end{enumerate}
\end{exercise}



\section{Recollement}
\label{section-glueing}

\begin{exercise}
\label{exercise-cover}
Supposons que $A$ soit un anneau et que $M$ soit un $A$-module.
Soit donnée une famille d'éléments $f_i$, $i \in I$, de $A$ telle que
$$
\Spec(A) = \bigcup D(f_i).
$$
\begin{enumerate}
\item Montrer que, si $M_{f_i}$ est un $A_{f_i}$-module de type fini,
alors $M$ est un $A$-module de type fini.
\item Montrer que, si $M_{f_i}$ est un $A_{f_i}$-module plat,
alors $M$ est un $A$-module plat.
(C'est presque tautologique lorsqu'on y réfléchit.)
\end{enumerate}
\end{exercise}

\begin{remark}
\label{remark-cover}
Dans le langage de la géométrie algébrique, cela signifie que la propriété
d'« être de type fini » ou d'« être plat » est locale pour la topologie de Zariski
(en un sens approprié). On peut aussi le montrer pour la propriété
d'« être de présentation finie ».
\end{remark}

\begin{exercise}
\label{exercise-cover-ring-map}
Supposons que $A \to B$ soit un homomorphisme d'anneaux.
Soient $f_i \in A$, $i \in I$, et $g_j \in B$, $j \in J$, deux familles
d'éléments telles que
$$
\Spec(A) = \bigcup D(f_i)
\quad\text{et}\quad
\Spec(B) = \bigcup D(g_j).
$$
Montrer que, si $A_{f_i} \to B_{f_ig_j}$ est de type fini pour tous $i, j$,
alors $A \to B$ est de type fini.
\end{exercise}




\section{Montée et descente}
\label{section-going-up}


\begin{definition}
\label{definition-GU-GD}
Soit $\phi : A \to B$ un homomorphisme d'anneaux. On dit
que le {\it théorème de montée} est valable pour $\phi$ si la
condition suivante est satisfaite :
\begin{itemize}
\item[(GU)] pour tous ${\mathfrak p}, {\mathfrak p}' \in \Spec(A)$ tels que
${\mathfrak p} \subset {\mathfrak p}'$, et pour tout $P \in \Spec(B)$ situé
au-dessus de ${\mathfrak p}$, il existe $P'\in \Spec(B)$ situé
au-dessus de ${\mathfrak p}'$ tel que $P \subset P'$.
\end{itemize}
De même, on dit que le {\it théorème de descente} est valable pour $\phi$
si la condition suivante est satisfaite :
\begin{itemize}
\item[(GD)] pour tous ${\mathfrak p}, {\mathfrak p}' \in \Spec(A)$ tels que
${\mathfrak p} \subset {\mathfrak p}'$, et pour tout
$P' \in \Spec(B)$ situé
au-dessus de ${\mathfrak p}'$, il existe $P\in \Spec(B)$ situé
au-dessus de ${\mathfrak p}$ tel que $P \subset P'$.
\end{itemize}
\end{definition}

\begin{exercise}
\label{exercise-GU-GD}
Dans chacun des cas suivants, déterminer lesquelles des conditions
(GU) et (GD) sont satisfaites, et expliquer pourquoi. (Utiliser toute proposition,
tout théorème ou tout lemme disponible, mais vérifier les hypothèses dans chaque cas.)
\begin{enumerate}
\item $k$ est un corps, $A = k$, $B = k[x]$.
\item $k$ est un corps, $A = k[x]$, $B = k[x, y]$.
\item $A = {\mathbf Z}$, $B = {\mathbf Z}[1/11]$.
\item $k$ est un corps algébriquement clos, $A = k[x, y]$,
$B = k[x, y, z]/(x^2-y, z^2-x)$.
\item $A = {\mathbf Z}$, $B = {\mathbf Z}[i, 1/(2 + i)]$.
\item $A = {\mathbf Z}$, $B = {\mathbf Z}[i, 1/(14 + 7i)]$.
\item $k$ est un corps algébriquement clos, $A = k[x]$,
$B = k[x, y, 1/(xy-1)]/(y^2-y)$.
\end{enumerate}
\end{exercise}

\begin{exercise}
\label{exercise-image}
Soit $A$ un anneau. Soit $B = A[x]$ l'algèbre polynomiale à
une indéterminée sur $A$. Soit $f = a_0 + a_1 x + \ldots + a_r x^r \in B = A[x]$.
Démontrer soigneusement que l'image de $D(f)$ dans $\Spec(A)$ est égale à
$D(a_0) \cup \ldots \cup D(a_r)$.
\end{exercise}

\begin{exercise}
\label{exercise-images}
Soit $k$ un corps algébriquement clos. Calculer l'image dans
$\Spec(k[x, y])$
des morphismes suivants :
\begin{enumerate}
\item $\Spec(k[x, yx^{-1}]) \to \Spec(k[x, y])$, où
$k[x, y] \subset k[x, yx^{-1}] \subset k[x, y, x^{-1}]$.
(Indication : pour éviter toute confusion, donner un autre nom à l'élément $yx^{-1}$.)
\item $\Spec(k[x, y, a, b]/(ax-by-1))\to \Spec(k[x, y])$.
\item $\Spec(k[t, 1/(t-1)]) \to \Spec(k[x, y])$, induit par $x
\mapsto t^2$,
et $y \mapsto t^3$.
\item $k = {\mathbf C}$ (corps des nombres complexes),
$\Spec(k[s, t]/(s^3 + t^3-1)) \to \Spec(k[x, y])$, où
$x\mapsto s^2$, $y \mapsto t^2$.
\end{enumerate}
\end{exercise}

\begin{remark}
\label{remark-elimination-theory}
La détermination de l'image comme ci-dessus se fait habituellement au moyen de la théorie de l'élimination.
\end{remark}




\section{Idéaux de Fitting}
\label{section-fitting-ideals}

\begin{exercise}
\label{exercise-fitting}
Soient $R$ un anneau et $M$ un $R$-module de type fini.
Choisir une présentation
$$
\bigoplus\nolimits_{j \in J} R \longrightarrow R^{\oplus n}
\longrightarrow M \longrightarrow 0.
$$
de $M$. Noter que l'homomorphisme $R^{\oplus n} \to M$ est défini par une famille
d'éléments $x_1, \ldots, x_n$ de $M$. Les éléments $x_i$
sont des {\it générateurs} de $M$. L'homomorphisme $\bigoplus_{j \in J} R \to R^{\oplus n}$
est défini par une matrice $n \times J$ notée $A$, à coefficients dans $R$.
Autrement dit, $A = (a_{ij})_{i = 1, \ldots, n, j \in J}$.
Les colonnes $(a_{1j}, \ldots, a_{nj})$, $j \in J$, de $A$
sont appelées les {\it relations}. Tout vecteur $(r_i) \in R^{\oplus n}$
tel que $\sum r_i x_i = 0$ est une combinaison linéaire des colonnes de $A$.
Bien entendu, tout $R$-module de type fini possède de nombreuses présentations différentes.
\begin{enumerate}
\item Montrer que l'idéal engendré par les mineurs $(n - k) \times (n - k)$ de
$A$ est indépendant du choix de la présentation.
Cet idéal est l'{\it idéal de Fitting d'indice $k$ de $M$}. On le note $Fit_k(M)$.
\item Montrer que
$Fit_0(M) \subset Fit_1(M) \subset Fit_2(M) \subset \ldots$.
(Indication : utiliser le développement d'un déterminant suivant une colonne.)
\item Montrer que les conditions suivantes sont équivalentes :
\begin{enumerate}
\item $Fit_{r - 1}(M) = (0)$ et $Fit_r(M) = R$, et
\item $M$ est localement libre de rang $r$.
\end{enumerate}
\end{enumerate}
\end{exercise}




\section{Fonctions de Hilbert}
\label{section-hilbert}

\begin{definition}
\label{definition-numerical-polynomial}
Un {\it polynôme numérique} est un polynôme $f(x) \in {\mathbf Q}[x]$
tel que $f(n) \in {\mathbf Z}$ pour tout entier $n$.
\end{definition}

\begin{definition}
\label{definition-graded-module}
Un {\it module gradué} $M$ sur un anneau $A$ est un $A$-module $M$
muni d'une décomposition en somme directe de la forme
$
\bigoplus\nolimits_{n \in {\mathbf Z}} M_n
$
en sous-$A$-modules. On dit que $M$ est {\it localement fini} si tous les
$M_n$ sont des $A$-modules de type fini. Supposons que $A$ soit un anneau noethérien et
que $\varphi$ soit une {\it fonction d'Euler-Poincar\'e} sur les $A$-modules de type fini.
Cela signifie que, pour tout $A$-module de type fini $M$, on se donne un
entier $\varphi(M) \in {\mathbf Z}$ et que, pour toute suite exacte courte
$$
0
\longrightarrow
M'
\longrightarrow
M
\longrightarrow
M''
\longrightarrow
0
$$
on a $\varphi(M) = \varphi(M') + \varphi(M'')$. La {\it fonction de Hilbert}
d'un module gradué localement fini $M$ (relativement à $\varphi$) est la
fonction $\chi_\varphi(M, n) = \varphi(M_n)$. On dit que $M$ admet un
{\it polynôme de Hilbert} s'il existe un polynôme numérique
$P_\varphi$ tel que $\chi_\varphi(M, n) = P_\varphi(n)$ pour tout entier $n$
assez grand.
\end{definition}

\begin{definition}
\label{definition-graded-algebra}
Une {\it $A$-algèbre graduée} est un $A$-module gradué
$B = \bigoplus_{n \geq 0} B_n$ muni d'une application $A$-bilinéaire
$$
B \times B \longrightarrow B, \ (b, b') \longmapsto bb'
$$
qui munit $B$ d'une structure de $A$-algèbre telle que $B_n \cdot B_m \subset B_{n + m}$.
Enfin, un {\it module gradué $M$ sur une $A$-algèbre graduée $B$} est la donnée
d'un $A$-module gradué $M$ muni d'une structure (compatible) de $B$-module
telle que $B_n \cdot M_d \subset M_{n + d}$. On peut alors définir les
{\it homomorphismes de modules/anneaux gradués}, les {\it sous-modules gradués},
les {\it idéaux gradués}, les {\it suites exactes de modules gradués}, etc.
\end{definition}

\begin{exercise}
\label{exercise-Euler-Poincare-field}
Soit le corps $A = k$. Quelles sont toutes les fonctions d'Euler-Poincar\'e
possibles sur les $A$-modules de type fini dans ce cas ?
\end{exercise}

\begin{exercise}
\label{exercise-Euler-Poincare-Z}
Soit $A ={\mathbf Z}$. Quelles sont toutes les fonctions d'Euler-Poincar\'e
possibles sur les $A$-modules de type fini dans ce cas ?
\end{exercise}

\begin{exercise}
\label{exercise-Euler-Poincare-node}
Soit $A = k[x, y]/(xy)$, où $k$ est un corps algébriquement clos. Quelles sont, dans ce cas, toutes les
fonctions d'Euler-Poincar\'e possibles sur les $A$-modules de type fini ?
\end{exercise}

\begin{exercise}
\label{exercise-kernel-locally-finite}
Supposons que $A$ soit noethérien. Montrer que le noyau d'un homomorphisme
de $A$-modules gradués localement finis est localement fini.
\end{exercise}

\begin{exercise}
\label{exercise-no-hilbert}
Soit $k$ un corps, posons $A = k$ et munissons $B = k[x, y]$ de la graduation
déterminée par $\deg(x) = 2$ et $\deg(y) = 3$. Soit $\varphi(M) = \dim_k(M)$.
Calculer la fonction de Hilbert de $B$ comme $k$-module gradué. Existe-t-il
un polynôme de Hilbert dans ce cas ?
\end{exercise}

\begin{exercise}
\label{exercise-no-hilbert-or-is-there}
Soit $k$ un corps, posons $A = k$ et munissons $B = k[x, y]/(x^2, xy)$ de la graduation
déterminée par $\deg(x) = 2$ et $\deg(y) = 3$. Soit $\varphi(M) = \dim_k(M)$.
Calculer la fonction de Hilbert de $B$ comme $k$-module gradué. Existe-t-il
un polynôme de Hilbert dans ce cas ?
\end{exercise}

\begin{exercise}
\label{exercise-hilbert-to-compute}
Soit $k$ un corps et posons $A = k$. Soit $\varphi(M) = \dim_k(M)$.
Fixons $d\in {\mathbf N}$. Considérons l'algèbre graduée sur $A$
$B = k[x, y, z]/(x^d + y^d + z^d)$, où $x, y, z$ sont tous de degré $1$.
Calculer la fonction de Hilbert de $B$. Existe-t-il un polynôme de Hilbert
dans ce cas ?
\end{exercise}



\section{Proj d'un anneau}
\label{section-proj-ring}

\begin{definition}
\label{definition-homogeneous-ideal}
Soit $R$ un anneau gradué. Un idéal {\it homogène} est simplement un idéal
$I \subset R$ qui est en outre un sous-module gradué de $R$. De manière équivalente,
c'est un idéal engendré par des éléments homogènes. De manière équivalente, si
$f \in I$ et si
$$
f = f_0 + f_1 + \ldots + f_n
$$
est la décomposition de $f$ en composantes homogènes dans $R$, alors $f_i \in I$
pour tout $i$.
\end{definition}

\begin{definition}
\label{definition-Proj-R}
On définit le {\it spectre homogène $\text{Proj}(R)$}
de l'anneau gradué $R$ comme l'ensemble des idéaux premiers homogènes
${\mathfrak p}$ de $R$ tels que $R_{+} \not \subset {\mathfrak p}$.
Noter que $\text{Proj}(R)$ est une partie de $\Spec(R)$ et porte donc naturellement
la topologie induite.
\end{definition}

\begin{definition}
\label{definition-Dplus-Vplus}
Soit $R = \oplus_{d \geq 0} R_d$ un anneau gradué, soit $f\in R_d$ et
supposons que $d \geq 1$. On définit {\it $R_{(f)}$} comme le sous-anneau de
$R_f$ formé des éléments de la forme $r/f^n$, où $r$ est homogène et
$\deg(r) = nd$. On définit en outre
$$
D_{+}(f) = \{ {\mathfrak p} \in \text{Proj}(R) | f \not\in {\mathfrak p} \}.
$$
Enfin, pour un idéal homogène $I \subset R$, on définit
$V_{+}(I) = V(I) \cap \text{Proj}(R)$.
\end{definition}

\begin{exercise}
\label{exercise-topology-proj}
Sur la topologie de $\text{Proj}(R)$. Avec les définitions et les notations
ci-dessus, démontrer les assertions suivantes.
\begin{enumerate}
\item Montrer que $D_{+}(f)$ est ouvert dans $\text{Proj}(R)$.
\item Montrer que $D_{+}(ff') = D_{+}(f) \cap D_{+}(f')$.
\item Soit $g = g_0 + \ldots + g_m$ un élément
de $R$ avec $g_i \in R_i$. Exprimer $D(g) \cap \text{Proj}(R)$
en fonction des $D_{+}(g_i)$, $i \geq 1$, et de $D(g_0) \cap \text{Proj}(R)$.
Aucune démonstration n'est nécessaire.
\item Soit $g\in R_0$ un élément homogène de degré $0$.
Exprimer $D(g) \cap \text{Proj}(R)$ en fonction des $D_{+}(f_\alpha)$
pour une famille convenable d'éléments homogènes $f_\alpha \in R$
de degré strictement positif.
\item Montrer que la famille $\{D_{+}(f)\}$ d'ouverts forme une
base de la topologie de $\text{Proj}(R)$.
\item
\label{item-bijection}
Montrer qu'il existe une bijection canonique $D_{+}(f) \to \Spec(R_{(f)})$.
(Indication : imiter la démonstration pour $\Spec$, mais, à un certain stade, faire intervenir le
radical d'un idéal.)
\item Montrer que l'application définie en (\ref{item-bijection}) est un homéomorphisme.
\item Donner un exemple d'un anneau $R$ tel que $\text{Proj}(R)$ ne soit pas
quasi-compact. Aucune démonstration n'est nécessaire.
\item Montrer que toute partie fermée $T \subset \text{Proj}(R)$ est de
la forme $V_{+}(I)$ pour un idéal homogène $I \subset R$.
\end{enumerate}
\end{exercise}

\begin{remark}
\label{remark-continuous-proj-spec}
Il existe une application continue $ \text{Proj}(R) \longrightarrow \Spec(R_0) $.
\end{remark}

\begin{exercise}
\label{exercise-iso-polynomial-ring-one-variable}
Si $R = A[X]$ avec $\deg(X) = 1$, montrer que l'application naturelle
$\text{Proj}(R) \to \Spec(A)$ est une bijection et même
un homéomorphisme.
\end{exercise}

\begin{exercise}
\label{exercise-blowing-up-I}
Éclatement : première partie.
Dans cet exercice, $R = Bl_I(A) = A \oplus I \oplus I^2 \oplus \ldots$.
Considérer l'application naturelle $b : \text{Proj}(R) \to \Spec(A)$.
Poser $U = \Spec(A) - V(I)$. Montrer que
$$
b : b^{-1}(U) \longrightarrow U
$$
est un homéomorphisme.
On peut donc identifier $U$ à une partie ouverte de $\text{Proj}(R)$.
Soit $Z \subset \Spec(A)$ un sous-schéma fermé irréductible
de point générique $\xi \in Z$. Supposer que $\xi \not\in V(I)$,
autrement dit $Z \not\subset V(I)$, autrement dit
$\xi \in U$, autrement dit $Z\cap U \not = \emptyset$. On définit
le {\it transformé pur} $Z'$ de $Z$ comme l'adhérence de l'unique
point $\xi'$ situé au-dessus de $\xi$. Autrement dit,
$Z'$ est l'adhérence dans $\text{Proj}(R)$ de la partie localement fermée
$Z\cap U \subset U \subset \text{Proj}(R)$.
\end{exercise}

\begin{exercise}
\label{exercise-blowing-up-II}
Éclatement : deuxième partie.
Soit $A = k[x, y]$, où $k$ est un corps, et soit $I = (x, y)$.
Soit $R$ l'algèbre d'éclatement associée à $A$ et $I$.
\begin{enumerate}
\item Montrer que les transformés purs de $Z_1 = V(\{x\})$ et de
$Z_2 = V(\{y\})$ sont disjoints.
\item Montrer que les transformés purs de $Z_1 = V(\{x\})$ et de
$Z_2 = V(\{x-y^2\})$ ne sont pas disjoints.
\item Trouver un idéal $J \subset A$ tel que $V(J) = V(I)$
et tel que les transformés purs de $Z_1 = V(\{x\})$ et de
$Z_2 = V(\{x-y^2\})$ dans l'éclatement le long de $J$ soient disjoints.
\end{enumerate}
\end{exercise}

\begin{exercise}
\label{exercise-proj-when-empty}
Soit $R$ un anneau gradué.
\begin{enumerate}
\item Montrer que $\text{Proj}(R)$ est vide si $R_n = (0)$ pour tout $n >> 0$.
\item Montrer que $\text{Proj}(R)$ est un espace topologique irréductible
si $R$ est intègre et si $R_{+}$ n'est pas nul. (Rappelons que l'espace
topologique vide n'est pas irréductible.)
\end{enumerate}
\end{exercise}

\begin{exercise}
\label{exercise-blowing-up-III}
Éclatement : troisième partie.
On reprend $A$, $I$, $U$ et $Z$ comme dans la définition du transformé pur.
Soit $Z = V({\mathfrak p})$ pour un idéal premier ${\mathfrak p}$. Poser $\bar A =
A/{\mathfrak p}$ et noter
$\bar I$ l'image de $I$ dans $\bar A$.
\begin{enumerate}
\item Montrer qu'il existe un homomorphisme d'anneaux surjectif
$R: = Bl_I(A) \to \bar R: = Bl_{\bar I}(\bar A)$.
\item Montrer que l'homomorphisme d'anneaux ci-dessus induit une application bijective
de $\text{Proj}(\bar R)$ sur le transformé pur $Z'$ de $Z$. (Cette question
n'est pas si facile. Indication : utiliser 5(b) ci-dessus.)
\item En déduire que le transformé pur $Z' = V_{+}(P)$, où
$P \subset R$ est l'idéal homogène défini par
$P_d = I^d \cap {\mathfrak p}$.
\item Supposer que $Z_1 = V({\mathfrak p})$ et
$Z_2 = V({\mathfrak q})$ sont des parties fermées irréductibles
définies par des idéaux premiers, que $Z_1 \not \subset Z_2$
et que $Z_2 \not \subset Z_1$. Montrer que l'éclatement de l'idéal
$I = {\mathfrak p} + {\mathfrak q}$ sépare les
transformés purs de $Z_1$ et de $Z_2$,
c'est-à-dire $Z_1' \cap Z_2' = \emptyset$. (Indication : considérer les idéaux homogènes
$P$ et $Q$ de la partie (c), puis $V(P + Q)$.)
\end{enumerate}
\end{exercise}


\section{Anneaux de Cohen-Macaulay de dimension 1}
\label{section-CM-dim-1}

\begin{definition}
\label{definition-CM}
On dit qu'un anneau local noethérien $A$ est un {\it anneau de Cohen-Macaulay}
de dimension $d$ s'il est de dimension $d$ et s'il existe un système
de paramètres $x_1, \ldots, x_d$ pour $A$ tel que $x_i$ soit un non-diviseur de zéro
dans $A/(x_1, \ldots, x_{i-1})$ pour $i = 1, \ldots, d$.
\end{definition}

\begin{exercise}
\label{exercise-CM-dim-1-I}
Anneaux de Cohen-Macaulay de dimension 1. Première partie : théorie.
\begin{enumerate}
\item Soit $(A, {\mathfrak m})$ un anneau local noethérien tel que $\dim A = 1$.
Montrer que, si $x\in {\mathfrak m}$ n'est pas un diviseur de zéro, alors
\begin{enumerate}
\item $\dim A/xA = 0$, autrement dit $A/xA$ est artinien,
autrement dit $\{x\}$ est un système de paramètres pour $A$.
\item $A$ n'a aucun idéal premier associé immergé.
\end{enumerate}
\item Réciproquement, soit $(A, {\mathfrak m})$ un anneau local noethérien de
dimension $1$. Montrer que, si $A$ n'a aucun idéal premier associé immergé, il existe
un non-diviseur de zéro dans ${\mathfrak m}$.
\end{enumerate}
\end{exercise}

\begin{exercise}
\label{exercise-CM-dim-1-II}
Anneaux de Cohen-Macaulay de dimension 1. Deuxième partie : exemples.
\begin{enumerate}
\item Soit $A$ l'anneau local en $(x, y)$ de $k[x, y]/(x^2, xy)$.
\begin{enumerate}
\item Montrer que $A$ est de dimension 1.
\item Démontrer que tout élément de ${\mathfrak m}\subset A$ est un
diviseur de zéro.
\item Trouver $z\in {\mathfrak m}$ tel que $\dim A/zA = 0$
(aucune démonstration n'est nécessaire).
\end{enumerate}
\item Soit $A$ l'anneau local en $(x, y)$ de $k[x, y]/(x^2)$.
Trouver un non-diviseur de zéro dans ${\mathfrak m}$ (aucune démonstration n'est nécessaire).
\end{enumerate}
\end{exercise}

\begin{exercise}
\label{exercise-embedding-dim-1}
Anneaux locaux de dimension de plongement $1$.
Supposer que $(A, {\mathfrak m}, k)$ soit un anneau local noethérien
de dimension de plongement $1$, c'est-à-dire
$$
\dim_k {\mathfrak m}/{\mathfrak m}^2 = 1.
$$
Montrer que la fonction $f(n) = \dim_k {\mathfrak m}^n/{\mathfrak m}^{n + 1}$
est soit constante de valeur $1$, soit telle que ses valeurs soient
$$
1, 1, \ldots, 1, 0, 0, 0, 0, 0, \ldots
$$
\end{exercise}

\begin{exercise}
\label{exercise-regular-local-dim-1}
Anneaux locaux réguliers de dimension $1$.
Supposer que $(A, {\mathfrak m}, k)$ soit un anneau local noethérien régulier de
dimension $1$. Rappelons que cela signifie que $A$ est de dimension $1$
et de dimension de plongement $1$, c'est-à-dire
$$
\dim_k {\mathfrak m}/{\mathfrak m}^2 = 1.
$$
Soit $x\in{\mathfrak m}$ un élément dont la classe dans ${\mathfrak m}/{\mathfrak
m}^2$ n'est pas nulle.
\begin{enumerate}
\item Montrer que, pour tout élément $y$
de ${\mathfrak m}$, il existe un entier $n$ tel que $y$ puisse s'écrire
$y = ux^n$, où $u\in A^\ast$ est une unité.
\item Montrer que $x$ est un non-diviseur de zéro dans $A$.
\item En déduire que $A$ est intègre.
\end{enumerate}
\end{exercise}

\begin{exercise}
\label{exercise-nonzerodivisor-graded}
Soit $(A, {\mathfrak m}, k)$ un anneau local noethérien, et soit
$Gr_{\mathfrak m}(A)$ son gradué associé.
\begin{enumerate}
\item Supposer que $x\in {\mathfrak m}^d$ a pour image un non-diviseur de zéro
$\bar x \in {\mathfrak m}^d/{\mathfrak m}^{d + 1}$ dans la composante de degré $d$ de
$Gr_{\mathfrak m}(A)$.
Montrer que $x$ est un non-diviseur de zéro.
\item Supposer que la profondeur de $A$ est au moins $1$.
Autrement dit, supposer qu'il existe un non-diviseur de zéro $y \in {\mathfrak m}$.
Dans ce cas, on peut améliorer le résultat : supposer seulement que $x\in {\mathfrak m}^d$ a pour image
l'élément $\bar x \in {\mathfrak m}^d/{\mathfrak m}^{d + 1}$ de degré $d$
de $Gr_{\mathfrak m}(A)$, lequel est un non-diviseur de zéro en degrés assez
élevés : $\exists N$ tel que, pour tout $n \geq N$, l'application
de multiplication par $\bar x$
$$
{\mathfrak m}^n/{\mathfrak m}^{n + 1} \longrightarrow
{\mathfrak m}^{n + d}/{\mathfrak m}^{n + d + 1}
$$
soit injective. Montrer alors que $x$ est un non-diviseur de zéro.
\end{enumerate}
\end{exercise}

\begin{exercise}
\label{exercise-embedding-2-dim-1}
Supposer que $(A, {\mathfrak m}, k)$ soit un anneau local noethérien de
dimension $1$. Supposer également que la dimension de plongement de $A$ soit
$2$, c'est-à-dire supposer que
$$
\dim_k {\mathfrak m}/{\mathfrak m}^2 = 2.
$$
Notation : $f(n) = \dim_k {\mathfrak m}^n/{\mathfrak m}^{n + 1}$.
Choisir des générateurs $x, y \in {\mathfrak m}$
et écrire $Gr_{\mathfrak m}(A) = k[\bar x, \bar y]/I$ pour un certain
idéal homogène $I$.
\begin{enumerate}
\item Montrer qu'il existe un élément homogène
$F\in k[\bar x, \bar y]$ tel que $I \subset (F)$, avec égalité
en tout degré assez élevé.
\item Montrer que $f(n) \leq n + 1$.
\item Montrer que, si $f(n) < n + 1$, alors $n \geq \deg(F)$.
\item Montrer que, si $f(n) < n + 1$, alors $f(n + 1) \leq f(n)$.
\item Montrer que $f(n) = \deg(F)$ pour tout $n >> 0$.
\end{enumerate}
\end{exercise}

\begin{exercise}
\label{exercise-CM-dim-1-embedding-dim-2}
Anneaux de Cohen-Macaulay de dimension 1 et de dimension de plongement 2.
Supposer que $(A, {\mathfrak m}, k)$ soit un anneau local noethérien de
Cohen-Macaulay
et de dimension $1$. Supposer également que la dimension de plongement de $A$ soit
$2$, c'est-à-dire supposer que
$$
\dim_k {\mathfrak m}/{\mathfrak m}^2 = 2.
$$
Notations : $f$, $F$, $x, y\in {\mathfrak m}$ et $I$ comme dans
l'exercice \ref{exercise-embedding-2-dim-1}. On pourra utiliser librement
tous les résultats des problèmes ci-dessus.
\begin{enumerate}
\item Supposer que $z\in {\mathfrak m}$ est un élément dont la classe
dans ${\mathfrak m}/{\mathfrak m}^2$ est la forme linéaire
$\alpha \bar x + \beta \bar y \in k[\bar x, \bar y]$
qui est première à $F$.
\begin{enumerate}
\item Montrer que $z$ est un non-diviseur de zéro dans $A$.
\item Poser $d = \deg(F)$.
Montrer que ${\mathfrak m}^n = z^{n + 1-d}{\mathfrak m}^{d-1}$
pour tout $n$ assez grand. (Indication : montrer d'abord que
$z^{n + 1-d}{\mathfrak m}^{d-1} \to {\mathfrak m}^n/{\mathfrak m}^{n + 1}$
est surjective en utilisant ce que l'on sait de $Gr_{\mathfrak m}(A)$. Puis utiliser le lemme de Nakayama.)
\end{enumerate}
\item Quelle condition sur $k$ garantit l'existence
d'un tel $z$ ? (Aucune démonstration n'est nécessaire; c'est trop facile.)

\noindent
Supposons désormais qu'il existe un $z$ comme ci-dessus. Cette hypothèse
est en fait sans restriction essentielle (au sens où l'on peut se ramener à
la situation où elle est satisfaite afin d'obtenir les résultats des
parties (d) et (e) ci-dessous).
\item Montrer maintenant que ${\mathfrak m}^\ell = z^{\ell - d + 1} {\mathfrak m}^{d-1}$
pour tout $\ell \geq d$.
\item En déduire que $I = (F)$.
\item En déduire que la fonction $f$ prend les valeurs
$$
2, 3, 4, \ldots, d-1, d, d, d, d, d, d, d, \ldots
$$
\end{enumerate}
\end{exercise}

\begin{remark}
\label{remark-CM-dim-1-embedding-dim-2}
Cela donne à penser qu'un anneau local noethérien de Cohen-Macaulay de dimension 1
et de dimension de plongement 2 est de la forme $B/FB$, où $B$ est un anneau local
régulier de dimension 2. C'est à peu près exact (sous des hypothèses
convenables de ``bonne tenue'' de l'anneau).
\end{remark}






\section{Une infinité de nombres premiers}
\label{section-many-primes}

\noindent
section regroupant des questions singulières sur les anneaux dans lesquels
les nombres premiers non inversibles sont en nombre infini.

\begin{exercise}
\label{exercise-not-in-Q}
Donner un exemple d'une ${\mathbf Z}$-algèbre $R$ de type fini
possédant les deux propriétés suivantes :
\begin{enumerate}
\item Il n'existe aucun homomorphisme d'anneaux $R \to {\mathbf Q}$.
\item Pour tout nombre premier $p$, il existe un idéal maximal
${\mathfrak m} \subset R$ tel que $R/{\mathfrak m} \cong {\mathbf F}_p$.
\end{enumerate}
\end{exercise}

\begin{exercise}
\label{exercise-strange-fp-1}
Pour $f \in {\mathbf Z}[x, u]$, on pose $f_p(x)
= f(x, x^p) \bmod p \in {\mathbf F}_p[x]$. Donner un exemple
d'un $f \in {\mathbf Z}[x, u]$ tel que les deux propriétés
suivantes soient satisfaites :
\begin{enumerate}
\item Il existe une infinité de nombres premiers $p$ tels que $f_p$
n'ait pas de racine dans ${\mathbf F}_p$.
\item Pour tout $p >> 0$, le polynôme $f_p$ possède un facteur
linéaire ou quadratique.
\end{enumerate}
\end{exercise}

\begin{exercise}
\label{exercise-strange-fp-2}
Pour $f \in {\mathbf Z}[x, y, u, v]$, on pose $f_p(x, y)
= f(x, y, x^p, y^p) \bmod p \in {\mathbf F}_p[x, y]$. Donner un exemple
« intéressant » d'un $f$ tel que $f_p$ soit réductible pour tout $p >> 0$.
Par exemple, $f = xv-yu$ avec $f_p = xy^p-x^py = xy(x^{p-1}-y^{p-1})$ est
« sans intérêt » ; tout $f$ ne dépendant que de $x, u$ est « sans intérêt »,
etc.
\end{exercise}

\begin{remark}
\label{remark-strange-fp}
Soit $h \in {\mathbf Z}[y]$ un polynôme unitaire de degré $d$.
Alors :
\begin{enumerate}
\item L'homomorphisme $A = {\mathbf Z}[x] \to B ={\mathbf Z}[y]$,
$x \mapsto h$ est fini et localement libre de rang $d$.
\item Pour tout nombre premier $p$, l'homomorphisme
$A_p = {\mathbf F}_p[x]\to B_p = {\mathbf F}_p[y]$,
$y \mapsto h(y) \bmod p$ est fini et localement libre de rang $d$.
\end{enumerate}
\end{remark}

\begin{exercise}
\label{exercise-strange-fp-3}
Soient $h, A, B, A_p, B_p$ comme dans la remarque. Pour $f \in {\mathbf Z}[x, u]$, on
pose $f_p(x) = f(x, x^p) \bmod p \in {\mathbf F}_p[x]$. Pour
$g \in {\mathbf Z}[y, v]$, on pose
$g_p(y) = g(y, y^p) \bmod p \in {\mathbf F}_p[y]$.
\begin{enumerate}
\item Donner un exemple de $h$ et de $g$ pour lesquels
il n'existe aucun $f$ ayant la propriété
$$
f_p  =  Norm_{B_p/A_p}(g_p).
$$
\item Montrer que, quels que soient $h$ et $g$ comme ci-dessus,
il existe un $f$ non nul tel que, pour tout $p$, on ait
$$
Norm_{B_p/A_p}(g_p)\quad\text{divides}\quad f_p .
$$
On peut, si on le souhaite, se restreindre au cas $h = y^n$, même avec $n = 2$,
mais l'énoncé est vrai en général.
\item Discuter le rapport avec les exercices 6 et 7 de la série
précédente.
\end{enumerate}
\end{exercise}

\begin{exercise}
\label{exercise-strange-fp-unsolved}
Problèmes non résolus. Ils sont peut-être très difficiles ou peut-être faciles.
Je l'ignore.
\begin{enumerate}
\item Existe-t-il un $f \in {\mathbf Z}[x, u]$ tel que $f_p$ soit
irréductible pour une infinité de $p$ ? (Indication : oui, c'est le cas pour
$f(x, u) = u - x - 1$ ainsi que pour $f(x, u) = u^2 - x^2 + 1$.)
\item Soit $f \in {\mathbf Z}[x, u]$ non nul, et supposons que
$\deg_x(f_p) = dp + d'$ pour tout $p$ assez grand. (Autrement dit, $\deg_u(f) = d$
et le coefficient $c$ de $u^d$ dans $f$ vérifie $\deg_x(c) = d'$.) Supposons que l'on
puisse écrire $d = d_1 + d_2$ et $d' = d'_1 + d'_2$ avec $d_1, d_2 > 0$
et $d'_1, d'_2 \geq 0$, de telle sorte que, pour tout $p$ assez grand,
il existe une factorisation
$$
f_p = f_{1, p} f_{2, p}
$$
avec $\deg_x(f_{1, p}) = d_1p + d'_1$. Est-il vrai que $f$ provient d'une
construction par la norme comme dans l'exercice 4 ? (Plus précisément, existe-t-il $h$ et
$g$ tels que $Norm_{B_p/A_p}(g_p)$ divise $f_p$ pour tout $p >> 0$ ?)
\item Question analogue à celle de (b), mais avec cette fois
$f \in {\mathbf Z}[x_1, x_2, u_1, u_2]$ irréductible et en supposant seulement que
$f_p(x_1, x_2) = f(x_1, x_2, x_1^p, x_2^p) \bmod p$ se factorise pour tout
$p >> 0$.
\end{enumerate}
\end{exercise}


















\section{Catégorie dérivée filtrée}
\label{section-filtered-derived}

\noindent
Pour traiter les exercices de cette section, lire au préalable les développements
d'Homologie, section \ref{homology-section-filtrations}. Nous dirons que
$A$ est un objet filtré de $\mathcal{A}$ pour signifier que $A$ est muni
d'une filtration $F$, que nous omettons de la notation.

\begin{exercise}
\label{exercise-split-injective}
Soit $\mathcal{A}$ une catégorie abélienne.
Soit $I$ un objet filtré de $\mathcal{A}$.
Supposons que la filtration de $I$ soit finie
et que chaque $\text{gr}^p(I)$ soit un objet injectif de $\mathcal{A}$.
Montrer qu'il existe un isomorphisme
$I \cong \bigoplus \text{gr}^p(I)$ respectant les filtrations, la filtration
$F^p(I)$ correspondant à $\bigoplus_{p' \geq p} \text{gr}^p(I)$.
\end{exercise}

\begin{exercise}
\label{exercise-filtered-injective}
Soit $\mathcal{A}$ une catégorie abélienne.
Soit $I$ un objet filtré de $\mathcal{A}$.
Supposons que la filtration de $I$ soit finie.
Montrer l'équivalence des propriétés suivantes :
\begin{enumerate}
\item Pour tout diagramme dont les flèches pleines sont données
$$
\xymatrix{
A \ar[r]_\alpha \ar[d] & B \ar@{-->}[ld] \\
I &
}
$$
d'objets filtrés, si
(\romannumeral1) les filtrations de $A$ et de $B$ sont finies,
et (\romannumeral2) $\text{gr}(\alpha)$ est injectif, alors la
flèche en pointillé existe et rend le diagramme commutatif.
\item Chaque $\text{gr}^p I$ est injectif.
\end{enumerate}
\end{exercise}

\noindent
Notons que, si $\alpha : A \to B$ est un morphisme d'objets filtrés
à filtrations finies, dire que $\text{gr}(\alpha)$ est injectif
revient à dire que $\alpha$ est un {\it monomorphisme strict}
dans la catégorie $\text{Fil}(\mathcal{A})$. En effet,
être un monomorphisme signifie que $\Ker(\alpha) = 0$, tandis qu'être strict signifie que
cela entraîne aussi $\Ker(\text{gr}(\alpha)) = 0$.
Voir Homologie, lemme \ref{homology-lemma-characterize-strict}.
(Nous n'employons le terme « injectif » pour un morphisme que dans une catégorie abélienne,
bien qu'il ait un sens dans toute catégorie additive admettant des noyaux.)
Les exercices ci-dessus justifient la définition suivante.

\begin{definition}
\label{definition-injective-filtered}
Soit $\mathcal{A}$ une catégorie abélienne.
Soit $I$ un objet filtré de $\mathcal{A}$.
Supposons que la filtration de $I$ soit finie.
Nous dirons que $I$ est {\it injectif filtré} si chaque $\text{gr}^p(I)$ est
un objet injectif de $\mathcal{A}$.
\end{definition}

\noindent
Nous adoptons la définition suivante afin de ne pas avoir à répéter
partout « muni d'une filtration finie ».

\begin{definition}
\label{definition-finite-filtration-category}
Soit $\mathcal{A}$ une catégorie abélienne.
Nous noterons {\it $\text{Fil}^f(\mathcal{A})$} la sous-catégorie pleine
de $\text{Fil}(\mathcal{A})$ formée des objets
$A \in \Ob(\text{Fil}(\mathcal{A}))$
dont la filtration est finie.
\end{definition}

\begin{exercise}
\label{exercise-inject-into-injective}
Soit $\mathcal{A}$ une catégorie abélienne.
Supposons que $\mathcal{A}$ admette suffisamment d'injectifs.
Soit $A$ un objet de $\text{Fil}^f(\mathcal{A})$.
Montrer qu'il existe un monomorphisme strict $\alpha : A \to I$
de $A$ dans un objet injectif filtré $I$ de $\text{Fil}^f(\mathcal{A})$.
\end{exercise}

\begin{definition}
\label{definition-filtered-quasi-isomorphism}
Soit $\mathcal{A}$ une catégorie abélienne.
Soit $\alpha : K^\bullet \to L^\bullet$ un morphisme de
complexes de $\text{Fil}(\mathcal{A})$. Nous dirons que
$\alpha$ est un {\it quasi-isomorphisme filtré} si,
pour tout $p \in \mathbf{Z}$, le morphisme
$\text{gr}^p(K^\bullet) \to \text{gr}^p(L^\bullet)$ est
un quasi-isomorphisme.
\end{definition}

\begin{definition}
\label{definition-filtered-acyclic}
Soit $\mathcal{A}$ une catégorie abélienne.
Soit $K^\bullet$ un complexe de $\text{Fil}^f(\mathcal{A})$.
Nous dirons que $K^\bullet$ est {\it acyclique filtré} si,
pour tout $p \in \mathbf{Z}$, le complexe $\text{gr}^p(K^\bullet)$ est
acyclique.
\end{definition}

\begin{exercise}
\label{exercise-filtered-quasi-isomorphism}
Soit $\mathcal{A}$ une catégorie abélienne.
Soit $\alpha : K^\bullet \to L^\bullet$ un morphisme de complexes bornés inférieurement
de $\text{Fil}^f(\mathcal{A})$. (Noter l'exposant $f$.)
Montrer l'équivalence des propriétés suivantes :
\begin{enumerate}
\item $\alpha$ est un quasi-isomorphisme filtré,
\item pour tout $p \in \mathbf{Z}$, l'application
$\alpha : F^pK^\bullet \to F^pL^\bullet$ est un quasi-isomorphisme,
\item pour tout $p \in \mathbf{Z}$, l'application
$\alpha : K^\bullet/F^pK^\bullet \to L^\bullet/F^pL^\bullet$
est un quasi-isomorphisme, et
\item le cône de $\alpha$ (voir
Catégories dérivées, définition \ref{derived-definition-cone})
est un complexe acyclique filtré.
\end{enumerate}
Montrer de plus que, si $\alpha$ est un quasi-isomorphisme filtré,
alors $\alpha$ est aussi un quasi-isomorphisme au sens usuel.
\end{exercise}

\begin{exercise}
\label{exercise-injective-resolution}
Soit $\mathcal{A}$ une catégorie abélienne.
Supposons que $\mathcal{A}$ admette suffisamment d'injectifs.
Soit $A$ un objet de $\text{Fil}^f(\mathcal{A})$.
Montrer qu'il existe un complexe
$I^\bullet$ de $\text{Fil}^f(\mathcal{A})$,
et un morphisme $A[0] \to I^\bullet$ tels que
\begin{enumerate}
\item chaque $I^p$ est injectif filtré,
\item $I^p = 0$ pour $p < 0$, et
\item $A[0] \to I^\bullet$ est un quasi-isomorphisme filtré.
\end{enumerate}
\end{exercise}

\begin{exercise}
\label{exercise-injective-resolution-complex}
Soit $\mathcal{A}$ une catégorie abélienne.
Supposons que $\mathcal{A}$ admette suffisamment d'injectifs.
Soit $K^\bullet$ un complexe borné inférieurement d'objets de
$\text{Fil}^f(\mathcal{A})$. Montrer qu'il existe un
quasi-isomorphisme filtré $\alpha : K^\bullet \to I^\bullet$
où $I^\bullet$ est un complexe de $\text{Fil}^f(\mathcal{A})$
dont les termes $I^n$ sont injectifs filtrés et qui est borné inférieurement.
En fait, on peut choisir $\alpha$ de sorte que chaque $\alpha^n$ soit
un monomorphisme strict.
\end{exercise}

\begin{exercise}
\label{exercise-morphisms-lift}
Soit $\mathcal{A}$ une catégorie abélienne.
Considérons un diagramme dont les flèches pleines sont données
$$
\xymatrix{
K^\bullet \ar[r]_\alpha \ar[d]_\gamma & L^\bullet \ar@{-->}[dl]^\beta \\
I^\bullet
}
$$
de complexes de $\text{Fil}^f(\mathcal{A})$. Supposons que
$K^\bullet$, $L^\bullet$ et $I^\bullet$ soient bornés inférieurement et que
chaque $I^n$ soit un objet injectif filtré.
Supposons aussi que $\alpha$ soit un quasi-isomorphisme filtré.
\begin{enumerate}
\item Il existe un morphisme de complexes $\beta$ rendant le diagramme
commutatif à homotopie près.
\item Si $\alpha$ est un monomorphisme strict en chaque degré,
on peut trouver un $\beta$ rendant le diagramme commutatif.
\end{enumerate}
\end{exercise}

\begin{exercise}
\label{exercise-acyclic-is-zero}
Soit $\mathcal{A}$ une catégorie abélienne.
Soient $K^\bullet$, $K^\bullet$ des complexes de $\text{Fil}^f(\mathcal{A})$.
Supposons que
\begin{enumerate}
\item $K^\bullet$ soit borné inférieurement et acyclique filtré, et
\item $I^\bullet$ soit borné inférieurement et formé d'objets injectifs filtrés.
\end{enumerate}
Alors tout morphisme $K^\bullet \to I^\bullet$ est homotope à zéro.
\end{exercise}

\begin{exercise}
\label{exercise-morphisms-equal-up-to-homotopy}
Soit $\mathcal{A}$ une catégorie abélienne.
Considérons un diagramme dont les flèches pleines sont données
$$
\xymatrix{
K^\bullet \ar[r]_\alpha \ar[d]_\gamma & L^\bullet \ar@{-->}[dl]^{\beta_i} \\
I^\bullet
}
$$
de complexes de $\text{Fil}^f(\mathcal{A})$.
Supposons que $K^\bullet$, $L^\bullet$ et $I^\bullet$ soient bornés inférieurement et que
chaque $I^n$ soit un objet injectif
filtré. Supposons aussi que $\alpha$ soit un quasi-isomorphisme filtré.
Deux morphismes quelconques $\beta_1, \beta_2$ rendant le diagramme commutatif
à homotopie près sont homotopes.
\end{exercise}







\section{Fonctions régulières}
\label{section-regular-functions}

\begin{exercise}
\label{exercise-extra-function}
Considérons la courbe affine $X$ définie par l'équation
$t^2 = s^5 + 8$ dans $\mathbf{C}^2$ muni des coordonnées $s, t$.
Soit $x \in X$ le point de coordonnées $(1, 3)$.
Posons $U = X \setminus \{x\}$. Montrer qu'il existe une
fonction régulière sur $U$ qui ne soit pas la restriction
d'une fonction régulière sur $\mathbf{C}^2$, autrement dit qui ne soit pas
la restriction d'un polynôme en $s$ et $t$ à $U$.
\end{exercise}

\begin{exercise}
\label{exercise-no-extra-function}
Soit $n \geq 2$. Soit $E \subset \mathbf{C}^n$ une partie finie.
Montrer que toute fonction régulière sur $\mathbf{C}^n \setminus E$
est un polynôme.
\end{exercise}

\begin{exercise}
\label{exercise-cone}
Soit $X \subset \mathbf{C}^n$ une variété affine. Disons que $X$ est un
{\it cône} si $x = (a_1, \ldots, a_n) \in X$ et $\lambda \in \mathbf{C}$
entraînent $(\lambda a_1, \ldots, \lambda a_n) \in X$.
Bien entendu, si $\mathfrak p \subset \mathbf{C}[x_1, \ldots, x_n]$
est un idéal premier engendré par des polynômes homogènes en $x_1, \ldots, x_n$,
alors la variété affine $X = V(\mathfrak p) \subset \mathbf{C}^n$ est un cône.
Montrer que, réciproquement, l'idéal premier
$\mathfrak p \subset \mathbf{C}[x_1, \ldots, x_n]$
correspondant à un cône peut être engendré par des polynômes homogènes
en $x_1, \ldots, x_n$.
\end{exercise}

\begin{exercise}
\label{exercise-extra-function-cone}
Donner un exemple d'une variété affine $X \subset \mathbf{C}^n$
qui soit un cône (voir l'exercice \ref{exercise-cone})
et d'une fonction régulière $f$ sur $U = X \setminus \{(0, \ldots, 0)\}$
qui ne soit pas la restriction d'une fonction polynomiale
sur $\mathbf{C}^n$.
\end{exercise}

\begin{exercise}
\label{exercise-regular-functions}
Dans cet exercice, nous cherchons à comprendre le comportement des fonctions régulières
sur des corps non algébriquement clos. Soit $k$ un corps.
Soit $Z \subset k^n$ une partie localement fermée pour la topologie de Zariski, c'est-à-dire qu'il
existe des idéaux $I \subset J \subset k[x_1, \ldots, x_n]$ tels que
$$
Z = \{a \in k^n \mid
f(a) = 0\ \forall\ f \in I,\ \exists\ g \in J,\ g(a) \not = 0\}.
$$
Une fonction $\varphi : Z \to k$ est dite {\it régulière} si, pour tout
$z \in Z$, il existe un voisinage ouvert de Zariski $z \in U \subset Z$
et des polynômes $f, g \in k[x_1, \ldots, x_n]$ tels que
$g(u) \not = 0$ pour tout $u \in U$ et que
$\varphi(u) = f(u)/g(u)$ pour tout $u \in U$.
\begin{enumerate}
\item Si $k = \bar k$ et $Z = k^n$, montrer que les fonctions régulières sont
données par des polynômes. (Ne traiter ceci que si vous n'avez pas déjà vu cet
argument.)
\item Si $k$ est fini, montrer que (a) toute fonction $\varphi$ est régulière,
(b) l'anneau des fonctions régulières est de dimension finie sur $k$.
(Si vous le souhaitez, vous pouvez prendre $Z = k^n$, voire $n = 1$.)
\item Si $k = \mathbf{R}$, donner un exemple d'une fonction régulière sur
$Z = \mathbf{R}$ qui ne soit pas donnée par un polynôme.
\item Si $k = \mathbf{Q}_p$, donner un exemple d'une fonction régulière sur
$Z = \mathbf{Q}_p$ qui ne soit pas donnée par un polynôme.
\end{enumerate}
\end{exercise}






\section{Faisceaux}
\label{section-sheaves}

\noindent
Un morphisme $f : X \to Y$ dans une catégorie $\mathcal{C}$ est un {\it monomorphisme}
si, pour tout couple de morphismes $a, b : W \to X$, on a
$f \circ a = f \circ b \Rightarrow a = b$. Dans la catégorie
des ensembles, un monomorphisme est une application injective.

\begin{exercise}
\label{exercise-mono-sheaves-sets}
Démontrer soigneusement qu'un morphisme de faisceaux d'ensembles est un monomorphisme
(dans la catégorie des faisceaux d'ensembles) si et seulement si les applications induites
sur toutes les fibres sont injectives.
\end{exercise}

\noindent
Un morphisme $f : X \to Y$ dans une catégorie $\mathcal{C}$ est un {\it isomorphisme}
s'il existe un morphisme $g : Y \to X$ tel que $f \circ g = \text{id}_Y$
et $g \circ f = \text{id}_X$. Dans la catégorie des ensembles,
un isomorphisme est une application bijective.

\begin{exercise}
\label{exercise-isomorphism-sheaves-sets}
Démontrer soigneusement qu'un morphisme de faisceaux d'ensembles est un isomorphisme
(dans la catégorie des faisceaux d'ensembles) si et seulement si les applications induites
sur toutes les fibres sont bijectives.
\end{exercise}

\noindent
Un morphisme $f : X \to Y$ dans une catégorie $\mathcal{C}$ est un {\it épimorphisme}
si, pour tout couple de morphismes $a, b : Y \to Z$, on a
$a \circ f = b \circ f \Rightarrow a = b$. Dans la
catégorie des ensembles, un épimorphisme est une application surjective.

\begin{exercise}
\label{exercise-epi-sheaves-sets}
Démontrer soigneusement qu'un morphisme de faisceaux d'ensembles est un épimorphisme
(dans la catégorie des faisceaux d'ensembles) si et seulement si les applications induites
sur toutes les fibres sont surjectives.
\end{exercise}

\begin{exercise}
\label{exercise-adjoint-push-pull}
Soit $f : X \to Y$ une application d'espaces topologiques.
Montrer que l'image directe $f_\ast$ et l'image réciproque $f^{-1}$ des faisceaux d'{\bf ensembles}
forment un couple de foncteurs adjoints.
\end{exercise}

\begin{exercise}
\label{exercise-j-shriek}
Soit $j : U \to X$ une immersion ouverte. Montrer que
\begin{enumerate}
\item L'image réciproque $j^{-1} : \Sh(X) \to \Sh(U)$ admet un adjoint à gauche
$j_{!} : \Sh(U) \to \Sh(X)$ appelé {\it prolongement par le vide}.
\item Caractériser les fibres de $j_{!}({\mathcal G})$ pour
$\mathcal{G} \in \Sh(U)$.
\item L'image réciproque $j^{-1} : \textit{Ab}(X) \to \textit{Ab}(U)$ admet un adjoint à gauche
$j_{!} : \textit{Ab}(U) \to \textit{Ab}(X)$ appelé {\it prolongement par zéro}.
\item Caractériser les fibres de $j_{!}({\mathcal G})$ pour
$\mathcal{G} \in \textit{Ab}(U)$.
\end{enumerate}
Observer que le prolongement par zéro diffère du prolongement par le
vide !
\end{exercise}

\begin{exercise}
\label{exercise-not-locally-generated-by-sections}
Soit $X = \mathbf{R}$ muni de la topologie usuelle.
Soit $\mathcal{O}_X = \underline{\mathbf{Z}/2\mathbf{Z}}_X$.
Soit $i : Z = \{0\} \to X$ l'inclusion et soit
$\mathcal{O}_Z = \underline{\mathbf{Z}/2\mathbf{Z}}_Z$.
Démontrer les assertions suivantes (les trois premières découlent des définitions, mais, si
celles-ci ne vous sont pas claires, il convient de les expliciter) :
\begin{enumerate}
\item $i_*\mathcal{O}_Z$ est un faisceau gratte-ciel.
\item Il existe un morphisme surjectif canonique
$\underline{\mathbf{Z}/2\mathbf{Z}}_X \to
i_*\underline{\mathbf{Z}/2\mathbf{Z}}_Z$.
Noter $\mathcal{I} \subset \mathcal{O}_X$ le noyau de ce morphisme.
\item $\mathcal{I}$ est un faisceau d'idéaux de $\mathcal{O}_X$.
\item Le faisceau $\mathcal{I}$ sur $X$ ne peut pas être localement engendré
par des sections (au sens de
Modules, définition \ref{modules-definition-locally-generated}.)
\end{enumerate}
\end{exercise}

\begin{exercise}
\label{exercise-quotient-j-shriek-Z}
Soit $X$ un espace topologique.
Soit ${\mathcal F}$ un faisceau abélien sur $X$. Montrer
que ${\mathcal F}$ est le quotient d'une somme directe (éventuellement très grande)
de faisceaux dont tous les termes sont de la forme
$$
j_{!}(\underline{{\mathbf Z}}_U)
$$
où $U \subset X$ est ouvert et où $\underline{{\mathbf Z}}_U$ désigne le
faisceau constant de valeur ${\mathbf Z}$ sur $U$.
\end{exercise}

\begin{remark}
\label{remark-direct-sum-stalk-abelian}
Soit $X$ un espace topologique.
Dans la catégorie des faisceaux abéliens, la somme directe
d'une famille de faisceaux $\{{\mathcal F}_i\}_{i\in I}$ est le faisceau associé au
préfaisceau $U \mapsto \oplus {\mathcal F}_i(U)$. Par conséquent, la fibre de
la somme directe en un point $x$ est la somme directe des fibres des
${\mathcal F}_i$ en $x$.
\end{remark}

\begin{exercise}
\label{exercise-product-over-points}
Soit $X$ un espace topologique. Supposons donnée une famille de
groupes abéliens $A_x$ indexée par $x \in X$. Montrer que la règle
$U \mapsto \prod_{x \in U} A_x$, munie des applications de restriction évidentes,
définit un faisceau $\mathcal{G}$ de groupes abéliens. Montrer, à l'aide d'un exemple,
que l'on n'a généralement pas $\mathcal{G}_x = A_x$ pour $x \in X$.
\end{exercise}

\begin{exercise}
\label{exercise-modified-product-over-points}
Soient $X$, $A_x$, $\mathcal{G}$ comme dans
l'exercice \ref{exercise-product-over-points}.
Soit $\mathcal{B}$ une base de la topologie de $X$, voir
Topologie, définition \ref{topology-definition-base}.
Pour $U \in \mathcal{B}$, soit $A_U$ un sous-groupe
$A_U \subset \mathcal{G}(U) = \prod_{x \in U} A_x$. Supposons que, si
$U \subset V$ avec $U, V \in \mathcal{B}$, les applications de restriction
envoient $A_V$ dans $A_U$. Pour tout ouvert $U \subset X$, posons
$$
\mathcal{F}(U) =
\left\{
(s_x)_{x \in U}
\middle|
\begin{matrix}
\text{ pour tout }x\text{ dans }U\text{ il existe } V \in \mathcal{B} \\
x \in V \subset U\text{ tel que } (s_y)_{y \in V} \in A_V
\end{matrix}
\right\}
$$
Montrer que $\mathcal{F}$ définit un faisceau de groupes abéliens sur $X$.
Montrer, à l'aide d'un exemple, que l'on n'a généralement pas
$\mathcal{F}(U) = A_U$ pour $U \in \mathcal{B}$.
\end{exercise}

\begin{exercise}
\label{exercise-exact-but-not-a-stalk-functor}
Donner un exemple d'un espace topologique $X$ et d'un foncteur
$$
F : \Sh(X) \longrightarrow \textit{Ensembles}
$$
qui soit exact (c'est-à-dire qui commute aux produits finis et aux égaliseurs, ainsi
qu'aux coproduits finis et aux coégaliseurs ; voir Catégories, section
\ref{categories-section-exact-functor}), mais pour lequel il n'existe aucun point $x \in X$
tel que $F$ soit isomorphe au foncteur fibre
$\mathcal{F} \mapsto \mathcal{F}_x$.
\end{exercise}




\section{Schémas}
\label{section-schemes}

\noindent
Soit $LRS$ la catégorie des espaces localement annelés.
Un schéma affine est un objet de $LRS$ isomorphe dans $LRS$ à
une paire de la forme $(\Spec(A), {\mathcal O}_{\Spec(A)})$.
Un schéma est un
objet $(X, {\mathcal O}_X)$ de $LRS$ tel que tout point $x\in X$
possède un voisinage ouvert $U \subset X$ pour lequel la paire
$(U, {\mathcal O}_X|_U)$ est un schéma affine.

\begin{exercise}
\label{exercise-one-point}
Trouver un espace localement annelé ayant $1$ point qui ne soit pas un schéma.
\end{exercise}

\begin{exercise}
\label{exercise-two-points}
Supposons que $X$ soit un schéma dont l'espace topologique
sous-jacent possède 2 points. Montrer que $X$ est un schéma affine.
\end{exercise}

\begin{exercise}
\label{exercise-discrete-finite-set-points}
Supposons que $X$ soit un schéma dont l'espace topologique sous-jacent est un
ensemble discret fini. Montrer que $X$ est un schéma affine.
\end{exercise}

\begin{exercise}
\label{exercise-three-points}
Montrer qu'il existe un schéma non affine possédant trois points.
\end{exercise}

\begin{exercise}
\label{exercise-quasi-compact-closed-point}
Supposons que $X$ soit un schéma quasi-compact non vide.
Montrer que $X$ possède un point fermé.
\end{exercise}

\begin{remark}
\label{remark-open-immersion}
Si $(X, {\mathcal O}_X)$ est un espace annelé et si $U \subset X$
est une partie ouverte, alors $(U, {\mathcal O}_X|_U)$ est un espace annelé. Notation :
${\mathcal O}_U = {\mathcal O}_X|_U$. Il existe un morphisme canonique
d'espaces annelés
$$
j : (U, {\mathcal O}_U) \longrightarrow (X, {\mathcal O}_X).
$$
Si $(X, {\mathcal O}_X)$ est un espace localement annelé, il en est de même de
$(U, {\mathcal O}_U)$, et
$j$ est un morphisme d'espaces localement annelés. Si $(X, {\mathcal O}_X)$
est un schéma,
il en est de même de $(U, {\mathcal O}_U)$, et $j$ est un morphisme de schémas. On dit
que
$(U, {\mathcal O}_U)$ est un {\it sous-schéma ouvert} de $(X, {\mathcal O}_X)$
et que
$j$ est une {\it immersion ouverte}. Plus généralement, tout morphisme
$j' : (V, {\mathcal O}_V) \to (X, {\mathcal O}_X)$ qui est {\it isomorphe}
à un
morphisme $j : (U, {\mathcal O}_U) \to (X, {\mathcal O}_X)$ comme ci-dessus est
appelé une
immersion ouverte.
\end{remark}

\begin{exercise}
\label{exercise-open-affine-not-affine}
Donner un exemple d'un schéma affine $(X, {\mathcal O}_X)$
et d'un ouvert $U \subset X$ tels que $(U, {\mathcal O}_X|U)$ ne soit pas un
schéma affine.
\end{exercise}

\begin{exercise}
\label{exercise-morphism-does-not-extend}
Donner un exemple d'une paire de schémas affines
$(X, {\mathcal O}_X)$, $(Y, {\mathcal O}_Y)$,
d'un sous-schéma ouvert $(U, {\mathcal O}_X|_U)$
de $X$ et d'un morphisme de schémas
$(U, {\mathcal O}_X|_U) \to (Y, {\mathcal O}_Y)$
qui ne se prolonge pas en un morphisme de schémas
$(X, {\mathcal O}_X) \to (Y, {\mathcal O}_Y)$.
\end{exercise}

\begin{exercise}
\label{exercise-closed-subscheme-does-not-extend}
(Cet exercice est assez difficile.)
Donner un exemple d'un schéma $X$, d'un sous-schéma ouvert $U \subset X$
et d'un sous-schéma fermé $Z \subset U$ tel que $Z$ ne se prolonge pas
en un sous-schéma fermé de $X$.
\end{exercise}

\begin{exercise}
\label{exercise-not-morphism-schemes}
Donner un exemple d'un schéma $X$, d'un corps $K$ et d'un
morphisme d'espaces annelés $\Spec(K) \to X$ qui
N'EST PAS un morphisme de schémas.
\end{exercise}

\begin{exercise}
\label{exercise-just-kidding}
Faire tous les exercices de \cite[chapitre II]{H},
sections 1 et 2...\ \ Je plaisante !
\end{exercise}

\begin{definition}
\label{definition-integral}
Un schéma $X$ est dit {\it intègre} si $X$ est non vide et si,
pour tout ouvert affine non vide $U \subset X$, l'anneau
$\Gamma(U, \mathcal{O}_X) = \mathcal{O}_X(U)$ est intègre.
\end{definition}

\begin{exercise}
\label{exercise-morphism-integral-schemes-surjective-stalks-not-closed}
Donner un exemple d'un morphisme de schémas {\it intègres}
$f : X \to Y$ tel que les homomorphismes induits ${\mathcal O}_{Y, f(x)}
\to {\mathcal O}_{X, x}$ soient surjectifs pour tout $x\in X$, mais que $f$
ne soit pas une immersion fermée.
\end{exercise}

\begin{exercise}
\label{exercise-fibre-product-affines-not-affine}
Donner un exemple d'un produit fibré $X \times_S Y$ tel que $X$ et $Y$
soient affines, mais que $X \times_S Y$ ne le soit pas.
\end{exercise}

\begin{remark}
\label{remark-separated-base-fibre-product-affines-affine}
Cela ne peut en fait pas se produire lorsque $S$ est séparé. Savez-vous pourquoi ?
\end{remark}

\begin{exercise}
\label{exercise-not-geometrically-integral}
Donner un exemple d'un schéma
$V$ qui soit un schéma intègre de dimension 1 et de type fini
sur ${\mathbf Q}$ tel que
$\Spec({\mathbf C}) \times_{\Spec({\mathbf Q})} V$
ne soit pas intègre.
\end{exercise}

\begin{exercise}
\label{exercise-not-geometrically-reduced}
Donner un exemple d'un schéma
$V$ qui soit un schéma intègre de dimension 1 et de type fini
sur un corps $k$ tel que $\Spec(k') \times_{\Spec(k)} V$
ne soit pas réduit pour une certaine extension finie de corps $k'/k$.
\end{exercise}

\begin{remark}
\label{remark-affine-dimension}
Si votre schéma est affine, sa dimension est
égale à la dimension de Krull de l'anneau sous-jacent. Vous pouvez donc
utiliser les résultats du semestre dernier pour calculer la dimension.
\end{remark}






\section{Morphismes}
\label{section-morphisms}

\noindent
Une question importante est de savoir, étant donné un morphisme $\pi : X \to S$,
si ce morphisme possède une section ou une section rationnelle.
Voici quelques exercices à titre d'exemple.

\begin{exercise}
\label{exercise-no-section}
Considérons le morphisme de schémas
$$
\pi :
X = \Spec(\mathbf{C}[x, t, 1/xt])
\longrightarrow
S = \Spec(\mathbf{C}[t]).
$$
\begin{enumerate}
\item Montrer qu'il n'existe aucun morphisme $\sigma : S \to X$
tel que $\pi \circ \sigma = \text{id}_S$.
\item Montrer qu'il existe un ouvert non vide $U \subset S$ et
un morphisme $\sigma : U \to X$ tel que $\pi \circ \sigma = \text{id}_U$.
\end{enumerate}
\end{exercise}

\begin{exercise}
\label{exercise-no-rational-section}
Considérons le morphisme de schémas
$$
\pi :
X = \Spec(\mathbf{C}[x, t]/(x^2 + t))
\longrightarrow
S = \Spec(\mathbf{C}[t]).
$$
Montrer qu'il n'existe aucun ouvert non vide $U \subset S$ et
aucun morphisme $\sigma : U \to X$ tel que $\pi \circ \sigma = \text{id}_U$.
\end{exercise}

\begin{exercise}
\label{exercise-has-rational-section}
Soient $A, B, C \in \mathbf{C}[t]$ des polynômes non nuls.
Considérons le morphisme de schémas
$$
\pi :
X = \Spec(\mathbf{C}[x, y, t]/(A + Bx^2 + Cy^2))
\longrightarrow
S = \Spec(\mathbf{C}[t]).
$$
Montrer qu'il existe un ouvert non vide $U \subset S$ et
un morphisme $\sigma : U \to X$ tel que $\pi \circ \sigma = \text{id}_U$.
(Indication : symboliquement, écrire $x = X/Z$, $y = Y/Z$ avec
$X, Y, Z \in \mathbf{C}[t]$ de degré $\leq d$, pour un certain $d$,
et déterminer la condition pour que cela résolve l'équation.
Montrer ensuite, au moyen de la théorie de la dimension, que si $d >> 0$, on peut trouver
des $X, Y, Z$ non nuls qui résolvent l'équation.)
\end{exercise}

\begin{remark}
\label{remark-tsen}
L'exercice \ref{exercise-has-rational-section}
est un cas particulier du « théorème de Tsen ».
L'exercice \ref{exercise-no-section-curve} montre que la
méthode est limitée aux équations de faible degré (aux coniques lorsque la base et
la fibre sont de dimension 1).
\end{remark}

\begin{exercise}
\label{exercise-no-section-curve}
Considérons le morphisme de schémas
$$
\pi :
X = \Spec(\mathbf{C}[x, y, t]
/(1 + t x^3 + t^2 y^3))
\longrightarrow
S = \Spec(\mathbf{C}[t])
$$
Montrer qu'il n'existe aucun ouvert non vide $U \subset S$ et
aucun morphisme $\sigma : U \to X$ tel que $\pi \circ \sigma = \text{id}_U$.
\end{exercise}

\begin{exercise}
\label{exercise-no-section-surface}
Considérons les schémas
$$
X = \Spec(\mathbf{C}[\{x_i\}_{i = 1}^{8}, s, t]
/(1 + s x_1^3 + s^2 x_2^3 +
t x_3^3 + st x_4^3 + s^2t x_5^3 +
t^2 x_6^3 + st^2 x_7^3 + s^2t^2 x_8^3))
$$
et
$$
S = \Spec(\mathbf{C}[s, t])
$$
ainsi que le morphisme de schémas
$$
\pi : X \longrightarrow S
$$
Montrer qu'il n'existe aucun ouvert non vide $U \subset S$ et
aucun morphisme $\sigma : U \to X$ tel que $\pi \circ \sigma = \text{id}_U$.
\end{exercise}

\begin{exercise}
\label{exercise-for-number-theorists}
(Pour les théoriciens des nombres.) Donner un exemple d'un sous-schéma fermé
$$
Z \subset \Spec\left({\mathbf Z}[x, \frac{1 }{ x(x-1)(2x-1)}]\right)
$$
tel que le morphisme $Z \to \Spec({\mathbf Z})$ soit fini
et surjectif.
\end{exercise}

\begin{exercise}
\label{exercise-quasi-section}
Si la théorie des nombres ne vous plaît pas, vous pouvez essayer la
variante où l'on considère
$$
\Spec\left({\mathbf F}_p[t, x, \frac{1 }{ x(x-t)(tx-1)}]\right)
\longrightarrow
\Spec({\mathbf F}_p[t])
$$
et chercher un sous-schéma fermé du schéma supérieur
dont le morphisme vers celui du bas soit fini et surjectif. (Il existe une
raison théorique de supposer ici le corps de base fini, bien que
cela ne soit peut-être pas nécessaire dans ce cas particulier.)
\end{exercise}

\begin{remark}
\label{remark-interpretation-skolem-noether}
Dans les exercices \ref{exercise-for-number-theorists} et
\ref{exercise-quasi-section}, on considère un morphisme $f : X \to S$
où $S$ est le spectre d'un anneau de Dedekind $A$, le morphisme $f$ étant plat et
surjectif, à fibre générique géométriquement irréductible. Dans les deux cas,
on conclut qu'il existe un morphisme fini et surjectif $S' \to S$
tel qu'après changement de base, $X_{S'} \to S'$ possède une section.
Ce résultat vaut lorsque $A$ est excellent, que ses corps résiduels
aux idéaux maximaux sont des extensions algébriques de corps finis,
et que $\Pic(A')$ est de torsion lorsque $A'$ est la clôture intégrale
de $A$ dans une extension finie de son corps des fractions; voir \cite{MB1}.
C'est par exemple le cas si $A = \mathbf{Z}$ ou $A = \mathbf{F}_p[x]$, ou pour une extension
finie de l'un de ces anneaux. Cependant, il
existe un anneau de Dedekind $A$ dont les corps résiduels aux
idéaux maximaux sont finis et dont le groupe de Picard possède un élément non de torsion; voir \cite{Goldman},
et le résultat est faux pour le spectre d'un tel anneau.
L'exercice \ref{exercise-no-quasi-section} donne un exemple géométrique
où cela ne fonctionne pas.
\end{remark}

\begin{exercise}
\label{exercise-no-quasi-section}
Démontrer qu'il existe un $f \in \mathbf{C}[x, t]$ qui n'est divisible
par $t - \alpha$ pour aucun $\alpha \in \mathbf{C}$ et tel qu'il
n'existe aucun $Z \subset \Spec(\mathbf{C}[x, t, 1/f])$
dont le morphisme vers $\Spec(\mathbf{C}[t])$ soit fini et surjectif.
(Je pense que $f(x, t) = (xt - 2)(x - t + 3)$ convient. Il n'est pas si facile,
me semble-t-il, de montrer que ce candidat possède la propriété requise.)
\end{exercise}

\begin{exercise}
\label{exercise-finite}
Soit $A \to B$ un homomorphisme d'anneaux de type fini. Supposons que
$\Spec(B) \to \Spec(A)$ se factorise par une immersion fermée
$\Spec(B) \to \mathbf{P}^n_A$ pour un certain $n$.
Démontrer que $A \to B$ est un homomorphisme fini, c'est-à-dire que
$B$ est un $A$-module de type fini.
Indication : si $A$ est noethérien (supposez-le simplement),
on peut utiliser le fait que $H^i(Z, \mathcal{O}_Z)$, pour $i \in \mathbf{Z}$,
est un $A$-module de type fini pour tout sous-schéma fermé
$Z \subset \mathbf{P}^n_A$.
\end{exercise}

\begin{exercise}
\label{exercise-projective-finite}
Soit $k$ un corps algébriquement clos.
Soit $f : X \to Y$ un morphisme de variétés
projectives tel que $f^{-1}(\{y\})$ soit fini pour
tout point fermé $y \in Y$.
Démontrer que $f$ est fini en tant que morphisme de schémas.
Indications : (a) la finitude est une propriété locale,
(b) essayer de se ramener à l'exercice \ref{exercise-finite}, et
(c) utiliser une immersion fermée $X \to \mathbf{P}^n_k$ pour obtenir
une immersion fermée $X \to \mathbf{P}^n_Y$ au-dessus de $Y$.
\end{exercise}






\section{Espaces tangents}
\label{section-tangent-space}

\begin{definition}
\label{definition-dual-numbers}
Pour tout anneau $R$, nous notons $R[\epsilon]$ l'anneau
des {\it nombres duaux}. Comme $R$-module, il est libre de
base $1$, $\epsilon$. La structure d'anneau s'obtient en posant
$\epsilon^2 = 0$.
\end{definition}

\begin{exercise}
\label{exercise-tangent-space-Zariski}
Soit $f : X \to S$ un morphisme de schémas.
Soit $x \in X$ un point, et soit $s = f(x)$.
Considérer le diagramme commutatif suivant, où les flèches pleines sont données
$$
\xymatrix{
\Spec(\kappa(x)) \ar[r] \ar[dr] \ar@/^1pc/[rr] &
\Spec(\kappa(x)[\epsilon]) \ar@{.>}[r] \ar[d]&
X \ar[d] \\
&
\Spec(\kappa(s)) \ar[r] &
S
}
$$
où la flèche courbe est le morphisme canonique de
$\Spec(\kappa(x))$ vers $X$.
Si $\kappa(x) = \kappa(s)$, montrer que l'ensemble des flèches
pointillées qui rendent le diagramme commutatif est en bijection
avec l'ensemble des applications linéaires
$$
\Hom_{\kappa(x)}(
\frac{\mathfrak m_x}{\mathfrak m_x^2 + \mathfrak m_s\mathcal{O}_{X, x}},
\kappa(x))
$$
Autrement dit, décrire une telle bijection.
(Cela vaut plus généralement si $\kappa(x) \supset \kappa(s)$ est une
extension algébrique séparable.)
\end{exercise}

\begin{definition}
\label{definition-tangent-space}
Soit $f : X \to S$ un morphisme de schémas.
Soit $x \in X$. Nous appelons l'ensemble des flèches pointillées
de l'exercice \ref{exercise-tangent-space-Zariski}
l'{\it espace tangent de $X$ sur $S$}
et nous le notons $T_{X/S, x}$. Un élément de cet
espace est appelé {\it vecteur tangent} à $X/S$ en $x$.
\end{definition}

\begin{exercise}
\label{exercise-simple-push-out}
Pour tout corps $K$, démontrer que le diagramme
$$
\xymatrix{
\Spec(K) \ar[r] \ar[d] & \Spec(K[\epsilon_1]) \ar[d] \\
\Spec(K[\epsilon_2]) \ar[r] &
\Spec(K[\epsilon_1, \epsilon_2]/(\epsilon_1\epsilon_2))
}
$$
est cocartésien dans la catégorie des schémas.
(Ici $\epsilon_i^2 = 0$ comme précédemment.)
\end{exercise}

\begin{exercise}
\label{exercise-tangent-space-vectors-space}
Soit $f : X \to S$ un morphisme de schémas.
Soit $x \in X$. Définir l'addition des vecteurs tangents,
à l'aide de l'exercice \ref{exercise-simple-push-out}
et d'un morphisme convenable
$$
\Spec(K[\epsilon])
\longrightarrow
\Spec(K[\epsilon_1, \epsilon_2]/(\epsilon_1\epsilon_2)).
$$
Définir de même la multiplication des vecteurs tangents par les scalaires (c'est plus simple).
Montrer que vos constructions munissent $T_{X/S, x}$ d'une structure
d'espace vectoriel sur $\kappa(x)$.
\end{exercise}

\begin{exercise}
\label{exercise-compute-TS}
Soit $k$ un corps. Considérer
le morphisme structural $f : X = \mathbf{A}^1_k \to \Spec(k) = S$.
\begin{enumerate}
\item Soit $x \in X$ un point fermé. Quelle est la dimension de
$T_{X/S, x}$ ?
\item Soit $\eta \in X$ le point générique. Quelle est la dimension
de $T_{X/S, \eta}$ ?
\item Considérer maintenant $X$ comme un schéma sur $\Spec(\mathbf{Z})$.
Quelles sont les dimensions de $T_{X/\mathbf{Z}, x}$ et de
$T_{X/\mathbf{Z}, \eta}$ ?
\end{enumerate}
\end{exercise}

\begin{remark}
\label{remark-tangent-space-relative}
L'exercice \ref{exercise-compute-TS} explique pourquoi il est nécessaire
de considérer l'espace tangent à $X$ relativement à $S$ pour obtenir une bonne notion.
\end{remark}

\begin{exercise}
\label{exercise-compute-TS-field}
Considérer le morphisme de schémas
$$
f : X = \Spec(\mathbf{F}_p(t))
\longrightarrow
\Spec(\mathbf{F}_p(t^p)) = S
$$
Calculer l'espace tangent de $X/S$ à l'unique point de $X$.
N'est-ce pas étrange ? Que se passe-t-il, à votre avis, si l'on considère le morphisme
de schémas correspondant à $\mathbf{F}_p[t^p] \to \mathbf{F}_p[t]$ ?
\end{exercise}

\begin{exercise}
\label{exercise-compute-TS-cusp}
Soit $k$ un corps.
Calculer l'espace tangent de $X/k$ au point $x = (0, 0)$
où $X = \Spec(k[x, y]/(x^2 - y^3))$.
\end{exercise}

\begin{exercise}
\label{exercise-map-tangent-spaces}
Soit $f : X \to Y$ un morphisme de schémas au-dessus de $S$.
Soit $x \in X$ un point. Poser $y = f(x)$.
Supposons que l'application naturelle $\kappa(y) \to \kappa(x)$ soit
bijective. Montrer, à partir de la définition,
que $f$ induit une application linéaire naturelle
$$
\text{d}f : T_{X/S, x} \longrightarrow T_{Y/S, y}.
$$
Identifier celle-ci à l'application obtenue sur les anneaux locaux à l'aide de
l'exercice \ref{exercise-tangent-space-Zariski} lorsque $\kappa(x) = \kappa(s)$.
\end{exercise}

\begin{exercise}
\label{exercise-Jacobian}
Soit $k$ un corps algébriquement clos.
Soit
\begin{eqnarray*}
f : \mathbf{A}_k^n & \longrightarrow & \mathbf{A}^m_k \\
(x_1, \ldots, x_n) & \longmapsto & (f_1(x_i), \ldots, f_m(x_i))
\end{eqnarray*}
un morphisme de schémas au-dessus de $k$. Il est défini par
$m$ polynômes $f_1, \ldots, f_m$ en $n$ variables.
Considérer la matrice
$$
A = \left( \frac{\partial f_j}{\partial x_i} \right)
$$
Soit $x \in \mathbf{A}^n_k$ un point fermé.
Poser $y =  f(x)$. Montrer que l'application sur les espaces tangents
$T_{\mathbf{A}^n_k/k, x} \to T_{\mathbf{A}^m_k/k, y}$
est donnée par la valeur de la matrice $A$ au point $x$.
\end{exercise}














\section{Faisceaux quasi-cohérents}
\label{section-quasi-coherent}

\begin{definition}
\label{definition-quasi-coherent}
Soit $X$ un schéma.
Un faisceau $\mathcal{F}$ de $\mathcal{O}_X$-modules est {\it quasi-cohérent}
si, pour tout ouvert affine $\Spec(R) = U \subset X$, la restriction
$\mathcal{F}|_U$ est de la forme $\widetilde M$ pour un certain $R$-module
$M$.
\end{definition}

\noindent
Il suffit de vérifier cette condition sur les membres d'un
recouvrement ouvert affine de $X$.
Voir Schémas, section \ref{schemes-section-quasi-coherent}
pour davantage de résultats.

\begin{definition}
\label{definition-specialization}
Soit $X$ un espace topologique. Soient $x, x' \in X$.
On dit que $x$ est une {\it spécialisation} de $x'$
si et seulement si $x \in \overline{\{x'\}}$.
\end{definition}

\begin{exercise}
\label{exercise-quasi-coherent-specialization-points}
Soit $X$ un schéma. Soient $x, x' \in X$. Soit $\mathcal{F}$
un faisceau quasi-cohérent de $\mathcal{O}_X$-modules.
Supposons que (a) $x$ soit une spécialisation de $x'$ et que (b)
$\mathcal{F}_{x'} \not = 0$. Montrer que $\mathcal{F}_x \not = 0$.
\end{exercise}

\begin{exercise}
\label{exercise-O-module-specialization-points}
Trouver un exemple d'un schéma $X$, de points $x, x' \in X$ et
d'un faisceau de $\mathcal{O}_X$-modules
$\mathcal{F}$ tels que (a) $x$ soit une spécialisation de $x'$ et que (b)
$\mathcal{F}_{x'} \not = 0$ et $\mathcal{F}_x = 0$.
\end{exercise}

\begin{definition}
\label{definition-Noetherian-scheme}
Un schéma $X$ est dit {\it localement noethérien} si et seulement si,
pour tout point $x \in X$, il existe un ouvert affine
$\Spec(R) = U \subset X$ tel que $R$ soit noethérien.
Un schéma est {\it noethérien} s'il est localement noethérien et quasi-compact.
\end{definition}

\noindent
Si $X$ est localement noethérien, tout ouvert affine de $X$
est le spectre d'un anneau noethérien, voir
Propriétés, lemme \ref{properties-lemma-locally-Noetherian}.

\begin{definition}
\label{definition-coherent}
Soit $X$ un schéma localement noethérien.
Soit $\mathcal{F}$ un faisceau quasi-cohérent de
$\mathcal{O}_X$-modules. On dit que $\mathcal{F}$ est {\it cohérent}
si, pour tout point $x \in X$, il existe un ouvert affine
$\Spec(R) = U \subset X$ tel que $\mathcal{F}|_U$
soit isomorphe à $\widetilde M$ pour un $R$-module $M$ de type fini.
\end{definition}

\begin{exercise}
\label{exercise-extend-quasi-coherent}
Soit $X = \Spec(R)$ un schéma affine.
\begin{enumerate}
\item Soit $f \in R$. Soit $\mathcal{G}$ un
faisceau quasi-cohérent de $\mathcal{O}_{D(f)}$-modules
sur le sous-schéma ouvert $D(f)$.
Montrer que $\mathcal{G} = \mathcal{F}|_{D(f)}$ pour un certain
faisceau quasi-cohérent de $\mathcal{O}_X$-modules
$\mathcal{F}$.
\item Soit $I \subset R$ un idéal.
Soit $i : Z \to X$ le sous-schéma fermé de $X$ correspondant
à $I$. Soit $\mathcal{G}$ un
faisceau quasi-cohérent de $\mathcal{O}_Z$-modules
sur le sous-schéma fermé $Z$.
Montrer que $\mathcal{G} = i^*\mathcal{F}$ pour un certain
faisceau quasi-cohérent de $\mathcal{O}_X$-modules $\mathcal{F}$.
(Pourquoi cela est-il trivial ?)
\item Supposons que $R$ soit noethérien.
Soit $f \in R$. Soit $\mathcal{G}$ un
faisceau cohérent de $\mathcal{O}_{D(f)}$-modules
sur le sous-schéma ouvert $D(f)$.
Montrer que $\mathcal{G} = \mathcal{F}|_{D(f)}$ pour un certain
faisceau cohérent de $\mathcal{O}_X$-modules $\mathcal{F}$.
\end{enumerate}
\end{exercise}

\begin{remark}
\label{remark-extend-off-open}
Si $U \to X$ est une immersion quasi-compacte, tout
faisceau quasi-cohérent sur $U$ est la restriction d'un
faisceau quasi-cohérent sur $X$.
Si $X$ est un schéma noethérien et si $U \subset X$ est ouvert,
tout faisceau cohérent sur $U$ est la restriction d'un
faisceau cohérent sur $X$.
L'exercice ci-dessus est bien sûr plus facile et ne devrait pas faire appel à ces résultats généraux.
\end{remark}
















\section{Proj et schémas projectifs}
\label{section-proj}

\begin{exercise}
\label{exercise-graded-ring-specified-result}
Donner des exemples d'anneaux gradués $S$ tels que
\begin{enumerate}
\item $\text{Proj}(S)$ soit affine et non vide, et
\item $\text{Proj}(S)$ soit intègre, non vide, mais non isomorphe
à ${\mathbf P}^n_A$ pour aucun $n\geq 0$ ni aucun anneau $A$.
\end{enumerate}
\end{exercise}

\begin{exercise}
\label{exercise-nonconstant-morphism-proj}
Donner un exemple de morphisme non constant
de schémas ${\mathbf P}^1_{\mathbf C} \to {\mathbf P}^5_{\mathbf C}$ au-dessus de
$\Spec({\mathbf C})$.
\end{exercise}

\begin{exercise}
\label{exercise-isomorphism-P1-conic}
Donner un exemple d'isomorphisme de schémas
$$
{\mathbf P}^1_{\mathbf C} \to
\text{Proj}({\mathbf C}[X_0, X_1, X_2]/(X_0^2 + X_1^2 + X_2^2))
$$
\end{exercise}

\begin{exercise}
\label{exercise-family-special-fibre-different}
Donner un exemple de morphisme de schémas
$f : X \to {\mathbf A}^1_{\mathbf C} = \Spec({\mathbf C}[T])$
tel que la fibre (schématique) $X_t$ de $f$ au-dessus de
$t \in {\mathbf A}^1_{\mathbf C}$ soit
(a) isomorphe à ${\mathbf P}^1_{\mathbf C}$ lorsque $t$
est un point fermé distinct de $0$, et
(b) non isomorphe à ${\mathbf P}^1_{\mathbf C}$ lorsque $t = 0$.
Nous appellerons $X_0$ la {\it fibre spéciale} du morphisme.
On peut le faire de très nombreuses façons. Essayer de donner des exemples qui
satisfassent (chacune des) contraintes supplémentaires suivantes (à moins
que cela ne soit impossible) :
\begin{enumerate}
\item Peut-on le faire avec une fibre spéciale projective ?
\item Peut-on le faire avec une fibre spéciale irréductible et projective ?
\item Peut-on le faire avec une fibre spéciale intègre et projective ?
\item Peut-on le faire avec une fibre spéciale lisse et projective ?
\item Peut-on le faire avec $f$ plat ?
Cela signifie simplement que, pour tout ouvert affine $\Spec(A) \subset X$,
l'homomorphisme d'anneaux induit $\mathbf{C}[t] \to A$ est plat, ce qui signifie ici
que tout polynôme non nul en $t$ est un non-diviseur de zéro dans $A$.
\item Peut-on le faire avec $f$ plat et projectif ?
\item Peut-on le faire avec $f$ plat et projectif et avec une fibre spéciale réduite ?
\item Peut-on le faire avec $f$ plat et projectif et avec une fibre spéciale irréductible ?
\item Peut-on le faire avec $f$ plat et projectif et avec une fibre spéciale intègre ?
\end{enumerate}
Que se passe-t-il, selon vous, lorsque l'on remplace
${\mathbf P}^1_{\mathbf C}$ par une autre variété sur ${\mathbf C}$ ?
(Cela peut devenir très difficile selon celle des variantes ci-dessus que l'on
demande.)
\end{exercise}

\begin{exercise}
\label{exercise-affine-onto-projective-space}
Soit $n \geq 1$ un entier positif quelconque.
Donner un exemple de morphisme surjectif
$X \to {\mathbf P}^n_{\mathbf C}$ avec $X$ affine.
\end{exercise}

\begin{exercise}
\label{exercise-morphism-proj}
Morphismes de $\text{Proj}$. Soient $R$ et $S$ des anneaux gradués.
Supposons donné un homomorphisme d'anneaux
$$
\psi : R \to S
$$
et un entier $e \geq 1$ tel que $\psi(R_d) \subset S_{de}$
pour tout $d \geq 0$. (Selon nos conventions, ce n'est pas un homomorphisme
d'anneaux gradués, sauf si $e = 1$.)
\begin{enumerate}
\item Pour quels éléments $\mathfrak p \in \text{Proj}(S)$ existe-t-il
un point correspondant bien défini dans $\text{Proj}(R)$ ? Autrement dit,
trouver un ouvert convenable $U \subset \text{Proj}(S)$ tel que $\psi$ définisse
une application continue $r_\psi : U \to \text{Proj}(R)$.
\item Donner un exemple où $U \not = \text{Proj}(S)$.
\item Donner un exemple où $U = \text{Proj}(S)$.
\item (Ne pas rédiger ceci.) Se convaincre que
l'application continue $U \to \text{Proj}(R)$ est canoniquement munie
d'un morphisme de faisceaux qui fait de $r_\psi$ un morphisme de schémas :
$$
\text{Proj}(S) \supset U \longrightarrow \text{Proj}(R).
$$
\item Que peut-on dire de ce morphisme si
$R = \bigoplus_{d \geq 0} S_{de}$ (considéré comme anneau gradué où
$S_e$, $S_{2e}$, etc. sont respectivement de degrés $1$, $2$, etc.) et si $\psi$
est l'inclusion ?
\end{enumerate}
\end{exercise}

\noindent
{\bf Notation.} Soit $R$ un anneau gradué comme ci-dessus et
soit $n \geq 0$ un entier. Soit $X = \text{Proj}(R)$. Il existe alors un unique
${\mathcal O}_X$-module quasi-cohérent ${\mathcal O}_X(n)$ sur $X$ tel que,
pour tout élément homogène $f \in R$ de degré positif, la restriction
${\mathcal O}_X |_{D_{+}(f)}$ soit le faisceau quasi-cohérent associé au
$R_{(f)} = (R_f)_0$-module $(R_f)_n$ ($ = $ les éléments homogènes de degré
$n$ dans $R_f = R[1/f]$). Voir \cite[page 116+]{H}. Noter qu'il existe
des morphismes naturels
$$
{\mathcal O}_X(n_1) \otimes_{{\mathcal O}_X} {\mathcal O}_X(n_2)
\longrightarrow
{\mathcal O}_X(n_1 + n_2)
$$

\begin{exercise}
\label{exercise-pathologies-proj}
Pathologies dans $\text{Proj}$.
Donner des exemples de $R$ comme ci-dessus tels que
\begin{enumerate}
\item ${\mathcal O}_X(1)$ ne soit pas un ${\mathcal O}_X$-module inversible.
\item ${\mathcal O}_X(1)$ soit inversible, mais que le
morphisme naturel ${\mathcal O}_X(1) \otimes_{{\mathcal O}_X} {\mathcal O}_X(1) \to
{\mathcal O}_X(2)$ ne soit PAS un isomorphisme.
\end{enumerate}
\end{exercise}

\begin{exercise}
\label{exercise-finitely-many-points-in-affine}
Soit $S$ un anneau gradué.
Soit $X = \text{Proj}(S)$.
Montrer que tout ensemble fini de points de $X$ est contenu dans un ouvert
affine standard.
\end{exercise}

\begin{exercise}
\label{exercise-prepare-glueing}
Soit $S$ un anneau gradué.
Soit $X = \text{Proj}(S)$.
Soient $Z, Z' \subset X$ deux sous-schémas fermés.
Soit $\varphi : Z \to Z'$ un isomorphisme.
Supposons que $Z \cap Z' = \emptyset$.
Montrer que, pour tout $z \in Z$, il existe un ouvert affine
$U \subset X$ tel que $z \in U$, $\varphi(z) \in U$ et
$\varphi(Z \cap U) = Z' \cap U$.
(Indication : utiliser l'exercice \ref{exercise-finitely-many-points-in-affine}
et un argument analogue à celui de
Schémas, lemme \ref{schemes-lemma-standard-open-two-affines}.)
\end{exercise}







\section{Morphismes de la droite projective}
\label{section-from-P1}

\noindent
Dans cette section, nous étudions les morphismes de $\mathbf{P}^1$
vers des schémas projectifs.

\begin{exercise}
\label{exercise-from-generic-point}
Soit $k$ un corps. Soit $k[t] \subset k(t)$
l'inclusion de l'anneau de polynômes dans son corps des fractions.
Soit $X$ un schéma de type fini sur $k$.
Montrer que, pour tout morphisme
$$
\varphi : \Spec(k(t)) \longrightarrow X
$$
au-dessus de $k$, il existe un élément non nul $f \in k[t]$ et un morphisme
$\psi : \Spec(k[t, 1/f]) \to X$ au-dessus de $k$
tels que $\varphi$ soit la composée
$$
\Spec(k(t)) \longrightarrow \Spec(k[t, 1/f]) \longrightarrow X
$$
\end{exercise}

\begin{exercise}
\label{exercise-from-generic-point-Pn}
Soit $k$ un corps. Soit $k[t] \subset k(t)$
l'inclusion de l'anneau de polynômes dans son corps des fractions.
Montrer que, pour tout morphisme
$$
\varphi : \Spec(k(t)) \longrightarrow \mathbf{P}^n_k
$$
au-dessus de $k$, il existe un morphisme $\psi : \Spec(k[t]) \to \mathbf{P}^n_k$
au-dessus de $k$ tel que $\varphi$ soit la composée
$$
\Spec(k(t)) \longrightarrow \Spec(k[t]) \longrightarrow \mathbf{P}^n_k
$$
Indication : l'image de $\varphi$ est contenue dans un ouvert standard $D_+(T_i)$
pour un certain $i$ ; montrer alors que l'on peut ``chasser les dénominateurs''.
\end{exercise}

\begin{exercise}
\label{exercise-from-generic-point-projective}
Soit $k$ un corps. Soit $k[t] \subset k(t)$
l'inclusion de l'anneau de polynômes dans son corps des fractions.
Soit $X$ un schéma projectif sur $k$.
Montrer que, pour tout morphisme
$$
\varphi : \Spec(k(t)) \longrightarrow X
$$
au-dessus de $k$, il existe un morphisme $\psi : \Spec(k[t]) \to X$
au-dessus de $k$ tel que $\varphi$ soit la composée
$$
\Spec(k(t)) \longrightarrow \Spec(k[t]) \longrightarrow X
$$
Indication : utiliser l'exercice \ref{exercise-from-generic-point-Pn}.
\end{exercise}

\begin{exercise}
\label{exercise-from-generic-point-P1-to-projective}
Soit $k$ un corps. Soit $X$ un schéma projectif sur $k$.
Soit $K$ le corps des fonctions rationnelles de $\mathbf{P}^1_k$ (voir
l'indication ci-dessous). Montrer que, pour tout morphisme
$$
\varphi : \Spec(K) \longrightarrow X
$$
au-dessus de $k$, il existe un morphisme $\psi : \mathbf{P}^1_k \to X$
au-dessus de $k$ tel que $\varphi$ soit la composée
$$
\Spec(k(t)) \longrightarrow \mathbf{P}^1_k \longrightarrow X
$$
Indication : utiliser l'exercice \ref{exercise-from-generic-point-projective}
pour chacun des deux ouverts du recouvrement affine
$\mathbf{P}^1_k = D_+(T_0) \cup D_+(T_1)$, et utiliser le fait que
$D_+(T_0)$ est le spectre d'un anneau de polynômes et que
$K$ est le corps des fractions de cet anneau de polynômes.
\end{exercise}







\section{Morphismes de surfaces vers des courbes}
\label{section-from-surfaces-to-curves}

\begin{exercise}
\label{exercise-points-projective-space}
Soit $R$ un anneau.
Soit $R \to k$ un homomorphisme de $R$ dans un corps.
Soit $n \geq 0$.
Montrer que
$$
\Mor_{\Spec(R)}(\Spec(k), \mathbf{P}^n_R)
=
(k^{n + 1} \setminus \{0\})/k^*
$$
où $k^*$ agit sur $k^{n + 1}$ par multiplication scalaire.
Désormais, nous noterons $(x_0 : \ldots : x_n)$ le
morphisme $\Spec(k) \to \mathbf{P}^n_k$ correspondant
à la classe d'équivalence de l'élément
$(x_0, \ldots, x_n) \in k^{n + 1} \setminus \{0\}$.
\end{exercise}

\begin{exercise}
\label{exercise-curve-projective-plane}
Soit $k$ un corps. Soit $Z \subset \mathbf{P}^2_k$ un
sous-schéma fermé irréductible et réduit.
Montrer que soit (a) $Z$ est un point fermé, soit (b) il existe
un polynôme homogène irréductible $F \in k[X_0, X_1, X_2]$ de degré $> 0$
tel que $Z = V_{+}(F)$, soit (c) $Z = \mathbf{P}^2_k$.
(Indication : se placer sur un ouvert affine standard.)
\end{exercise}

\begin{exercise}
\label{exercise-bezout}
Soit $k$ un corps. Soient $Z_1, Z_2 \subset \mathbf{P}^2_k$ des
sous-schémas fermés irréductibles de la forme $V_{+}(F)$
pour certains polynômes homogènes irréductibles $F_i \in k[X_0, X_1, X_2]$ de degré $> 0$.
Montrer que $Z_1 \cap Z_2$ est non vide.
(Indication : utiliser la théorie de la dimension pour estimer la dimension de
l'anneau local de $k[X_0, X_1, X_2]/(F_1, F_2)$ au point $0$.)
\end{exercise}

\begin{exercise}
\label{exercise-no-nonconstant-morphism-proj}
Montrer qu'il n'existe aucun morphisme non constant de schémas
$\mathbf{P}^2_{\mathbf{C}} \to \mathbf{P}^1_{\mathbf{C}}$
au-dessus de $\Spec(\mathbf{C})$. Ici, un {\it morphisme constant} est
un morphisme dont l'image est réduite à un seul point.
(Indication : si le morphisme n'est pas constant, considérer les fibres au-dessus de
$0$ et de $\infty$, puis montrer qu'elles doivent se rencontrer, ce qui donne une contradiction.)
\end{exercise}

\begin{exercise}
\label{exercise-nonconstant-morphism}
Soit $k$ un corps.
Supposons que $X \subset \mathbf{P}^3_k$ soit un sous-schéma fermé
défini par une seule équation homogène $F \in k[X_0, X_1, X_2, X_3]$.
En d'autres termes,
$$
X = \text{Proj}(k[X_0, X_1, X_2, X_3]/(F)) \subset \mathbf{P}^3_k
$$
comme cela a été expliqué dans le cours. Supposons que
$$
F = X_0 G + X_1 H
$$
pour des polynômes homogènes $G, H \in k[X_0, X_1, X_2, X_3]$
de degré positif. Montrer que, si $X_0, X_1, G, H$ n'ont aucun zéro commun,
alors il existe un morphisme non constant
$$
X \longrightarrow \mathbf{P}^1_k
$$
de schémas au-dessus de $\Spec(k)$
qui, sur les points à valeurs dans un corps (voir l'exercice \ref{exercise-points-projective-space}),
s'écrit $(x_0 : x_1 : x_2 : x_3) \mapsto (x_0 : x_1)$ lorsque
$x_0$ ou $x_1$ est non nul.
\end{exercise}





\section{Faisceaux inversibles}
\label{section-invertible-sheaves}

\begin{definition}
\label{definition-invertible-sheaf}
Soit $X$ un espace localement annelé.
Un {\it ${\mathcal O}_X$-module inversible} sur $X$
est un faisceau de ${\mathcal O}_X$-modules ${\mathcal L}$ tel que tout point
possède un voisinage ouvert $U \subset X$ tel que ${\mathcal L}|_U$
soit isomorphe à ${\mathcal O}_U$ comme ${\mathcal O}_U$-module.
On dit que ${\mathcal L}$ est trivial s'il est isomorphe à
${\mathcal O}_X$ comme ${\mathcal O}_X$-module.
\end{definition}

\begin{exercise}
\label{exercise-general-facts-invertible}
Généralités.
\begin{enumerate}
\item Montrer qu'un ${\mathcal O}_X$-module inversible sur
un schéma $X$ est quasi-cohérent.
\item Supposons que $X\to Y$ soit un morphisme d'espaces localement annelés,
et que ${\mathcal L}$ soit un ${\mathcal O}_Y$-module inversible.
Montrer que $f^\ast {\mathcal L}$ est un ${\mathcal O}_X$-module inversible.
\end{enumerate}
\end{exercise}

\begin{exercise}
\label{exercise-invertible-algebra}
Algèbre.
\begin{enumerate}
\item Montrer qu'un ${\mathcal O}_X$-module inversible sur
un schéma affine $\Spec(A)$ correspond à un $A$-module $M$ qui est
(i) de type fini, (ii) projectif, (iii) localement libre de rang 1,
et par conséquent (iv) plat et (v) de présentation finie. (On pourra
citer librement des résultats du cours du semestre dernier ou de livres d'algèbre.)
\item Supposons que $A$ soit un anneau intègre et que $M$ soit
un module comme en (a). Montrer que $M$ est isomorphe, comme $A$-module,
à un idéal $I \subset A$ tel que $IA_{\mathfrak p}$ soit principal pour
tout idéal premier ${\mathfrak p}$.
\end{enumerate}
\end{exercise}

\begin{definition}
\label{definition-invertible-module}
Soit $R$ un anneau. Un {\it module inversible $M$} est un $R$-module
$M$ tel que $\widetilde M$ soit un faisceau inversible sur le
spectre de $R$. On dit que $M$ est {\it trivial} si $M \cong R$ comme
$R$-module.
\end{definition}

\noindent
Autrement dit, $M$ est inversible si et seulement s'il
satisfait à toutes les conditions suivantes :
il est plat, de présentation finie, projectif et
localement libre de rang 1. (Il suffit bien entendu qu'il
soit localement libre de rang 1.)

\begin{exercise}
\label{exercise-simple-examples-invertible}
Exemples simples.
\begin{enumerate}
\item
\label{item-affine-line}
Soit $k$ un corps. Soit $A = k[x]$.
Montrer que $X = \Spec(A)$ ne possède que des
${\mathcal O}_X$-modules inversibles triviaux. Autrement dit, montrer que tout
$A$-module inversible est libre de rang 1.
\item
\label{item-affine-line-with-0-and-1-identified}
Soit $A$ l'anneau
$$
A = \{ f\in k[x] \mid f(0) = f(1) \}.
$$
Montrer qu'il existe un $A$-module inversible non trivial, sauf si
$k = {\mathbf F}_2$. (Indication : voir $\Spec(A)$ comme obtenu en identifiant
$0$ et $1$ dans ${\mathbf A}^1_k = \Spec(k[x])$.)
\item
\label{item-affine-line-with-cusp}
Même question qu'en (\ref{item-affine-line-with-0-and-1-identified})
pour l'anneau $A = k[x^2, x^3] \subset k[x]$
(cette fois, $k = {\mathbf F}_2$ convient également).
\end{enumerate}
\end{exercise}

\begin{exercise}
\label{exercise-higher-dimension-invertible}
Dimensions supérieures.
\begin{enumerate}
\item Prouver que tout faisceau inversible sur l'espace affine de dimension
deux est trivial. Plus précisément, soit
${\mathbf A}^2_k = \Spec(k[x, y])$ où $k$ est un corps.
Montrer que tout faisceau inversible sur ${\mathbf A}^2_k$ est trivial.
(Indication : une méthode consiste à considérer le module
$M$ correspondant, à examiner $M \otimes_{k[x, y]} k(x)[y]$, puis
à utiliser
l'exercice \ref{exercise-simple-examples-invertible} (\ref{item-affine-line})
pour lui trouver un générateur ; il reste alors encore à réfléchir.
Une autre méthode consiste à utiliser
l'exercice \ref{exercise-invertible-algebra}
et ce que l'on sait des idéaux de l'anneau
de polynômes : les idéaux premiers de hauteur un sont engendrés par un polynôme
irréductible ; il reste alors encore à réfléchir.)
\item Prouver que tout faisceau inversible sur tout sous-schéma
ouvert de l'espace affine de dimension deux est trivial. Plus précisément, soit
$U \subset {\mathbf A}^2_k$ un sous-schéma ouvert, où $k$ est un corps.
Montrer que tout faisceau inversible sur $U$ est trivial. Indication : montrer que tout
faisceau inversible sur $U$ se prolonge à ${\mathbf A}^2_k$. Ce n'est pas facile,
mais on peut le trouver dans \cite{H}.
\item Trouver un exemple de faisceau inversible non trivial
sur un cône épointé au-dessus d'un corps. Plus précisément,
soit $k$ un corps et soit $C = \Spec(k[x, y, z]/(xy-z^2))$.
Soit $U = C \setminus \{ (x, y, z) \}$. Trouver un faisceau inversible non trivial
sur $U$. Indication : il peut être plus facile de calculer le
groupe des classes d'isomorphisme de faisceaux inversibles sur $U$ que d'en
trouver simplement un. Noter que $U$ est recouvert par les ouverts
$\Spec(k[x, y, z, 1/x]/(xy-z^2))$ et
$\Spec(k[x, y, z, 1/y]/(xy-z^2))$,
qui sont « faciles » à traiter.
\end{enumerate}
\end{exercise}

\begin{definition}
\label{definition-picard-group}
Soit $X$ un espace localement annelé.
Le {\it groupe de Picard de $X$} est l'ensemble $\Pic(X)$
des classes d'isomorphisme de $\mathcal{O}_X$-modules inversibles,
la loi d'addition étant donnée par le produit tensoriel.
Voir Modules, définition \ref{modules-definition-pic}.
Pour un anneau $R$, on pose $\Pic(R) = \Pic(\Spec(R))$.
\end{definition}

\begin{exercise}
\label{exercise-traverso}
Soit $R$ un anneau.
\begin{enumerate}
\item Montrer que si $R$ est un anneau intègre normal noethérien, alors
$\Pic(R) = \Pic(R[t])$. [Indication : il existe un homomorphisme
$R[t] \to R$, $t \mapsto 0$ qui est un inverse à gauche de l'homomorphisme
$R \to R[t]$. Il suffit donc de montrer que tout
$R[t]$-module inversible $M$ tel que $M/tM \cong R$ est libre de rang $1$.
Soit $K$ le corps des fractions de $R$.
Choisir une trivialisation $K[t] \to M \otimes_{R[t]} K[t]$, ce qui est possible par
l'exercice \ref{exercise-simple-examples-invertible} (\ref{item-affine-line}).
L'ajuster pour qu'elle induise la trivialisation
ci-dessus de $M/tM$. Montrer qu'elle est en fait une trivialisation de
$M$ sur $R[t]$ (c'est ici qu'intervient la normalité).]
\item Soit $k$ un corps. Montrer que
$\Pic(k[x^2, x^3, t]) \not = \Pic(k[x^2, x^3])$.
\end{enumerate}
\end{exercise}













\section{Cohomologie de {\v C}ech}
\label{section-cech-cohomology}

\begin{exercise}
\label{exercise-cech-cohomology}
Cohomologie de {\v C}ech. Ici, $k$ est un corps.
\begin{enumerate}
\item Soit $X$ un schéma muni d'un recouvrement ouvert
${\mathcal U} : X = U_1 \cup U_2$, avec $U_1 = \Spec(k[x])$,
$U_2 =  \Spec(k[y])$,
$U_1 \cap U_2 = \Spec(k[z, 1/z])$ et les immersions ouvertes
$U_1 \cap U_2 \to U_1$ resp.\ $U_1 \cap U_2 \to U_2$ étant déterminées
par $x \mapsto z$ resp.\ $y \mapsto z$ (et c'est bien ce que je veux dire).
(Nous avons vu en cours qu'un tel $X$ existe ; c'est la droite affine
à origine dédoublée.) Calculer ${\check H}^1({\mathcal U}, {\mathcal
O})$ ;
par exemple,\ en donner une base comme espace vectoriel sur $k$.
\item Pour chaque élément de
${\check H}^1({\mathcal U}, {\mathcal O})$,
construire une suite exacte de faisceaux de ${\mathcal O}_X$-modules
$$
0 \to {\mathcal O}_X \to E \to {\mathcal O}_X \to 0
$$
telle que le cobord $\delta(1) \in {\check H}^1({\mathcal U},
{\mathcal O})$
soit l'élément donné. (Il faut notamment donner un sens à cet énoncé.
Voir aussi ci-dessous.
On peut également montrer abstraitement qu'un tel objet doit exister.)
\end{enumerate}
\end{exercise}

\begin{definition}
\label{definition-delta}
(Définition de delta.) Supposons que
$$
0 \to {\mathcal F}_1 \to {\mathcal F}_2 \to {\mathcal F}_3 \to 0
$$
soit une suite exacte courte de faisceaux abéliens sur un espace topologique $X$.
L'application cobord
$\delta : H^0(X, {\mathcal F}_3) \to {\check H}^1(X, {\mathcal F}_1)$
est définie comme suit. Soit un élément $\tau \in H^0(X, {\mathcal F}_3)$.
Choisir un recouvrement ouvert ${\mathcal U} : X = \bigcup_{i\in I} U_i$ tel
que, pour tout $i$, il existe une section $\tilde \tau_i \in {\mathcal F}_2$
relevant la restriction de $\tau$ à $U_i$. Considérer alors la correspondance
$$
(i_0, i_1) \longmapsto
\tilde \tau_{i_0}|_{U_{i_0i_1}} - \tilde \tau_{i_1}|_{U_{i_0i_1}}.
$$
C'est clairement un 1-cobord dans le complexe de {\v C}ech
${\check C}^\ast({\mathcal U}, {\mathcal F}_2)$. Mais on observe que
(en considérant ${\mathcal F}_1$ comme un sous-faisceau de ${\mathcal F}_2$) le membre de droite
est toujours une section de ${\mathcal F}_1$ sur $U_{i_0i_1}$. On voit donc
que cette correspondance définit une 1-cochaîne dans le complexe
${\check C}^\ast({\mathcal U}, {\mathcal F}_2)$. La classe de cohomologie
de cette 1-cochaîne est par définition {\it $\delta(\tau)$}.
\end{definition}



\section{Cohomologie}
\label{section-cohomology}

\begin{exercise}
\label{exercise-cohomology-not-zero}
Soit $X = \mathbf{R}$ muni de la topologie euclidienne usuelle.
En utilisant seulement les propriétés formelles de la cohomologie (fonctorialité
et suite exacte longue de cohomologie), montrer qu'il
existe un faisceau $\mathcal{F}$ sur $X$ tel que $H^1(X, \mathcal{F})$ soit non nul.
\end{exercise}

\begin{exercise}
\label{exercise-mayer-vietoris}
Soit $X = U \cup V$ un espace topologique écrit comme la
réunion de deux ouverts. On a alors une suite exacte longue
de Mayer--Vietoris
$$
0 \to
H^0(X, \mathcal{F}) \to
H^0(U, \mathcal{F}) \oplus H^0(V, \mathcal{F}) \to
H^0(U \cap V, \mathcal{F}) \to
H^1(X, \mathcal{F}) \to \ldots
$$
Quelle propriété des faisceaux injectifs est essentielle à la construction
de la suite exacte longue de Mayer--Vietoris ? Pourquoi est-elle satisfaite ?
\end{exercise}

\begin{exercise}
\label{exercise-cohomology-two-point-space}
Soit $X$ un espace topologique.
\begin{enumerate}
\item Montrer que $H^i(X, \mathcal{F})$ est nul pour $i > 0$
si $X$ possède au plus $2$ points.
\item Que se passe-t-il si $X$ possède $3$ points ?
\end{enumerate}
\end{exercise}

\begin{exercise}
\label{exercise-cohomology-spec-local-ring}
Soit $X$ le spectre d'un anneau local. Montrer que
$H^i(X, \mathcal{F})$ est nul pour $i > 0$ et pour tout
faisceau de groupes abéliens $\mathcal{F}$.
\end{exercise}

\begin{exercise}
\label{exercise-affine-morphism}
Soit $f : X \to Y$ un morphisme affine de schémas.
Prouver que $H^i(X, \mathcal{F}) = H^i(Y, f_*\mathcal{F})$
pour tout $\mathcal{O}_X$-module quasi-cohérent $\mathcal{F}$.
On pourra imposer des conditions supplémentaires à $X$ et à $Y$
et utiliser la coïncidence de la cohomologie de {\v C}ech et de la cohomologie
pour les faisceaux quasi-cohérents et les recouvrements ouverts affines des schémas séparés.
\end{exercise}

\begin{exercise}
\label{exercise-affine-morphism-to-Pn}
Soit $A$ un anneau. Soit $\mathbf{P}^n_A = \text{Proj}(A[T_0, \ldots, T_n])$
l'espace projectif sur $A$. Soit
$\mathbf{A}^{n + 1}_A = \Spec(A[T_0, \ldots, T_n])$ et soit
$$
U = \bigcup\nolimits_{i = 0, \ldots, n} D(T_i) \subset \mathbf{A}^{n + 1}_A
$$
le complémentaire de l'image de l'immersion fermée
$0 : \Spec(A) \to \mathbf{A}^{n + 1}_A$.
Construire un morphisme affine surjectif
$$
f : U \longrightarrow \mathbf{P}^n_A
$$
et prouver que $f_*\mathcal{O}_U =
\bigoplus_{d \in \mathbf{Z}} \mathcal{O}_{\mathbf{P}^n_A}(d)$.
Plus généralement, montrer que, pour un $A[T_0, \ldots, T_n]$-module gradué $M$,
on a
$$
f_*(\widetilde{M}|_U) =
\bigoplus\nolimits_{d \in \mathbf{Z}} \widetilde{M(d)}
$$
où, au membre de gauche, figure le faisceau quasi-cohérent
$\widetilde{M}$ associé à $M$ sur $\mathbf{A}^{n + 1}_A$
et, au membre de droite, les faisceaux quasi-cohérents
$\widetilde{M(d)}$ associés au module gradué $M(d)$.
\end{exercise}

\begin{exercise}
\label{exercise-compute-Pn}
Soit $A$ un anneau et soit $\mathbf{P}^n_A = \text{Proj}(A[T_0, \ldots, T_n])$
l'espace projectif sur $A$.
Calculer soigneusement la cohomologie des tordus de Serre
$\mathcal{O}_{\mathbf{P}^n_A}(d)$ du faisceau
structural sur $\mathbf{P}^n_A$. On pourra utiliser la cohomologie de {\v C}ech
et la coïncidence de la cohomologie de {\v C}ech et de la cohomologie
pour les faisceaux quasi-cohérents et les recouvrements ouverts affines des schémas séparés.
\end{exercise}

\begin{exercise}
\label{exercise-cohomology-hypersurface}
Soit $A$ un anneau et soit $\mathbf{P}^n_A = \text{Proj}(A[T_0, \ldots, T_n])$
l'espace projectif sur $A$. Soit
$F \in A[T_0, \ldots, T_n]$ homogène de degré $d$.
Soit $X \subset \mathbf{P}^n_A$ le sous-schéma fermé
correspondant à l'idéal gradué $(F)$ de $A[T_0, \ldots, T_n]$.
Que peut-on dire de $H^i(X, \mathcal{O}_X)$ ?
\end{exercise}

\begin{exercise}
\label{exercise-characterize-finite-product-fields}
Soit $R$ un anneau tel que, pour tout foncteur exact à gauche
$F : \text{Mod}_R \to \textit{Ab}$, on ait
$R^iF = 0$ pour $i > 0$. Montrer que $R$ est un produit fini
de corps.
\end{exercise}






\section{Compléments de cohomologie}
\label{section-more-cohomology}

\begin{exercise}
\label{exercise-cohomology-coordinate-axes}
Soit $k$ un corps. Soit $X \subset \mathbf{P}^n_k$
la « croix des axes de coordonnées ». Autrement dit, soit $X$
défini par les équations homogènes
$$
T_i T_j = 0\text{ pour }i > j > 0
$$
où, comme d'habitude, on écrit $\mathbf{P}^n_k = \text{Proj}(k[T_0, \ldots, T_n])$.
Autrement dit, $X$ est le sous-schéma fermé correspondant au quotient
$k[T_0, \ldots, T_n]/(T_iT_j; i > j > 0)$ de l'anneau de polynômes.
Calculer $H^i(X, \mathcal{O}_X)$ pour tout $i$. Indication : utiliser la cohomologie de {\v C}ech.
\end{exercise}

\begin{exercise}
\label{exercise-compute-cohomology-punctured}
Soit $A$ un anneau. Soit $I = (f_1, \ldots, f_t)$
un idéal de type fini de $A$.
Soit $U \subset \Spec(A)$ le complémentaire de $V(I)$.
Pour tout $A$-module $M$, écrire un complexe
de $A$-modules (en fonction de $A$, $f_1, \ldots, f_t$, $M$)
dont les groupes de cohomologie donnent $H^n(U, \widetilde{M})$.
\end{exercise}

\begin{exercise}
\label{exercise-compute-cohomology-affine-space-punctured}
Soit $k$ un corps. Soit $U \subset \mathbf{A}^d_k$ le
complémentaire du point fermé $0$ de $\mathbf{A}^d_k$.
Calculer $H^n(U, \mathcal{O}_U)$ pour tout $n$.
\end{exercise}

\begin{exercise}
\label{exercise-find-curve-genus-one-hundred}
Soit $k$ un corps. Trouver explicitement un schéma $X$ projectif sur $k$,
de dimension $1$, tel que $H^0(X, \mathcal{O}_X) = k$ et
$\dim_k H^1(X, \mathcal{O}_X) = 100$.
\end{exercise}

\begin{exercise}
\label{exercise-degree-2-cover}
Soit $f : X \to Y$ un morphisme fini localement libre de degré $2$.
Supposons que $X$ et $Y$ soient des schémas intègres et que $2$ soit inversible
dans le faisceau structural de $Y$, c'est-à-dire que $2 \in \Gamma(Y, \mathcal{O}_Y)$
soit inversible. Montrer que l'homomorphisme de $\mathcal{O}_Y$-modules
$$
f^\sharp : \mathcal{O}_Y \longrightarrow f_*\mathcal{O}_X
$$
possède un inverse à gauche, c'est-à-dire qu'il existe un homomorphisme de $\mathcal{O}_Y$-modules
$\tau : f_*\mathcal{O}_X \to \mathcal{O}_Y$
tel que $\tau \circ f^\sharp = \text{id}$.
En déduire que $H^n(Y, \mathcal{O}_Y) \to H^n(X, \mathcal{O}_X)$
est injectif\footnote{Il existe effectivement un morphisme fini localement libre
$X \to Y$ entre schémas intègres, de degré $2$, pour lequel l'homomorphisme
$H^1(Y, \mathcal{O}_Y) \to H^1(X, \mathcal{O}_X)$ n'est pas injectif.}.
\end{exercise}

\begin{exercise}
\label{exercise-pic}
Soit $X$ un schéma (ou un espace localement annelé). La règle
$U \mapsto \mathcal{O}_X(U)^*$ définit un faisceau de groupes
noté $\mathcal{O}_X^*$.
Expliquer brièvement pourquoi le groupe de Picard de $X$
(définition \ref{definition-picard-group}) est
égal à $H^1(X, \mathcal{O}_X^*)$.
\end{exercise}

\begin{exercise}
\label{exercise-pic-nontrivial}
Donner un exemple de schéma affine $X$ tel que $\Pic(X)$ soit non trivial.
En déduire, à l'aide de l'exercice \ref{exercise-pic}, que $H^1(X, -)$ n'est pas le
foncteur nul pour un tel $X$.
\end{exercise}

\begin{exercise}
\label{exercise-kill-cohomology-complement}
Soit $A$ un anneau. Soit $I = (f_1, \ldots, f_t)$ un idéal de type fini
de $A$. Soit $U \subset \Spec(A)$ le complémentaire de $V(I)$.
Étant donnés un $\mathcal{O}_{\Spec(A)}$-module quasi-cohérent $\mathcal{F}$
et $\xi \in H^p(U, \mathcal{F})$ avec $p > 0$, montrer qu'il existe
$n > 0$ tel que $f_i^n \xi = 0$ pour $i = 1, \ldots, t$.
Indication : une méthode possible consiste à utiliser le complexe
trouvé dans l'exercice \ref{exercise-compute-cohomology-punctured}.
\end{exercise}

\begin{exercise}
\label{exercise-h0-complement}
Soit $A$ un anneau. Soit $I = (f_1, \ldots, f_t)$ un
idéal de type fini de $A$. Soit $U \subset \Spec(A)$
le complémentaire de $V(I)$. Soit $M$ un $A$-module
dont la $I$-torsion est nulle, c'est-à-dire
$0 = \Ker((f_1, \ldots, f_t) : M \to M^{\oplus t})$.
Montrer qu'il existe un isomorphisme canonique
$$
H^0(U, \widetilde{M}) = \colim \Hom_A(I^n, M).
$$
Attention : ce n'est pas trivial.
\end{exercise}

\begin{exercise}
\label{exercise-Noetherian}
Soit $A$ un anneau noethérien. Soit $I$ un idéal de $A$.
Soit $M$ un $A$-module. Soit $M[I^\infty]$ l'ensemble des éléments de
torsion annulés par une puissance de $I$, défini par
$$
M[I^\infty] = \{x \in M \mid
\text{ il existe un }n \geq 1\text{ tel que }I^nx = 0\}
$$
Posons $M' = M/M[I^\infty]$. La torsion annulée par une puissance de $I$ dans $M'$ est alors nulle.
Montrer que
$$
\colim \Hom_A(I^n, M) = \colim \Hom_A(I^n, M').
$$
Attention : ce n'est pas trivial. Indications : (1) essayer de se ramener au cas où $M$ est de type fini,
(2) montrer que tout élément de $\Ext^1_A(I^n, N)$
s'envoie sur zéro dans $\Ext^1_A(I^{n + m}, N)$ pour un certain $m > 0$
si $N = M[I^\infty]$ et si $M$ est de type fini, (3) montrer la même assertion
comme en (2) pour $\Hom_A(I^n, N)$, (3) considérer la suite exacte longue
de cohomologie
$$
0 \to \Hom_A(I^n, M[I^\infty]) \to \Hom_A(I^n, M) \to \Hom_A(I^n, M')
\to \Ext^1_A(I^n, M[I^\infty])
$$
pour $M$ de type fini et la comparer à la suite relative à $I^{n + m}$ pour conclure.
\end{exercise}








\section{Retour sur la cohomologie}
\label{section-more-more-cohomology}


\begin{exercise}
\label{exercise-nonkillable}
Construire un exemple d'un corps $k$, d'une courbe $X$ sur $k$,
d'un $\mathcal{O}_X$-module inversible $\mathcal{L}$ et
d'une classe de cohomologie $\xi \in H^1(X, \mathcal{L})$ vérifiant
la propriété suivante : pour tout morphisme fini surjectif
$\pi : Y \to X$ de schémas, l'élément $\xi$
a pour image réciproque un élément non nul de $H^1(Y, \pi^*\mathcal{L})$.
Indication : construire $X$, $k$, $\mathcal{L}$ de telle sorte que
l'on ait une suite exacte courte
$0 \to \mathcal{L} \to \mathcal{O}_X \to i_*\mathcal{O}_Z \to 0$
où $Z \subset X$ soit un sous-schéma fermé constitué de
plus de $1$ point fermé. Examiner ensuite ce qui se passe par
image réciproque.
\end{exercise}

\begin{exercise}
\label{exercise-cohomology-cycle-curves}
Soit $k$ un corps algébriquement clos.
Soit $X$ un schéma projectif de dimension $1$.
Supposons que $X$ contienne un cycle de courbes, c'est-à-dire qu'il
existe un $n \geq 2$, des sous-schémas fermés intègres
de dimension $1$ deux à deux distincts
$C_1, \ldots, C_n$, ainsi que des points fermés deux à deux distincts
$x_1, \ldots, x_n \in X$, tels que $x_n \in C_n \cap C_1$
et $x_i \in C_i \cap C_{i + 1}$
pour $i = 1, \ldots, n - 1$.
Démontrer que $H^1(X, \mathcal{O}_X)$ est non nul.
Indication : soit $\mathcal{F}$ l'image du morphisme
$\mathcal{O}_X \to \bigoplus \mathcal{O}_{C_i}$,
et montrer que $H^1(X, \mathcal{F})$ est non nul
en utilisant $\kappa(x_i) = k$ et $H^0(C_i, \mathcal{O}_{C_i}) = k$.
Utiliser aussi $H^2(X, -) = 0$, d'après le théorème de Grothendieck.
\end{exercise}

\begin{exercise}
\label{exercise-surface}
Soit $X$ une surface projective sur un corps algébriquement clos $k$.
Démontrer qu'il existe un sous-schéma fermé strict $Z \subset X$ tel que
$H^1(Z, \mathcal{O}_Z)$ soit non nul. Indication : utiliser
l'exercice \ref{exercise-cohomology-cycle-curves}.
\end{exercise}

\begin{exercise}
\label{exercise-surface-better}
Soit $X$ une surface projective sur un corps algébriquement clos $k$.
Montrer que, pour tout $n \geq 0$,
il existe un sous-schéma fermé strict $Z \subset X$ tel que
$\dim_k H^1(Z, \mathcal{O}_Z) > n$.
Indiquer seulement comment procéder en modifiant les arguments
de l'exercice \ref{exercise-surface} et de
\ref{exercise-cohomology-cycle-curves} ; ne pas donner tous les détails.
\end{exercise}

\begin{exercise}
\label{exercise-surface-h2-nonzero}
Soit $X$ une surface projective sur un corps algébriquement clos $k$.
Démontrer qu'il existe un $\mathcal{O}_X$-module cohérent $\mathcal{F}$
tel que $H^2(X, \mathcal{F})$ soit non nul. Indication : utiliser le résultat
de l'exercice \ref{exercise-surface-better} ainsi qu'une
suite exacte judicieusement choisie.
\end{exercise}

\begin{exercise}
\label{exercise-kunneth}
Soient $X$ et $Y$ des schémas sur un corps $k$ (on pourra supposer
$X$ et $Y$ suffisamment raisonnables, par exemple qcqs ou projectifs sur $k$).
Posons $X \times Y = X \times_{\Spec(k)} Y$ et notons
$p : X \times Y \to X$ et $q : X \times Y \to Y$ les projections. Pour
un $\mathcal{O}_X$-module quasi-cohérent $\mathcal{F}$ et
un $\mathcal{O}_Y$-module quasi-cohérent $\mathcal{G}$, démontrer que
$$
H^n(X \times Y,
p^*\mathcal{F} \otimes_{\mathcal{O}_{X \times Y}} q^*\mathcal{G}) =
\bigoplus\nolimits_{a + b = n}
H^a(X, \mathcal{F}) \otimes_k H^b(Y, \mathcal{G})
$$
ou montrer seulement que l'égalité vaut après passage aux dimensions.
Des points supplémentaires seront accordés aux solutions « élégantes ».
\end{exercise}

\begin{exercise}
\label{exercise-Oab}
Soit $k$ un corps. Soit $X = \mathbf{P}|^1 \times \mathbf{P}^1$
le produit de la droite projective sur $k$ par elle-même,
muni des projections $p : X \to \mathbf{P}^1_k$ et $q : X \to \mathbf{P}^1_k$.
Posons
$$
\mathcal{O}(a, b) = p^*\mathcal{O}_{\mathbf{P}^1_k}(a)
\otimes_{\mathcal{O}_X} q^*\mathcal{O}_{\mathbf{P}^1_k}(b)
$$
Calculer les dimensions de $H^i(X, \mathcal{O}(a, b))$
pour tous $i, a, b$. Indication : utiliser l'exercice \ref{exercise-kunneth}.
\end{exercise}










\section{Cohomologie et polynômes de Hilbert}
\label{section-cohomology-hilbert-polynomials}

\begin{situation}
\label{situation-hilbert-polynomial}
Soit $k$ un corps. Soit $X = \mathbf{P}^n_k$
l'espace projectif de dimension $n$. Soit $\mathcal{F}$
un $\mathcal{O}_X$-module cohérent. Rappelons que
$$
\chi(X, \mathcal{F}) =
\sum\nolimits_{i = 0}^n (-1)^i \dim_k H^i(X, \mathcal{F})
$$
Rappelons que le {\it polynôme de Hilbert} de $\mathcal{F}$ est la fonction
$$
t \longmapsto \chi(X, \mathcal{F}(t))
$$
Rappelons aussi que
$\mathcal{F}(t) = \mathcal{F} \otimes_{\mathcal{O}_X} \mathcal{O}_X(t)$
où $\mathcal{O}_X(t)$ est le $t$-ième tordu du faisceau structural
comme dans Constructions, définition \ref{constructions-definition-twist}.
Dans Variétés, sous-section \ref{varieties-subsection-hilbert}, nous avons
démontré que le polynôme de Hilbert est un polynôme en $t$.
\end{situation}

\begin{exercise}
\label{exercise-hilbert-pol-easy}
Dans la situation \ref{situation-hilbert-polynomial}.
\begin{enumerate}
\item Si $P(t)$ est le polynôme de Hilbert de $\mathcal{F}$,
quel est le polynôme de Hilbert de $\mathcal{F}(-13)$ ?
\item Si $P_i$ est le polynôme de Hilbert de $\mathcal{F}_i$,
quel est le polynôme de Hilbert de $\mathcal{F}_1 \oplus \mathcal{F}_2$ ?
\item Si $P_i$ est le polynôme de Hilbert de $\mathcal{F}_i$ et si
$\mathcal{F}$ est le noyau d'un morphisme surjectif
$\mathcal{F}_1 \to \mathcal{F}_2$, quel est le polynôme de Hilbert de
$\mathcal{F}$ ?
\end{enumerate}
\end{exercise}

\begin{exercise}
\label{exercise-find-given-hilbert-pol-dim-1}
Dans la situation \ref{situation-hilbert-polynomial}, supposons $n \geq 1$.
Trouver un faisceau cohérent dont le polynôme de Hilbert est $t - 101$.
\end{exercise}

\begin{exercise}
\label{exercise-find-given-hilbert-pol-dim-2}
Dans la situation \ref{situation-hilbert-polynomial}, supposons $n \geq 2$.
Trouver un faisceau cohérent dont le polynôme de Hilbert est
$t^2/2 + t/2 - 1$. (Cet exercice est assez délicat ; il suffit
de montrer qu'un tel faisceau cohérent existe.)
\end{exercise}

\begin{exercise}
\label{exercise-bound-genus-in-degree}
Dans la situation \ref{situation-hilbert-polynomial}, supposons $n \geq 2$
et que $k$ soit algébriquement clos.
Soit $C \subset X$ un sous-schéma fermé intègre de dimension $1$.
Autrement dit, $C$ est une courbe projective.
Soit $d t + e$ le polynôme de Hilbert de $\mathcal{O}_C$
considéré comme faisceau cohérent sur $X$.
\begin{enumerate}
\item Donner une borne supérieure de $e$. (Indications : utiliser le fait que
$\mathcal{O}_C(t)$ n'a de cohomologie qu'aux degrés $0$ et $1$
et étudier $H^0(C, \mathcal{O}_C)$.)
\item Choisir une section globale $s$ de $\mathcal{O}_X(1)$
qui coupe $C$ transversalement, c'est-à-dire telle qu'il
existe des points fermés deux à deux distincts $c_1, \ldots, c_r \in C$ et
une suite exacte courte
$$
0 \to \mathcal{O}_C \xrightarrow{s} \mathcal{O}_C(1) \to
\bigoplus\nolimits_{i = 1, \ldots, r} k_{c_i} \to 0
$$
où $k_{c_i}$ est le faisceau gratte-ciel de valeur $k$ en $c_i$.
(Une telle section $s$ existe ; on l'admettra.)
Montrer que $r = d$. (Indication : tordre la suite et examiner ce que l'on obtient.)
\item En tordant la suite exacte courte, on obtient une application $k$-linéaire
$\varphi_t : \Gamma(C, \mathcal{O}_C(t)) \to \bigoplus_{i = 1, \ldots, d} k$
pour tout $t$. Montrer que cette application est surjective si $t \geq d - 1$.
\item Donner une borne inférieure de $e$ en fonction de $d$. (Indication : montrer que
$H^1(C, \mathcal{O}_C(d - 2)) = 0$ à l'aide du résultat de (3), puis utiliser
un résultat d'annulation.)
\end{enumerate}
\end{exercise}

\begin{exercise}
\label{exercise-three-quadrics-in-plane}
Dans la situation \ref{situation-hilbert-polynomial}, supposons $n = 2$.
Soient $s_1, s_2, s_3 \in \Gamma(X, \mathcal{O}_X(2))$ trois
équations quadratiques. Considérons le faisceau cohérent
$$
\mathcal{F} = \Coker\left(\mathcal{O}_X(-2)^{\oplus 3}
\xrightarrow{s_1, s_2, s_3} \mathcal{O}_X\right)
$$
Énumérer les polynômes de Hilbert possibles d'un tel $\mathcal{F}$.
(Essayer de visualiser les intersections de quadriques dans le plan projectif.)
\end{exercise}





\section{Courbes}
\label{section-curves}

\begin{exercise}
\label{exercise-closed-subset-vanshing}
Soit $k$ un corps algébriquement clos. Soit $X$ une courbe projective
sur $k$. Soit $\mathcal{L}$ un $\mathcal{O}_X$-module inversible.
Soient $s_0, \ldots, s_n \in H^0(X, \mathcal{L})$
des sections globales de $\mathcal{L}$.
Démontrer qu'il existe un sous-schéma fermé naturel
$$
Z \subset
\mathbf{P}^n \times X
$$
tel que le point fermé
$((\lambda_0 : \ldots : \lambda_n), x)$
appartienne à $Z$ si et seulement si la section
$\lambda_0 s_0 + \ldots + \lambda_n s_n$
s'annule en $x$. (Indication : décrire $Z$ localement sur des ouverts affines.)
\end{exercise}

\begin{exercise}
\label{exercise-diagonal-cartier}
Soit $k$ un corps algébriquement clos. Soit $X$ une courbe lisse sur $k$.
Soit $r \geq 1$. Montrer que le sous-ensemble fermé
$$
D \subset X \times X^r = X^{r + 1}
$$
dont les points fermés sont les uplets $(x, x_1, \ldots, x_r)$ tels que
$x = x_i$ pour un certain $i$, possède un faisceau d'idéaux inversible. Autrement dit,
montrer que
$D$ est un diviseur de Cartier effectif. Indications : se ramener au cas $r = 1$ et
utiliser le fait que $X$ est une courbe lisse pour obtenir une propriété de la diagonale
(consulter les ouvrages sur ce point).
\end{exercise}

\begin{exercise}
\label{exercise-one-divisor-in-another}
Soit $k$ un corps algébriquement clos. Soit $X$ une courbe projective lisse
sur $k$. Soit $T$ un schéma de type fini sur $k$ et soient
$$
D_1 \subset X \times T
\quad\text{et}\quad
D_2 \subset X \times T
$$
deux diviseurs de Cartier effectifs tels que, pour $t \in T$, les
fibres $D_{i, t} \subset X_t$ ne soient pas denses (c'est-à-dire ne contiennent pas
le point générique de la courbe $X_t$).
Démontrer qu'il existe un sous-schéma fermé canonique
$Z \subset T$ tel qu'un point fermé $t \in T$ appartienne à $Z$
si et seulement si, pour les fibres schématiques $D_{1, t}$, $D_{2, t}$
de $D_1$, $D_2$, on a
$$
D_{1, t} \subset D_{2, t}
$$
comme sous-schémas fermés de $X_t$.
Indications : montrer que, quitte à restreindre $T$, on peut
supposer $T = \Spec(A)$ affine et qu'il existe un ouvert affine $U \subset X$
tel que $D_i \subset U \times T$. Puis montrer que
$M_1 = \Gamma(D_1, \mathcal{O}_{D_1})$ est un $A$-module localement libre de rang fini
(il faudra ici un peu d'algèbre non triviale --- demander à vos amis).
Après avoir restreint $T$, on peut supposer que $M_1$ est un $A$-module libre.
Considérer alors
$$
\Gamma(U \times T, \mathcal{I}_{D_2}) \to M_1 = A^{\oplus N}
$$
et définir $Z$ comme le sous-schéma fermé découpé par l'idéal engendré
par les coordonnées des vecteurs de l'image de ce morphisme. Expliquer pourquoi cela fonctionne
(cela demandera peut-être encore un peu d'algèbre commutative).
\end{exercise}

\begin{exercise}
\label{exercise-closed-subset-vanshing-r}
Soit $k$ un corps algébriquement clos. Soit $X$ une courbe projective lisse
sur $k$. Soit $\mathcal{L}$ un $\mathcal{O}_X$-module inversible.
Soient $s_0, \ldots, s_n \in H^0(X, \mathcal{L})$
des sections globales de $\mathcal{L}$. Soit $r \geq 1$.
Démontrer qu'il existe un sous-schéma fermé naturel
$$
Z \subset
\mathbf{P}^n \times X \times \ldots \times X = \mathbf{P}^n \times X^r
$$
tel que le point fermé
$((\lambda_0 : \ldots : \lambda_n), x_1, \ldots, x_r)$
appartienne à $Z$ si et seulement si la section
$s_\lambda = \lambda_0 s_0 + \ldots + \lambda_n s_n$
s'annule sur le diviseur $D = x_1 + \ldots + x_r$, c'est-à-dire que
la section $s_\lambda$ appartient à $\mathcal{L}(-D)$.
Indication : expliquer comment cela résulte de la combinaison des résultats des
exercices \ref{exercise-diagonal-cartier} et
\ref{exercise-one-divisor-in-another}.
\end{exercise}

\begin{exercise}
\label{exercise-effective}
Soit $k$ un corps algébriquement clos. Soit $X$ une courbe projective lisse
sur $k$. Soit $\mathcal{L}$ un $\mathcal{O}_X$-module inversible.
Montrer qu'il existe un sous-ensemble fermé naturel
$$
Z \subset X^r
$$
tel qu'un point fermé $(x_1, \ldots, x_r)$ de $X^r$ appartienne à $Z$
si et seulement si $\mathcal{L}(-x_1 - \ldots -x_r)$ admet une section
globale non nulle. Indication : utiliser l'exercice \ref{exercise-closed-subset-vanshing-r}.
\end{exercise}

\begin{exercise}
\label{exercise-difference-effective}
Soit $k$ un corps algébriquement clos. Soit $X$ une courbe projective lisse
sur $k$. Soient $r \geq s$ des entiers. Montrer qu'il existe un sous-ensemble fermé
naturel
$$
Z \subset X^r \times X^s
$$
tel qu'un point fermé $(x_1, \ldots, x_r, y_1, \ldots, y_s)$ de
$X^r \times X^s$ appartienne à $Z$
si et seulement si $x_1 + \ldots + x_r - y_1 - \ldots - y_s$
est linéairement équivalent à un diviseur effectif.
Indication : choisir un module inversible auxiliaire $\mathcal{L}$
de degré suffisamment grand pour que $\mathcal{L}(-D)$ admette une section
non nulle pour tout diviseur effectif $D$ de degré $r$.
Utiliser ensuite deux fois le résultat de l'exercice \ref{exercise-effective}.
\end{exercise}

\begin{exercise}
\label{exercise-genus-7}
Choisir un corps algébriquement clos $k$ de son choix.
Déterminer aussi complètement que possible tous les $\mathfrak g^r_d$
susceptibles d'exister sur une courbe de genre $7$.
Au cours de cette étude, essayer aussi de
\begin{enumerate}
\item déterminer dans quels cas le $\mathfrak g^r_d$ est sans point de base, et
\item déterminer dans quels cas le $\mathfrak g^r_d$ définit une immersion fermée
dans $\mathbf{P}^r$.
\end{enumerate}
Faire de même si l'on suppose la courbe « générale » (adopter
la notion de généralité de son choix -- cela peut être plus facile que la question précédente).
Faire de même si l'on suppose la courbe hyperelliptique.
Faire de même si l'on suppose la courbe trigonale
(et non hyperelliptique). Etc.
\end{exercise}





\section{Modules}
\label{section-moduli}

\noindent
Dans cette section, nous examinons quelques approches naïves des problèmes de modules
d'objets de géométrie algébrique.

\medskip\noindent
Soit $k$ un corps algébriquement clos. Supposons que $M$ soit un
problème de modules sur $k$. Nous ne définirons pas ici avec précision
ce que cela signifie, mais le sens devrait être clair dans chaque exercice.
Pour comprendre ce qui suit, il suffit de connaître
les objets de $M$ sur $k$, les isomorphismes entre
ces objets de $M$ sur $k$, et les familles d'objets de $M$
sur une variété. On dit alors que le {\it nombre de modules de $M$}
est $d \geq 0$ si les conditions suivantes sont vérifiées :
\begin{enumerate}
\item il existe un nombre fini de familles $X_i \to V_i$, $i = 1, \ldots, n$
telles que tout objet de $M$ sur $k$ soit isomorphe à une fibre de l'une
d'entre elles et que $\max \dim(V_i) = d$, et
\item il est impossible de le faire avec une valeur plus petite de $d$.
\end{enumerate}
Il ne s'agit en réalité que d'une approximation très grossière de notions meilleures
que l'on trouve dans la littérature.

\begin{exercise}
\label{exercise-number-moduli-lines-conics}
Soit $k$ un corps algébriquement clos. Soient $d \geq 1$ et $n \geq 1$.
Convenons que le problème de modules des hypersurfaces de degré $d$ dans $P^n$ est donné par
\begin{enumerate}
\item un objet est une hypersurface $X \subset \mathbf{P}^n_k$
de degré $d$,
\item un isomorphisme entre deux objets $X \subset \mathbf{P}^n_k$
et $Y \subset \mathbf{P}^n_k$ est un élément $g \in \text{PGL}_n(k)$
tel que $g(X) = Y$, et
\item une famille d'hypersurfaces sur une variété $V$ est un sous-schéma
fermé $X \subset \mathbf{P}^n_V$ tel que, pour tout $v \in V$,
la fibre schématique $X_v$ de $X \to V$ soit
une hypersurface dans $\mathbf{P}^n_v$.
\end{enumerate}
Calculer (si possible -- les questions deviennent de plus en plus difficiles)
\begin{enumerate}
\item le nombre de modules lorsque $n = 1$ et $d$ est arbitraire,
\item le nombre de modules lorsque $n = 2$ et $d = 1$,
\item le nombre de modules lorsque $n = 2$ et $d = 2$,
\item le nombre de modules lorsque $n \geq 1$ et $d = 2$,
\item le nombre de modules lorsque $n = 2$ et $d = 3$,
\item le nombre de modules lorsque $n = 3$ et $d = 3$, et
\item le nombre de modules lorsque $n = 2$ et $d = 4$.
\end{enumerate}
\end{exercise}

\begin{exercise}
\label{exercise-number-moduli-hyperelliptic}
Soit $k$ un corps algébriquement clos. Soit $g \geq 2$.
Convenons que le problème de modules des courbes hyperelliptiques de genre $g$ est donné par
\begin{enumerate}
\item un objet est une courbe hyperelliptique projective lisse $X$
de genre $g$,
\item un isomorphisme entre deux objets $X$ et $Y$ est un
isomorphisme $X \to Y$ de schémas sur $k$, et
\item une famille de courbes hyperelliptiques de genre $g$ sur une variété $V$ est un
morphisme propre et plat\footnote{On peut supprimer cette hypothèse sans changer
la réponse à la question.} $X \to Y$ tel que toutes les fibres
schématiques de $X \to V$ soient des courbes hyperelliptiques projectives lisses
de genre $g$.
\end{enumerate}
Montrer que le nombre de modules est $2g - 1$.
\end{exercise}





\section{Ext globaux}
\label{section-global-exts}

\begin{exercise}
\label{exercise-ext-line-plane-p3}
Soit $k$ un corps. Soit $X = \mathbf{P}^3_k$. Soient $L \subset X$
et $P \subset X$ respectivement une droite et un plan, considérés comme des sous-schémas fermés
définis respectivement par $1$ et par $2$ équations linéaires. Calculer
les dimensions de
$$
\Ext^i_X(\mathcal{O}_L, \mathcal{O}_P)
$$
pour tout $i$.
Traiter à la fois le cas où $L$ est contenu dans $P$ et
celui où $L$ n'est pas contenu dans $P$.
\end{exercise}

\begin{exercise}
\label{exercise-ext-point}
Soit $k$ un corps. Soit $X = \mathbf{P}^n_k$. Soit $Z \subset X$
un point fermé rationnel sur $k$, considéré comme un sous-schéma fermé.
Par exemple, le point de coordonnées homogènes
$(1 : 0 : \ldots : 0)$. Calculer les dimensions de
$$
\Ext^i_X(\mathcal{O}_Z, \mathcal{O}_Z)
$$
pour tout $i$.
\end{exercise}

\begin{exercise}
\label{exercise-ext-pairings}
Soit $X$ un espace annelé. Définir les applications de cup-produit
$$
\Ext^i_X(\mathcal{G}, \mathcal{H})
\times
\Ext^j_X(\mathcal{F}, \mathcal{G})
\longrightarrow
\Ext^{i + j}_X(\mathcal{F}, \mathcal{H})
$$
pour des $\mathcal{O}_X$-modules $\mathcal{F}, \mathcal{G}, \mathcal{H}$.
(Indication : il s'agit d'une construction extrêmement générale.)
\end{exercise}

\begin{exercise}
\label{exercise-move-out-locally-free}
Soit $X$ un espace annelé. Soit $\mathcal{E}$ un
$\mathcal{O}_X$-module localement libre de rang fini, de dual
$\mathcal{E}^\vee = \SheafHom_{\mathcal{O}_X}(\mathcal{E}, \mathcal{O}_X)$.
Démontrer les assertions suivantes
\begin{enumerate}
\item $\SheafExt^i_{\mathcal{O}_X}(
\mathcal{F} \otimes_{\mathcal{O}_X} \mathcal{E}, \mathcal{G}) =
\SheafExt^i_{\mathcal{O}_X}(
\mathcal{F},
\mathcal{E}^\vee \otimes_{\mathcal{O}_X} \mathcal{G}) =
\SheafExt^i_{\mathcal{O}_X}(
\mathcal{F}, \mathcal{G}) \otimes_{\mathcal{O}_X} \mathcal{E}^\vee$, et
\item $\Ext^i_X(
\mathcal{F} \otimes_{\mathcal{O}_X} \mathcal{E}, \mathcal{G}) =
\Ext^i_X(\mathcal{F},
\mathcal{E}^\vee \otimes_{\mathcal{O}_X} \mathcal{G})$.
\end{enumerate}
Ici, $\mathcal{F}$ et $\mathcal{G}$ sont des $\mathcal{O}_X$-modules. En déduire
que
$$
\Ext^i_X(\mathcal{E}, \mathcal{G}) =
H^i(X, \mathcal{E}^\vee \otimes_{\mathcal{O}_X} \mathcal{G})
$$
\end{exercise}

\begin{exercise}
\label{exercise-trace-on-ext}
Soit $X$ un espace annelé. Soit $\mathcal{E}$ un
$\mathcal{O}_X$-module localement libre de rang fini. Construire un homomorphisme trace
$$
\Ext^i_X(\mathcal{E}, \mathcal{E}) \to H^i(X, \mathcal{O}_X)
$$
pour tout $i$. Le généraliser en un homomorphisme trace
$$
\Ext^i_X(\mathcal{E}, \mathcal{E} \otimes_{\mathcal{O}_X} \mathcal{F})
\to H^i(X, \mathcal{F})
$$
pour tout $\mathcal{O}_X$-module $\mathcal{F}$.
\end{exercise}

\begin{exercise}
\label{exercise-duality}
Soit $k$ un corps. Soit $X = \mathbf{P}^d_k$. Posons
$\omega_{X/k} = \mathcal{O}_X(-d - 1)$.
Démontrer que, pour des modules localement libres de rang fini
$\mathcal{E}$, $\mathcal{F}$, le cup-produit sur les Ext
composé avec l'homomorphisme trace sur les Ext
$$
\Ext^i_X(\mathcal{E}, \mathcal{F} \otimes_{\mathcal{O}_X} \omega_{X/k})
\times
\Ext^{d - i}_X(\mathcal{F}, \mathcal{E})
\to
\Ext_X^d(\mathcal{F}, \mathcal{F} \otimes_{\mathcal{O}_X} \omega_{X/k})
\to
H^d(X, \omega_{X/k}) = k
$$
définit un accouplement non dégénéré. Indication : on peut soit redémontrer la dualité
dans ce cadre, soit se ramener à la cohomologie des faisceaux et appliquer le
théorème de dualité de Serre démontré dans le cours.
\end{exercise}







\section{Diviseurs}
\label{section-divisors}

\noindent
Nous rassemblons ici, pour plus de commodité, toutes les définitions utiles.

\begin{definition}
\label{definition-divisor}
Dans toute cette définition, soit $S$ un schéma quelconque et soit
$X$ un schéma noethérien intègre.
\begin{enumerate}
\item Un {\it diviseur de Weil} sur $X$ est une combinaison linéaire formelle
$\Sigma n_i[Z_i]$ de diviseurs premiers $Z_i$ à coefficients entiers.
\item Un {\it diviseur premier} est un sous-schéma fermé $Z \subset X$,
intègre et de point générique $\xi \in Z$, tel que
${\mathcal O}_{X, \xi}$ soit de dimension $1$. Nous utiliserons la notation
${\mathcal O}_{X, Z} = {\mathcal O}_{X, \xi}$
lorsque $\xi \in Z \subset X$ est comme ci-dessus. Noter que ${\mathcal O}_{X, Z} \subset
K(X)$ est un sous-anneau du corps des fonctions de $X$.
\item Le {\it diviseur de Weil associé à une fonction rationnelle
$f \in K(X)^\ast$} est la somme $\Sigma v_Z(f)[Z]$. Ici, $v_Z(f)$ est
défini comme suit
\begin{enumerate}
\item Si $f \in {\mathcal O}_{X, Z}^\ast$, alors $v_Z(f) = 0$.
\item Si $f \in {\mathcal O}_{X, Z}$, alors
$$
v_Z(f) = \text{length}_{{\mathcal O}_{X, Z}}({\mathcal O}_{X, Z}/(f)).
$$
\item Si $f = \frac{a}{b}$ avec $a, b \in {\mathcal O}_{X, Z}$,
alors
$$
v_Z(f) = \text{length}_{{\mathcal O}_{X, Z}}({\mathcal O}_{X, Z}/(a)) -
\text{length}_{{\mathcal O}_{X, Z}}({\mathcal O}_{X, Z}/(b)).
$$
\end{enumerate}
\item Un {\it diviseur de Cartier effectif} sur un schéma $S$
est un sous-schéma fermé $D \subset S$ tel que tout point $d\in D$
possède un voisinage ouvert affine $\Spec(A) = U \subset S$ dans $S$
tel que $D \cap U = \Spec(A/(f))$, où $f \in A$ est un élément régulier.
\item Le {\it diviseur de Weil $[D]$ associé à un diviseur
de Cartier effectif $D \subset X$} de notre schéma noethérien intègre
$X$ est défini comme la somme $\Sigma v_Z(D)[Z]$, où
$v_Z(D)$ est défini comme suit
\begin{enumerate}
\item Si le point générique $\xi$ de $Z$ n'appartient pas à $D$,
alors $v_Z(D) = 0$.
\item Si le point générique $\xi$ de $Z$ appartient à $D$,
alors
$$
v_Z(D) = \text{length}_{{\mathcal O}_{X, Z}}({\mathcal O}_{X, Z}/(f))
$$
où $f \in {\mathcal O}_{X, Z} = {\mathcal O}_{X, \xi}$ est l'élément régulier
qui définit $D$ dans un voisinage affine de $\xi$ (comme en (4) ci-dessus).
\end{enumerate}
\item Soit $S$ un schéma. Le {\it faisceau des anneaux totaux de
fractions ${\mathcal K}_S$} est le faisceau de ${\mathcal O}_S$-algèbres obtenu par
le faisceau associé au préfaisceau ${\mathcal K}'$ défini comme suit.
Pour $U \subset S$ ouvert, posons ${\mathcal K}'(U) = S_U^{-1}{\mathcal O}_S(U)$,
où $S_U \subset {\mathcal O}_S(U)$ est la partie multiplicative
formée des sections $f \in {\mathcal O}_S(U)$ telles que le germe
de $f$ dans ${\mathcal O}_{S, u}$ soit régulier pour tout $u\in U$.
En particulier, tous les éléments de $S_U$ sont réguliers.
Ainsi, ${\mathcal O}_S$ est un sous-faisceau de ${\mathcal K}_S$, et l'on obtient une
suite exacte
$$
0 \to {\mathcal O}_S^\ast \to {\mathcal K}_S^\ast \to
{\mathcal K}_S^\ast/{\mathcal O}_S^\ast \to 0.
$$
\item Un {\it diviseur de Cartier} sur un schéma $S$ est une section globale
du faisceau quotient ${\mathcal K}_S^\ast/{\mathcal O}_S^\ast$.
\item Le {\it diviseur de Weil associé à un diviseur de Cartier}
$\tau \in \Gamma(X, {\mathcal K}_X^\ast/{\mathcal O}_X^\ast)$ sur notre schéma
noethérien intègre
$X$ est la somme $\Sigma v_Z(\tau)[Z]$, où $v_Z(\tau)$ est défini
par la procédure suivante
\begin{enumerate}
\item Si le germe de $\tau$ au point générique $\xi$
de $Z$ est nul -- autrement dit, si l'image de $\tau$ dans la fibre
$({\mathcal K}^\ast/{\mathcal O}^\ast)_\xi$ est ``nulle'' --, alors $v_Z(\tau) = 0$.
\item Choisir un voisinage ouvert affine $\Spec(A) = U \subset X$
tel que $\tau|_U$ soit l'image d'une section $f \in {\mathcal K}(U)$
et que, de plus, $f = a/b$ avec $a, b \in A$. On pose alors
$$
v_Z(f) = \text{length}_{{\mathcal O}_{X, Z}}({\mathcal O}_{X, Z}/(a)) -
\text{length}_{{\mathcal O}_{X, Z}}({\mathcal O}_{X, Z}/(b)).
$$
\end{enumerate}
\end{enumerate}
\end{definition}

\begin{remarks}
\label{remarks-divisors}
Voici quelques remarques immédiates.
\begin{enumerate}
\item Sur un schéma noethérien intègre $X$, le
faisceau ${\mathcal K}_X$ est constant de valeur le corps des fonctions $K(X)$.
\item Pour que les définitions ci-dessus aient un sens, il faut
montrer que
$$
\text{length}_{\mathcal O}({\mathcal O}/(ab)) =
\text{length}_{\mathcal O}({\mathcal O}/(a)) +
\text{length}_{\mathcal O}({\mathcal O}/(b))
$$
pour toute paire $(a, b)$ d'éléments non nuls d'un anneau local noethérien intègre
de dimension 1, noté ${\mathcal O}$. Cela sera démontré dans le cours.
\end{enumerate}
\end{remarks}

\begin{exercise}
\label{exercise-effective-cartier-cartier}
(Sur un schéma quelconque.)
Décrire comment associer un diviseur de Cartier à un diviseur de Cartier effectif.
\end{exercise}

\begin{exercise}
\label{exercise-rational-function-cartier}
(Sur un schéma intègre.)
Décrire comment associer un diviseur de Cartier $D$ à une fonction rationnelle
$f$ de telle sorte que les diviseurs de Weil associés à $D$ et à $f$ coïncident.
(C'est tautologique.)
\end{exercise}

\begin{exercise}
\label{exercise-weil-not-cartier}
Donner un exemple de diviseur de Weil sur une variété qui ne soit
associé à aucun diviseur de Cartier.
\end{exercise}

\begin{exercise}
\label{exercise-weil-Q-cartier}
Donner un exemple de diviseur de Weil $D$ sur une variété qui ne soit
associé à aucun diviseur de Cartier, mais
tel que $nD$ soit associé à un diviseur de Cartier
pour un certain $n > 1$.
\end{exercise}

\begin{exercise}
\label{exercise-weil-not-Q-cartier}
Donner un exemple de diviseur de Weil $D$ sur une variété qui ne soit
associé à aucun diviseur de Cartier et
tel que $nD$ NE SOIT associé à aucun diviseur de Cartier
pour tout $n > 1$.
(Indication : considérer un cône, par exemple $X : xy - zw = 0$ dans
$\mathbf{A}^4_k$. Essayer de montrer que $D = [x = 0, z = 0]$ convient.)
\end{exercise}

\begin{exercise}
\label{exercise-cartier-not-difference-effective-cartier}
Sur un schéma séparé $X$ de type fini sur un corps :
donner un exemple de diviseur de Cartier qui ne soit pas la différence de
deux diviseurs de Cartier effectifs.
Indication : trouver un $X$ qui ne possède aucun diviseur de Cartier effectif non vide,
par exemple le schéma construit dans
\cite[III, exercice 5.9]{H}. Il existe même un exemple
où $X$ est une variété : c'est la variété de
l'exercice \ref{exercise-nonprojective}.
\end{exercise}

\begin{exercise}
\label{exercise-nonprojective}
Exemple de variété propre non projective.
Soit $k$ un corps. Soit $L \subset \mathbf{P}^3_k$ une droite et
soit $C \subset \mathbf{P}^3_k$ une conique non singulière. Supposons que
$C \cap L = \emptyset$. Choisissons un
isomorphisme $\varphi : L \to C$. Soit $X$ la $k$-variété obtenue
en recollant $C$ à $L$ au moyen de $\varphi$. Autrement dit, il existe
un morphisme propre birationnel surjectif
$$
\pi : \mathbf{P}^3_k \longrightarrow X
$$
et un ouvert $U \subset X$ tel que $\pi : \pi^{-1}(U) \to U$ soit
un isomorphisme, $\pi^{-1}(U) = \mathbf{P}^3_k \setminus (L \cup C)$,
et tel que $\pi|_L = \pi|_C \circ \varphi$. (Ces conditions ne
définissent pas encore $X$ de façon unique. Pour cela, il faut préciser le
faisceau structural de $X$ au voisinage des points de $Z = X \setminus U$.)
Montrer que $X$ existe et que c'est une variété propre, mais non projective.
(Indication : pour l'existence, utiliser le résultat de
l'exercice \ref{exercise-prepare-glueing}. Pour la non-projectivité, utiliser
$\Pic(\mathbf{P}^3_k) = \mathbf{Z}$ pour montrer que $X$ ne peut pas posséder
de faisceau inversible ample.)
\end{exercise}



\section{Différentielles}
\label{section-differentials}

\noindent
{\bf Définitions et résultats.} Différentielles de K\"ahler.
\begin{enumerate}
\item Soit $R \to A$ un homomorphisme d'anneaux. Le {\it module des différentielles de K\"ahler
de $A$ sur $R$} se note $\Omega_{A/R}$.
Il est engendré par les éléments $\text{d}a$, $a \in A$,
soumis aux relations :
$$
\text{d}(a_1 + a_2) = \text{d}a_1 + \text{d}a_2,\quad
\text{d}(a_1a_2) = a_1\text{d}a_2 + a_2\text{d}a_1,\quad
\text{d}r = 0
$$
La $R$-dérivation universelle canonique $\text{d} : A \to \Omega_{A/R}$
envoie $a\mapsto \text{d}a$.
\item Considérons la suite exacte
$$
0 \to I \to A \otimes_R A \to A \to 0
$$
qui définit l'idéal $I$. Il existe une dérivation canonique
$\text{d} : A \to I/I^2$ qui envoie $a$ sur la classe de
$a \otimes 1 - 1 \otimes a$. C'est une autre présentation du
module des différentielles de $A$ sur $R$; autrement dit,
$$
(I/I^2, \text{d}) \cong (\Omega_{A/R}, \text{d}).
$$
\item Pour des parties multiplicatives $S_R \subset R$ et
$S_A \subset A$ telles que l'image de $S_R$ soit contenue dans $S_A$, on a
$$
\Omega_{S_A^{-1}A / S_R^{-1}R} =
S_A^{-1}\Omega_{A/R}.
$$
\item Si $A$ est une $R$-algèbre de présentation finie, alors
$\Omega_{A/R}$ est un $A$-module de présentation finie. Ainsi, dans
ce cas, les idéaux de {\it Fitting} de $\Omega_{A/R}$ sont définis.
\item Soit $f : X \to S$ un morphisme de schémas. Il existe
un faisceau quasi-cohérent de ${\mathcal O}_X$-modules $\Omega_{X/S}$
et une dérivation ${\mathcal O}_S$-linéaire
$$
\text{d} : {\mathcal O}_X \longrightarrow \Omega_{X/S}
$$
tels que, pour tous ouverts affines $\Spec(A) = U \subset X$,
$\Spec(R) = V \subset S$,
avec $f(U) \subset V$, on ait
$$
\Gamma(\Spec(A), \Omega_{X/S}) = \Omega_{A/R}
$$
d'une manière compatible avec $\text{d}$.
\end{enumerate}

\begin{exercise}
\label{exercise-dual-numbers}
Soit $k[\epsilon]$ l'anneau des nombres duaux
sur le corps $k$, c'est-à-dire que $\epsilon^2 = 0$.
\begin{enumerate}
\item Considérons l'homomorphisme d'anneaux
$$
R = k[\epsilon] \to A = k[x, \epsilon]/(\epsilon x)
$$
Montrer que les idéaux de Fitting de $\Omega_{A/R}$ sont (en commençant par
l'idéal de Fitting d'indice zéro)
$$
(\epsilon), A, A, \ldots
$$
\item Considérons l'homomorphisme $R = k[t] \to
A = k[x, y, t]/(x(y-t)(y-1), x(x-t))$. Montrer que les idéaux de Fitting de
$\Omega_{A/R}$ dans $A$ sont (supposer, pour simplifier, que $k$ est de caractéristique
nulle)
$$
x(2x-t)(2y-t-1)A, \ (x, y, t)\cap (x, y-1, t), \ A, \ A, \ldots
$$
Ainsi, l'idéal de Fitting d'indice $0$ est défini par un unique élément de $A$,
l'idéal de Fitting d'indice $1$ définit deux points fermés de $\Spec(A)$, et
tous les autres sont triviaux.
\item Considérons l'homomorphisme $R = k[t] \to A = k[x, y, t]/(xy-t^n)$.
Calculer les idéaux de Fitting de $\Omega_{A/R}$.
\end{enumerate}
\end{exercise}

\begin{remark}
\label{remark-fitting-omega-not-sings}
On utilise couramment le $k$-ième idéal de Fitting de $\Omega_{X/S}$
pour définir le schéma des singularités du morphisme $X \to S$ lorsque $X$ est de dimension relative
$k$ sur $S$. Toutefois, comme le montre la partie (a), cette définition demande de la prudence
lorsque la dimension des fibres de la famille n'est pas ``constante'', par exemple lorsque
la famille n'est pas plate. Comme le montre la partie (b), la platitude ne suffit pas non plus
(et oui, cette famille est plate). Dans les ``bons cas'' -- comme en (c) --, pour les
familles de courbes, on s'attend à ce que l'idéal de Fitting d'indice $0$ soit nul et
que celui d'indice $1$ définisse, schématiquement, le lieu singulier.
\end{remark}

\begin{exercise}
\label{exercise-formally-smooth}
Supposons que $R$ soit un anneau et que
$$
A = R[x_1, \ldots, x_n]/(f_1, \ldots, f_n).
$$
Noter que nous supposons $A$ présenté par autant
d'équations que de variables. La matrice des dérivées
partielles
$$
( \partial f_i / \partial x_j )
$$
est $n \times n$, donc carrée. Supposons que
son déterminant soit inversible comme élément de $A$. Noter que
c'est exactement la condition qui exprime que $\Omega_{A/R} = (0)$
dans ce cas de $n$ générateurs et de $n$ relations.
Soit $\pi : B' \to B$ une surjection de $R$-algèbres
dont le noyau $J$ est de carré nul (comme idéal de $B'$).
Soit $\varphi : A \to B$ un homomorphisme de $R$-algèbres.
Montrer qu'il existe un unique homomorphisme de $R$-algèbres
$\varphi' : A \to B'$ tel que $\varphi = \pi \circ \varphi'$.
\end{exercise}

\begin{exercise}
\label{exercise-formally-smooth-one-equation}
Généraliser le résultat de l'exercice \ref{exercise-formally-smooth}
au cas où $A = R[x, y]/(f)$.
\end{exercise}

\begin{exercise}
\label{exercise-Jacobian-criterion}
Soit $k$ un corps, soient $f_1, \ldots, f_c \in k[x_1, \ldots, x_n]$, et
soit $A = k[x_1, \ldots, x_n]/(f_1, \ldots, f_c)$.
Supposons que $f_j(0, \ldots, 0) = 0$. Cela signifie que
$\mathfrak m = (x_1, \ldots, x_n)A$ est un idéal maximal.
Démontrer que l'anneau local $A_\mathfrak m$ est régulier
si la matrice
$$
(\partial f_j/ \partial x_i)|_{(x_1, \ldots, x_n) = (0, \ldots, 0)}
$$
est de rang $c$. Quelle est alors la dimension de $A_\mathfrak m$ ?
Montrer que la réciproque est fausse en donnant un exemple où
$A_\mathfrak m$ est régulier mais où le rang est strictement inférieur à $c$;
quelle est la dimension de $A_\mathfrak m$ dans cet exemple ?
\end{exercise}




\section{Schémas, examen final, automne 2007}
\label{section-final-exam-fall-2007}

\noindent
Voici les questions de l'examen final d'un cours sur les schémas,
donné au printemps 2007 à l'Université Columbia.

\begin{exercise}[Définitions]
\label{exercise-definitions}
Donner la définition des notions suivantes.
\begin{enumerate}
\item $X$ est un {\it schéma}
\item le morphisme de schémas $f : X \to Y$ est {\it fini}
\item le morphisme de schémas $f : X \to Y$ est {\it de type fini}
\item le schéma $X$ est {\it noethérien}
\item le ${\mathcal O}_X$-module ${\mathcal L}$ sur
le schéma $X$ est {\it inversible}
\item le {\it genre} d'une courbe non singulière
projective sur un corps algébriquement clos
\end{enumerate}
\end{exercise}

\begin{exercise}
\label{exercise-kill-global-sections}
Soit $X = \Spec({\mathbf Z}[x, y])$, et soit ${\mathcal F}$ un
${\mathcal O}_X$-module quasi-cohérent.
Supposons que ${\mathcal F}$ soit nul après restriction à
cet
ouvert affine standard $D(x)$.
\begin{enumerate}
\item Montrer que toute section globale $s$ de ${\mathcal F}$ est annulée par une
puissance de $x$, c'est-à-dire $x^ns = 0$ pour un certain $n\in {\mathbf N}$.
\item Pensez-vous que le même énoncé reste vrai si l'on ne suppose pas que ${\mathcal F}$
est quasi-cohérent ?
\end{enumerate}
\end{exercise}

\begin{exercise}
\label{exercise-empty-fibre-empty}
Supposons que $X \to \Spec(R)$ soit un morphisme propre et que
$R$ soit un anneau de valuation discrète de corps résiduel $k$. Supposons que
$X \times_{\Spec(R)} \Spec(k)$ soit le schéma vide. Montrer que
$X$ est le schéma vide.
\end{exercise}

\begin{exercise}
\label{exercise-curve-p1-p1}
Considérons la
variété projective\footnote{Le plongement projectif est
$((X_0, X_1), (Y_0, Y_1))\mapsto (X_0Y_0, X_0Y_1, X_1Y_0, X_1Y_1)$
ou, autrement dit, $(x, y)\mapsto (1, y, x, xy)$.}
$$
{\mathbf P}^1 \times {\mathbf P}^1 = {\mathbf P}^1_{{\mathbf C}}
\times_{\Spec({\mathbf C})} {\mathbf P}^1_{\mathbf C}
$$
sur le corps des nombres complexes ${\mathbf C}$. Elle est recouverte par quatre morceaux
affines,
correspondant aux couples de morceaux affines standards de ${\mathbf P}^1_{\mathbf
C}$. Par exemple,
prenons les coordonnées homogènes $X_0, X_1$ sur le premier facteur et
$Y_0, Y_1$ sur le second. Posons $x = X_1/X_0$, et $y = Y_1/Y_0$. Les 4
morceaux affines sont alors les spectres des anneaux
$$
{\mathbf C}[x, y], \quad
{\mathbf C}[x^{-1}, y], \quad
{\mathbf C}[x, y^{-1}], \quad
{\mathbf C}[x^{-1}, y^{-1}].
$$
Soit $X \subset {\mathbf P}^1 \times {\mathbf P}^1$ le sous-schéma fermé
qui est
l'adhérence de la partie fermée du premier morceau affine définie par l'équation
$$
y^3(x^4 + 1) = x^4 -1.
$$
\begin{enumerate}
\item Montrer que $X$ est contenu dans la réunion du premier et
du dernier des 4 morceaux affines.
\item Montrer que $X$ est une courbe projective non singulière.
\item Considérons le morphisme $pr_2 : X \to {\mathbf P}^1$ (projection sur
le premier facteur). Sur le premier morceau affine, il est donné par $(x, y) \mapsto x$.
Expliquer brièvement pourquoi il est de degré $3$.
\item Calculer les points de ramification et les indices de ramification
du morphisme $pr_2 : X \to {\mathbf P}^1$.
\item Calculer le genre de $X$.
\end{enumerate}
\end{exercise}

\begin{exercise}
\label{exercise-finite-type-over-Z}
Soit $X \to \Spec({\mathbf Z})$ un morphisme de type fini.
Supposons qu'il existe une infinité de nombres premiers $p$ tels que
$X \times_{\Spec({\mathbf Z})} \Spec({\mathbf F}_p)$ ne soit pas vide.
\begin{enumerate}
\item Montrer que $X \times_{\Spec({\mathbf Z})}\Spec(\mathbf{Q})$
n'est pas vide.
\item Pensez-vous que le même énoncé reste vrai si l'on remplace la condition
« de type fini » par la condition « localement de type fini » ?
\end{enumerate}
\end{exercise}




\section{Schémas, examen final, printemps 2009}
\label{section-final-exam-spring-2009}

\noindent
Voici les questions de l'examen final d'un cours sur les schémas,
donné au printemps 2009 à l'Université Columbia.

\begin{exercise}
\label{exercise-Noetherian-coherent}
Soit $X$ un schéma noethérien.
Soit $\mathcal{F}$ un faisceau cohérent sur $X$.
Soit $x \in X$ un point.
Supposons que $\text{Supp}(\mathcal{F}) = \{ x \}$.
\begin{enumerate}
\item Montrer que $x$ est un point fermé de $X$.
\item Montrer que $H^0(X, \mathcal{F})$ n'est pas nul.
\item Montrer que $\mathcal{F}$ est engendré par ses sections globales.
\item Montrer que $H^p(X, \mathcal{F}) = 0$ pour $p > 0$.
\end{enumerate}
\end{exercise}

\begin{remark}
\label{remark-invertible-projective-space}
Soit $k$ un corps.
Soit $\mathbf{P}^2_k = \text{Proj}(k[X_0, X_1, X_2])$.
Tout faisceau inversible sur $\mathbf{P}^2_k$ est isomorphe à
$\mathcal{O}_{\mathbf{P}^2_k}(n)$ pour un certain $n \in \mathbf{Z}$.
Rappelons que
$$
\Gamma(\mathbf{P}^2_k, \mathcal{O}_{\mathbf{P}^2_k}(n)) =
k[X_0, X_1, X_2]_n
$$
est la composante homogène de degré $n$ de l'anneau de polynômes.
Pour un faisceau quasi-cohérent $\mathcal{F}$ sur $\mathbf{P}^2_k$, posons
$\mathcal{F}(n) =
\mathcal{F}
\otimes_{\mathcal{O}_{\mathbf{P}^2_k}}
\mathcal{O}_{\mathbf{P}^2_k}(n)$
comme d'habitude.
\end{remark}

\begin{exercise}
\label{exercise-nonsplit-vectorbundle}
Soit $k$ un corps.
Soit $\mathcal{E}$ un fibré vectoriel sur $\mathbf{P}^2_k$, c'est-à-dire un
$\mathcal{O}_{\mathbf{P}^2_k}$-module localement libre de rang fini.
On dit que $\mathcal{E}$ est {\it scindé} si $\mathcal{E}$ est isomorphe à
une somme directe de $\mathcal{O}_{\mathbf{P}^2_k}$-modules inversibles.
\begin{enumerate}
\item Montrer que $\mathcal{E}$ est scindé si et seulement si $\mathcal{E}(n)$
est scindé.
\item Montrer que, si $\mathcal{E}$ est scindé, alors
$H^1({\mathbf{P}^2_k}, \mathcal{E}(n)) = 0$
pour tout $n \in \mathbf{Z}$.
\item Soit
$$
\varphi :
\mathcal{O}_{\mathbf{P}^2_k}
\longrightarrow
\mathcal{O}_{\mathbf{P}^2_k}(1) \oplus
\mathcal{O}_{\mathbf{P}^2_k}(1) \oplus
\mathcal{O}_{\mathbf{P}^2_k}(1)
$$
défini par les formes linéaires
$L_0, L_1, L_2 \in \Gamma(\mathbf{P}^2_k, \mathcal{O}_{\mathbf{P}^2_k}(1))$.
Supposons que $L_i \not = 0$ pour un certain $i$.
Quelle condition doivent vérifier $L_0, L_1, L_2$ pour que
le conoyau de $\varphi$ soit un fibré vectoriel ?
Pourquoi ?
\item Donner un exemple d'un tel $\varphi$.
\item Montrer que $\Coker(\varphi)$ n'est pas scindé (lorsqu'il s'agit
d'un fibré vectoriel).
\end{enumerate}
\end{exercise}

\begin{remark}
\label{remark-recall-dimension-theory}
On pourra utiliser librement les faits suivants de théorie de la dimension
(et en ajouter d'autres au besoin).
\begin{enumerate}
\item La dimension d'un schéma est la borne supérieure des longueurs des chaînes
de parties fermées irréductibles.
\item La dimension d'un schéma de type fini sur un corps est le maximum
des dimensions de ses ouverts affines.
\item La dimension d'un schéma noethérien est le maximum des dimensions
de ses composantes irréductibles.
\item La dimension d'un schéma affine
coïncide avec la dimension de l'anneau correspondant.
\item Soit $k$ un corps et soit $A$ une $k$-algèbre de type fini.
Si $A$ est intègre et $x \not = 0$, alors $\dim(A) = \dim(A/xA) + 1$.
\end{enumerate}
\end{remark}

\begin{exercise}
\label{exercise-irreducible-fibres-same-dimension-irreducible}
Soit $k$ un corps.
Soit $X$ un schéma projectif réduit sur $k$.
Soit $f : X \to \mathbf{P}^1_k$ un morphisme de schémas sur $k$.
Supposons qu'il existe un entier $d \geq 0$ tel que,
pour tout point $t \in \mathbf{P}^1_k$, la fibre $X_t = f^{-1}(t)$
soit irréductible de dimension $d$. (Rappelons qu'un espace irréductible
n'est pas vide.)
\begin{enumerate}
\item Montrer que $\dim(X) = d + 1$.
\item Soit $X_0 \subset X$ une composante irréductible de $X$ de dimension
$d + 1$. Démontrer que, pour tout $t \in \mathbf{P}^1_k$, la fibre
$X_{0, t}$ est de dimension $d$.
\item Que peut-on déduire de ce qui précède au sujet de $X_t$ et de $X_{0, t}$ ?
\item Montrer que $X$ est irréductible.
\end{enumerate}
\end{exercise}

\begin{remark}
\label{remark-chi}
Étant donnés un schéma projectif $X$ sur un corps $k$ et
un faisceau cohérent $\mathcal{F}$ sur $X$, nous posons
$$
\chi(X, \mathcal{F}) =
\sum\nolimits_{i \geq 0} (-1)^i\dim_k H^i(X, \mathcal{F}).
$$
\end{remark}

\begin{exercise}
\label{exercise-complete-intersection}
Soit $k$ un corps.
Écrivons $\mathbf{P}^3_k = \text{Proj}(k[X_0, X_1, X_2, X_3])$.
Soit $C \subset \mathbf{P}^3_k$ une
{\it courbe intersection complète de type $(5, 6)$}.
Cela signifie qu'il existe $F \in k[X_0, X_1, X_2, X_3]_5$ et
$G \in k[X_0, X_1, X_2, X_3]_6$ tels que
$$
C = \text{Proj}(k[X_0, X_1, X_2, X_3]/(F, G))
$$
soit une variété de dimension $1$. (Une variété est en particulier réduite et irréductible,
mais on peut, si on le souhaite, supposer que $C$ est non singulière.)
Soit $i : C \to \mathbf{P}^3_k$ l'immersion fermée associée.
Le fait d'être une intersection complète implique aussi que
$$
\xymatrix{
0 \ar[r] &
\mathcal{O}_{\mathbf{P}^3_k}(-11)
\ar[r]^-{
\left(
\begin{matrix}
-G \\
F
\end{matrix}
\right)
} &
\mathcal{O}_{\mathbf{P}^3_k}(-5) \oplus \mathcal{O}_{\mathbf{P}^3_k}(-6)
\ar[r]^-{(F, G)} &
\mathcal{O}_{\mathbf{P}^3_k} \ar[r] &
i_*\mathcal{O}_C \ar[r] &
0
}
$$
est une suite exacte de faisceaux. Utiliser ces faits pour :
\begin{enumerate}
\item calculer $\chi(C, i^*\mathcal{O}_{\mathbf{P}^3_k}(n))$ pour
tout $n \in \mathbf{Z}$, et
\item calculer la dimension de $H^1(C, \mathcal{O}_C)$.
\end{enumerate}
\end{exercise}

\begin{exercise}
\label{exercise-glueing}
Soit $k$ un corps.
Considérons les anneaux
\begin{align*}
A & = k[x, y]/(xy) \\
B & = k[u, v]/(uv) \\
C & = k[t, t^{-1}] \times k[s, s^{-1}]
\end{align*}
et les morphismes de $k$-algèbres
$$
\begin{matrix}
A \longrightarrow C, &
x \mapsto (t, 0), &
y \mapsto (0, s) \\
B \longrightarrow C, &
u \mapsto (t^{-1}, 0), &
v \mapsto (0, s^{-1})
\end{matrix}
$$
Ces morphismes induisent effectivement des isomorphismes
$A_{x + y} \to C$ et $B_{u + v} \to C$. Par conséquent, les morphismes $A \to C$
et $B \to C$ identifient $\Spec(C)$ à des ouverts de
$\Spec(A)$ et de $\Spec(B)$. Soit $X$ le schéma obtenu
en recollant $\Spec(A)$ et $\Spec(B)$ le long de $\Spec(C)$ :
$$
X = \Spec(A) \amalg_{\Spec(C)} \Spec(B).
$$
Comme nous l'avons vu dans le cours, un tel schéma existe et possède des
ouverts affines $\Spec(A) \subset X$ et $\Spec(B) \subset X$
dont l'intersection est exactement $\Spec(C)$, identifié à un ouvert de
chacun d'eux au moyen des morphismes ci-dessus.
\begin{enumerate}
\item Pourquoi $X$ est-il séparé ?
\item Pourquoi $X$ est-il de type fini sur $k$ ?
\item Calculer $H^1(X, \mathcal{O}_X)$, ou, autrement dit, quelle est sa dimension ?
\item Quelle description plus géométrique peut-on donner de $X$ ?
\end{enumerate}
\end{exercise}





\section{Schémas, examen final, automne 2010}
\label{section-final-exam-fall-2010}

\noindent
Voici les questions de l'examen final d'un cours sur les schémas,
donné à l'automne 2010 à l'Université Columbia.

\begin{exercise}[Définitions]
\label{exercise-definitions-fall-2010}
Donner la définition des notions suivantes.
\begin{enumerate}
\item un schéma séparé,
\item un morphisme quasi-compact de schémas,
\item un morphisme affine de schémas,
\item une partie multiplicative d'un anneau,
\item un schéma noethérien,
\item une variété.
\end{enumerate}
\end{exercise}

\begin{exercise}
\label{exercise-prime-avoidance}
Évitement des idéaux premiers.
\begin{enumerate}
\item Soit $A$ un anneau. Soit $I \subset A$ un idéal et
soient $\mathfrak q_1$, $\mathfrak q_2$ des idéaux premiers tels que
$I \not \subset \mathfrak q_i$. Montrer que
$I \not \subset \mathfrak q_1 \cup \mathfrak q_2$.
\item Quelle est une interprétation géométrique de (1) ?
\item Soit $X = \text{Proj}(S)$ pour un anneau gradué $S$.
Soient $x_1, x_2 \in X$. Montrer qu'il existe un ouvert standard
$D_{+}(F)$ qui contient à la fois $x_1$ et $x_2$.
\end{enumerate}
\end{exercise}

\begin{exercise}
\label{exercise-compose-affine}
Pourquoi le composé de morphismes affines est-il affine ?
\end{exercise}

\begin{exercise}[Exemples]
\label{exercise-examples}
Donner des exemples des objets suivants :
\begin{enumerate}
\item Un schéma projectif réductible sur un corps $k$.
\item Un schéma à 100 points.
\item Un morphisme de schémas qui n'est pas affine.
\end{enumerate}
\end{exercise}

\begin{exercise}
\label{exercise-chevalley-hilbert-nullstellensatz}
Le théorème de Chevalley et le théorème des zéros de Hilbert.
\begin{enumerate}
\item Soit $\mathfrak p \subset \mathbf{Z}[x_1, \ldots, x_n]$
un idéal maximal. Que déduit-on du théorème de Chevalley au sujet de
$\mathfrak p \cap \mathbf{Z}$ ?
\item Que déduit-on alors du théorème des zéros de Hilbert au sujet de
$\kappa(\mathfrak p)$ ?
\end{enumerate}
\end{exercise}

\begin{exercise}
\label{exercise-P0}
Soit $A$ un anneau. Soit $S = A[X]$ une $A$-algèbre graduée
où $X$ est de degré $1$. Montrer que $\text{Proj}(S) \cong \Spec(A)$
en tant que schémas sur $A$.
\end{exercise}

\begin{exercise}
\label{exercise-finite-is-projective}
Soit $A \to B$ un homomorphisme fini d'anneaux. Montrer que
$\Spec(B)$ est un schéma H-projectif sur $\Spec(A)$.
\end{exercise}

\begin{exercise}
\label{exercise-not-geometrically-irreducible}
Donner un exemple d'un schéma $X$ sur un corps $k$ tel que
$X$ soit irréductible et que, pour une certaine extension finie $k'/k$,
le changement de base $X_{k'} = X \times_{\Spec(k)} \Spec(k')$
soit connexe mais réductible.
\end{exercise}





\section{Schémas, examen final, printemps 2011}
\label{section-final-exam-spring-2011}

\noindent
Voici les questions de l'examen final d'un cours sur les schémas,
donné au printemps 2011 à l'Université Columbia.

\begin{exercise}[Définitions]
\label{exercise-definitions-spring-2011}
Donner la définition des notions en italique.
\begin{enumerate}
\item un schéma {\it séparé},
\item un morphisme de schémas {\it universellement fermé},
\item {\it $A$ domine $B$} pour des anneaux locaux $A, B$ contenus dans un
même corps,
\item la {\it dimension} d'un schéma $X$,
\item la {\it codimension} d'un sous-schéma fermé irréductible $Y$
d'un schéma $X$,
\end{enumerate}
\end{exercise}

\begin{exercise}[Résultats]
\label{exercise-results-spring-2011}
Pour chacun des points suivants, énoncer un résultat formellement équivalent au fait étudié
dans le cours.
\begin{enumerate}
\item Le critère valuatif de propreté pour un morphisme
$X \to Y$ de variétés, par exemple.
\item La relation entre $\dim(X)$ et le corps des fonctions
$k(X)$ de $X$ pour une variété $X$ sur un corps $k$.
\item Compléter : La catégorie des courbes projectives non singulières sur
$k$ et des morphismes non constants est antiéquivalente à $\ldots\ldots\ldots$.
\item La normalisation de Noether.
\item Le critère jacobien.
\end{enumerate}
\end{exercise}

\begin{exercise}
\label{exercise-genus-plane-curve}
Soit $k$ un corps.
Soit $F \in k[X_0, X_1, X_2]$ une forme homogène de degré $d$.
Supposons que $C = V_{+}(F) \subset \mathbf{P}^2_k$ soit une courbe lisse sur $k$.
Notons $i : C \to \mathbf{P}^2_k$ l'immersion fermée associée.
\begin{enumerate}
\item Montrer qu'il existe une suite exacte courte
$$
0 \to \mathcal{O}_{\mathbf{P}^2_k}(-d)
\to \mathcal{O}_{\mathbf{P}^2_k} \to i_*\mathcal{O}_C \to 0
$$
de faisceaux cohérents sur $\mathbf{P}^2_k$ : indiquer quelles sont les flèches
et expliquer brièvement pourquoi cette suite est exacte.
\item En déduire que $H^0(C, \mathcal{O}_C) = k$.
\item Calculer le genre de $C$.
\item Supposons maintenant que $P = (0 : 0 : 1)$ n'appartienne pas à $C$. Démontrer que
$\pi : C \to \mathbf{P}^1_k$ défini par $(a_0 : a_1 : a_2) \mapsto (a_0 : a_1)$
est de degré $d$.
\item Supposons que $k$ soit algébriquement clos, que tous les indices de ramification
(les ``$e_i$'') soient égaux à $1$ ou $2$, et que la caractéristique de $k$
ne soit pas égale à $2$. Combien de points de ramification
le morphisme $\pi : C \to \mathbf{P}^1_k$ possède-t-il ?
\item En fonction de $F$, quelles équations proposeriez-vous pour
l'ensemble des points de ramification de $\pi$ ?
\item Pouvez-vous deviner $K_C$ ?
\end{enumerate}
\end{exercise}

\begin{exercise}
\label{exercise-Pic-triangle}
Soit $k$ un corps. Soit $X$ un ``triangle'' sur $k$, c'est-à-dire
que l'on obtient $X$ en recollant trois copies de $\mathbf{A}^1_k$ les unes aux autres en
identifiant $0$ de la première copie à $1$ de la deuxième, $0$ de la
deuxième copie à $1$ de la troisième et $0$ de la troisième à $1$ de
la première. Il se trouve que $X$ est isomorphe à
$\Spec(k[x, y]/(xy(x + y + 1)))$ ; on pourra utiliser ce fait.
Calculer le groupe de Picard de $X$.
\end{exercise}

\begin{exercise}
\label{exercise-birational-morphism-curves-ample}
Soit $k$ un corps. Soit $\pi : X \to Y$ un morphisme fini birationnel
de courbes, où $X$ est une courbe projective non singulière sur $k$.
Il résulte du cours que $Y$ est une courbe propre et
que $\pi$ est le morphisme de normalisation de $Y$.
Nous avons aussi vu dans le cours qu'il existe un ouvert dense $V \subset Y$
tel que $U = \pi^{-1}(V)$ soit un ouvert dense de $X$ et que $\pi : U \to V$
soit un isomorphisme.
\begin{enumerate}
\item Montrer qu'il existe un diviseur de Cartier effectif $D \subset X$
tel que $D \subset U$ et que $\mathcal{O}_X(D)$ soit ample sur $X$.
\item Soit $D$ comme en (1). Montrer que $E = \pi(D)$ est un diviseur de Cartier
effectif sur $Y$.
\item Indiquer brièvement pourquoi
\begin{enumerate}
\item le morphisme $\mathcal{O}_Y \to \pi_*\mathcal{O}_X$ a un conoyau cohérent
$Q$ dont le support est contenu dans $Y \setminus V$, et
\item pour tout $n$, il existe un morphisme correspondant
$\mathcal{O}_Y(nE) \to \pi_*\mathcal{O}_X(nD)$
dont le conoyau est isomorphe à $Q$.
\end{enumerate}
\item Montrer que
$\dim_k H^0(X, \mathcal{O}_X(nD)) - \dim_k H^0(Y, \mathcal{O}_Y(nE))$
est bornée (par quoi ?) et en déduire que le faisceau inversible
$\mathcal{O}_Y(nE)$ possède beaucoup de sections pour $n$ grand (pourquoi ?).
\end{enumerate}
\end{exercise}





\section{Schémas, examen final, automne 2011}
\label{section-final-exam-fall-2011}

\noindent
Ces questions ont été posées lors de l'examen final d'un cours
d'algèbre commutative, donné à l'automne 2011 à l'Université Columbia.


\begin{exercise}[Définitions]
\label{exercise-definitions-fall-2011}
Donner la définition des notions en italique.
\begin{enumerate}
\item un anneau {\it noethérien},
\item un schéma {\it noethérien},
\item un homomorphisme d'anneaux {\it fini},
\item un morphisme de schémas {\it fini},
\item la {\it dimension} d'un anneau.
\end{enumerate}
\end{exercise}

\begin{exercise}[Résultats]
\label{exercise-results-fall-2011}
Énoncer un résultat formellement équivalent à celui qui a été étudié
dans le cours.
\begin{enumerate}
\item Le théorème principal de Zariski.
\item La normalisation de Noether.
\item Le théorème chinois.
\item La montée pour les homomorphismes d'anneaux finis.
\end{enumerate}
\end{exercise}

\begin{exercise}
\label{exercise-dimension-of-ring}
Soit $(A, \mathfrak m, \kappa)$ un anneau local noethérien
dont le corps résiduel est de caractéristique différente de $2$.
Supposons que $\mathfrak m$ soit engendré par trois éléments
$x, y, z$ et que $x^2 + y^2 + z^2 = 0$ dans $A$.
\begin{enumerate}
\item Quelles sont les valeurs possibles de $\dim(A)$ ?
\item Donner un exemple montrant que chacune de ces valeurs est possible.
\item Montrer que $A$ est intègre si $\dim(A) = 2$.
(Indication : considérer
$\bigoplus_{n \geq 0} \mathfrak m^n/\mathfrak m^{n + 1}$.)
\end{enumerate}
\end{exercise}

\begin{exercise}
\label{exercise-localization}
Soit $A$ un anneau.
Soient $S \subset T \subset A$ des parties multiplicatives.
Supposons que
$$
\{\mathfrak q \mid \mathfrak q \cap S = \emptyset\} =
\{\mathfrak q \mid \mathfrak q \cap T = \emptyset\}.
$$
Montrer que $S^{-1}A \to T^{-1}A$ est un isomorphisme.
\end{exercise}

\begin{exercise}
\label{exercise-locus-of-rank-1}
Soit $k$ un corps algébriquement clos. Soit
$$
V_0 = \{ A \in \text{Mat}(3 \times 3, k) \mid \text{rank}(A) = 1\}
\subset \text{Mat}(3 \times 3, k) = k^9.
$$
\begin{enumerate}
\item Montrer que $V_0$ est l'ensemble des points fermés d'une
partie localement fermée (pour la topologie de Zariski) $V \subset \mathbf{A}^9_k$.
\item $V$ est-il irréductible ?
\item Que vaut $\dim(V)$ ?
\end{enumerate}
\end{exercise}

\begin{exercise}
\label{exercise-not-complete-intersection}
Démontrer que l'idéal $(x^2, xy, y^2)$ de $\mathbf{C}[x, y]$
ne peut pas être engendré par $2$ éléments.
\end{exercise}

\begin{exercise}
\label{exercise-finite-projection}
Soit $f \in \mathbf{C}[x, y]$ un polynôme non constant.
Montrer qu'il existe $\alpha, \beta \in \mathbf{C}$ tels que
l'homomorphisme de $\mathbf{C}$-algèbres
$$
\mathbf{C}[t] \longrightarrow \mathbf{C}[x, y]/(f),\quad
t \longmapsto \alpha x + \beta y
$$
soit fini.
\end{exercise}

\begin{exercise}
\label{exercise-union-of-two-affines}
Montrer que, pour tous points $p_1, \ldots, p_n \in \mathbf{C}^2$,
le schéma $\mathbf{A}^2_\mathbf{C} \setminus \{p_1, \ldots, p_n\}$
est réunion de deux ouverts affines.
\end{exercise}

\begin{exercise}
\label{exercise-surjection-A1-P1}
Montrer qu'il existe un morphisme surjectif de schémas
$\mathbf{A}^1_\mathbf{C} \to \mathbf{P}^1_\mathbf{C}$.
(Surjectif signifie simplement surjectif sur les ensembles de points sous-jacents.)
\end{exercise}

\begin{exercise}
\label{exercise-almost-surjective}
Soit $k$ un corps algébriquement clos. Soit $A \subset B$ une extension
d'anneaux intègres qui sont tous deux des $k$-algèbres de type fini.
Montrer que l'image de $\Spec(B) \to \Spec(A)$
contient un ouvert non vide de $\Spec(A)$ en suivant les étapes ci-dessous :
\begin{enumerate}
\item Le démontrer lorsque $A \to B$ est en outre fini.
\item Le démontrer lorsque le corps des fractions de $B$ est une extension finie
du corps des fractions de $A$.
\item Réduire l'énoncé au cas précédent.
\end{enumerate}
\end{exercise}





\section{Schémas, examen final, automne 2013}
\label{section-final-exam-fall-2013}

\noindent
Ces questions ont été posées lors de l'examen final d'un cours
d'algèbre commutative, donné à l'automne 2013 à l'Université Columbia.

\begin{exercise}[Définitions]
\label{exercise-definitions-fall-2013}
Donner la définition des notions en italique.
\begin{enumerate}
\item un {\it idéal radical} d'un anneau,
\item un homomorphisme d'anneaux {\it de type fini},
\item une {\it différentielle au sens de Weil},
\item un {\it schéma}.
\end{enumerate}
\end{exercise}

\begin{exercise}[Résultats]
\label{exercise-results-fall-2013}
Énoncer un résultat formellement équivalent à celui qui a été étudié
dans le cours.
\begin{enumerate}
\item un résultat sur les polynômes de Hilbert des modules gradués.
\item la dimension d'un anneau local noethérien $(R, \mathfrak m)$ et de
$\bigoplus_{n \geq 0} \mathfrak m^n/\mathfrak m^{n + 1}$.
\item Riemann-Roch.
\item le théorème de Clifford.
\item le théorème de Chevalley.
\end{enumerate}
\end{exercise}

\begin{exercise}
\label{exercise-surjective-after-localization}
Soit $A \to B$ un homomorphisme d'anneaux. Soit $S \subset A$ une partie multiplicative.
Supposons que $A \to B$ soit de type fini et que $S^{-1}A \to S^{-1}B$ soit
surjectif. Montrer qu'il existe $f \in S$ tel que $A_f \to B_f$
soit surjectif.
\end{exercise}

\begin{exercise}
\label{exercise-injective-local-ring-map-not-surjective}
Donner un exemple d'un homomorphisme local injectif $A \to B$
d'anneaux locaux tel que $\Spec(B) \to \Spec(A)$ ne soit pas surjectif.
\end{exercise}

\begin{situation}[Notation : courbe plane]
\label{situation-curve-in-the-plane}
Soit $k$ un corps algébriquement clos. Soit
$F(X_0, X_1, X_2) \in k[X_0, X_1, X_2]$ un polynôme irréductible
homogène de degré $d$. Posons
$$
D = V(F) \subset \mathbf{P}^2
$$
la courbe plane projective définie par l'annulation de $F$.
Posons $x = X_1/X_0$ et $y = X_2/X_0$, ainsi que
$f(x, y) = X_0^{-d}F(X_0, X_1, X_2) = F(1, x, y)$.
Notons $K$ le corps des fractions de l'anneau intègre $k[x, y]/(f)$.
Soit $C$ la courbe abstraite correspondant à $K$.
Rappelons (d'après le cours) qu'il existe un morphisme surjectif $C \to D$
qui est bijectif au-dessus du lieu non singulier de $D$ et un
isomorphisme si $D$ est non singulière.
Posons $f_x = \partial f/\partial x$ et $f_y = \partial f/\partial y$.
Enfin, notons $\omega = \text{d}x/f_y = - \text{d}y/f_x$ cet
élément de $\Omega_{K/k}$ étudié dans le cours.
Notons $K_C$ le diviseur des zéros et des pôles de $\omega$.
\end{situation}

\begin{exercise}
\label{exercise-node-in-the-plane}
Dans la situation \ref{situation-curve-in-the-plane}, supposons
$d \geq 3$ et que la courbe $D$ ait exactement un point singulier,
à savoir $P = (1 : 0 : 0)$. Supposons en outre que le développement
$$
f(x, y) = xy + h.o.t
$$
soit valable au voisinage de $P = (0, 0)$. Alors $C$ a deux points $v$ et $w$ au-dessus de
$P$, caractérisés par
$$
v(x) = 1, v(y) > 1
\quad\text{et}\quad
w(x) > 1, w(y) = 1
$$
\begin{enumerate}
\item Montrer que l'élément
$\omega = \text{d}x/f_y = - \text{d}y/f_x$ de $\Omega_{K/k}$
a un pôle simple en $v$ et en $w$. (Le comportement de
$\omega$ aux points non singuliers est celui étudié dans le cours.)
\item Nous avons montré dans le cours que $\omega$ s'annule à l'ordre
$d - 3$ le long du diviseur $X_0 = 0$ ramené sur $C$ par le morphisme
$C \to D$. En combinant ce fait avec l'information de (1), quel est le degré
du diviseur des zéros et des pôles de $\omega$ sur $C$ ?
\item Quel est le genre de la courbe $C$ ?
\end{enumerate}
\end{exercise}

\begin{exercise}
\label{exercise-smooth-plane-curve-linear-system}
Dans la situation \ref{situation-curve-in-the-plane}, supposons $d = 5$
et que la courbe $C = D$ soit non singulière. Nous avons montré dans le cours
que le genre de $C$ vaut $6$ et que le système linéaire $K_C$ est donné par
$$
L(K_C) = \{h\omega \mid h \in k[x, y],\ \deg(h) \leq 2\}
$$
où $\deg$ désigne le degré total\footnote{On obtient $\leq 2$ parce que
$d - 3 = 5 - 3 = 2$.}. Soient $P_1, P_2, P_3, P_4, P_5 \in D$
des points deux à deux distincts situés dans l'ouvert affine $X_0 \not = 0$. Notons
$\sum P_i = P_1 + P_2 + P_3 + P_4 + P_5$ le diviseur correspondant sur $C$.
\begin{enumerate}
\item Décrire $L(K_C - \sum P_i)$ en termes de polynômes.
\item Quelles sont les valeurs possibles de $l(\sum P_i)$ ?
\end{enumerate}
\end{exercise}

\begin{exercise}
\label{exercise-rational-curve-high-degree}
Donner un $F$ comme dans la situation \ref{situation-curve-in-the-plane}
avec $d = 100$ tel que le genre de $C$ soit $0$.
\end{exercise}

\begin{exercise}
\label{exercise-high-degree-curve-quadratic-equation}
Soit $k$ un corps algébriquement clos. Soit $K/k$ une extension de corps
finiment engendrée de degré de transcendance $1$.
Soit $C$ la courbe abstraite correspondant à $K$.
Soit $V \subset K$ un $g^r_d$ et soit $\Phi : C \to \mathbf{P}^r$
le morphisme correspondant. Montrer que l'image de $C$
est contenue dans une quadrique\footnote{Une quadrique est une hypersurface de degré $2$,
c'est-à-dire le lieu des zéros dans $\mathbf{P}^r$ d'un polynôme homogène
de degré $2$.} si $V$ est un système linéaire complet
et si $d$ est assez grand par rapport au genre de $C$.
(Bonus : donner une bonne borne pour le degré nécessaire.)
\end{exercise}

\begin{exercise}
\label{exercise-surjective-map-a2-p2}
Notations de la situation \ref{situation-curve-in-the-plane}.
Soit $U \subset \mathbf{P}^2_k$ le sous-schéma ouvert
dont le complémentaire est $D$. Décrire la $k$-algèbre
$A = \mathcal{O}_{\mathbf{P}^2_k}(U)$. Donner une borne supérieure du
nombre de générateurs de $A$ comme $k$-algèbre.
\end{exercise}



\section{Schémas, examen final, printemps 2014}
\label{section-final-exam-spring-2014}

\noindent
Ces questions ont été posées lors de l'examen final d'un cours sur les
schémas, donné au printemps 2014 à l'Université Columbia.

\begin{exercise}[Définitions]
\label{exercise-definitions-spring-2014}
Soit $(X, \mathcal{O}_X)$ un schéma. Donner la définition des
notions en italique.
\begin{enumerate}
\item l'{\it anneau local de $X$ en un point $x$},
\item un faisceau {\it quasi-cohérent} de $\mathcal{O}_X$-modules,
\item un faisceau {\it cohérent} de $\mathcal{O}_X$-modules (supposer
que $X$ est localement noethérien),
\item un {\it ouvert affine} de $X$,
\item un {\it morphisme fini de schémas} $X \to Y$.
\end{enumerate}
\end{exercise}

\begin{exercise}[Théorèmes]
\label{exercise-results-spring-2014}
Énoncer précisément un résultat non trivial étudié dans le cours qui se rapporte
à chacun des sujets suivants.
\begin{enumerate}
\item l'invariance birationnelle des plurigenres des variétés,
\item le fait que la propriété d'un morphisme d'être affine est locale,
\item la topologie d'une fibre schématique d'un morphisme, et
\item le critère valuatif de propreté.
\end{enumerate}
\end{exercise}

\begin{exercise}
\label{exercise-miss-curve}
Soit $X = \mathbf{A}^2_\mathbf{C}$, où $\mathbf{C}$ est le corps
des nombres complexes. Une {\it droite} désignera un
sous-schéma fermé de $X$ défini par une équation linéaire $ax + by + c = 0$ pour
certains $a, b, c \in \mathbf{C}$ tels que $(a, b) \not = (0, 0)$.
Une {\it courbe} désignera un sous-schéma fermé irréductible (donc non vide)
$C \subset X$ de dimension $1$.
Une {\it conique} désignera une courbe définie par une
équation quadratique $ax^2 + bxy + cy^2 + dx + ey + f = 0$
pour certains $a, b, c, d, e, f \in \mathbf{C}$ tels que
$(a, b, c) \not = (0, 0, 0)$.
\begin{enumerate}
\item Trouver une courbe $C$ telle que toute droite ait une intersection non vide avec $C$.
\item Trouver une courbe $C$ telle que toute droite et toute conique aient une
intersection non vide avec $C$.
\item Montrer que, pour toute courbe $C$, il existe une autre courbe
telle que $C \cap C' = \emptyset$.
\end{enumerate}
\end{exercise}

\begin{exercise}
\label{exercise-normal-bundle-exceptional-curve}
Soit $k$ un corps. Soit $b : X \to \mathbf{A}^2_k$ l'éclatement
du plan affine à l'origine. Autrement dit, si
$\mathbf{A}^2_k = \Spec(k[x, y])$, alors
$X = \text{Proj}(\bigoplus_{n \geq 0} \mathfrak m^n)$
où $\mathfrak m = (x, y) \subset k[x, y]$.
Démontrer les assertions suivantes :
\begin{enumerate}
\item la fibre schématique $E$ de $b$ au-dessus de l'origine
est isomorphe à $\mathbf{P}^1_k$,
\item $E$ est un diviseur de Cartier effectif sur $X$,
\item la restriction de $\mathcal{O}_X(-E)$
à $E$ est un faisceau inversible de degré $1$.
\end{enumerate}
(Rappelons que $\mathcal{O}_X(-E)$ est le faisceau d'idéaux de $E$ dans $X$.)
\end{exercise}

\begin{exercise}
\label{exercise-surjective-map-affine-variety-projective-variety}
Soit $k$ un corps. Soit $X$ une variété projective sur $k$.
Montrer qu'il existe une variété affine $U$ sur $k$ et un morphisme surjectif
de variétés $U \to X$.
\end{exercise}

\begin{exercise}
\label{exercise-vandermonde}
Soit $k$ un corps de caractéristique $p > 0$ différente de $2,3$.
Considérons le sous-schéma fermé $X$ de $\mathbf{P}^n_k$ défini par
$$
\sum\nolimits_{i = 0, \ldots, n} X_i = 0,\quad
\sum\nolimits_{i = 0, \ldots, n} X_i^2 = 0,\quad
\sum\nolimits_{i = 0, \ldots, n} X_i^3 = 0
$$
Pour quelles paires $(n, p)$ cette variété est-elle singulière ?
\end{exercise}






\section{Algèbre commutative, examen final, automne 2016}
\label{section-final-exam-fall-2016}

\noindent
Ces questions ont été posées lors de l'examen final d'un cours
d'algèbre commutative, donné à l'automne 2016 à l'Université Columbia.

\begin{exercise}[Définitions]
\label{exercise-definitions-fall-2016}
Soit $R$ un anneau.
Donner la définition des notions en italique.
\begin{enumerate}
\item l'{\it anneau local de $R$ en un idéal premier $\mathfrak p$},
\item un $R$-module {\it de type fini},
\item un $R$-module {\it de présentation finie},
\item $R$ est {\it régulier},
\item $R$ est {\it caténaire},
\item $R$ est {\it de Cohen-Macaulay}.
\end{enumerate}
\end{exercise}

\begin{exercise}[Théorèmes]
\label{exercise-results-fall-2016}
Énoncer précisément un résultat non trivial étudié dans le cours qui se rapporte
à chacun des sujets suivants.
\begin{enumerate}
\item les anneaux réguliers,
\item les idéaux premiers associés des modules de Cohen-Macaulay,
\item la dimension d'un anneau intègre de type fini sur un corps, et
\item le théorème de Chevalley.
\end{enumerate}
\end{exercise}

\begin{exercise}
\label{exercise-finite-injective}
Soit $A \to B$ un homomorphisme d'anneaux tel que
\begin{enumerate}
\item $A$ soit local, d'idéal maximal $\mathfrak m$,
\item $A \to B$ soit fini\footnote{Rappelons que cela signifie que $B$
est de type fini en tant que $A$-module.},
\item $A \to B$ soit injectif (nous considérons $A$ comme un sous-anneau de $B$).
\end{enumerate}
Montrer qu'il existe un idéal premier $\mathfrak q \subset B$ tel que
$\mathfrak m = A \cap \mathfrak q$.
\end{exercise}

\begin{exercise}
\label{exercise-find-components}
Soit $k$ un corps. Soit $R = k[x, y, z, w]$.
Considérons l'idéal $I = (xy, xz, xw)$.
Quelles sont les composantes irréductibles de
$V(I) \subset \Spec(R)$ et quelles sont leurs dimensions ?
\end{exercise}

\begin{exercise}
\label{exercise-no-nonconstant-morphism}
Soit $k$ un corps. Soient $A = k[x, x^{-1}]$ et $B = k[y]$.
Montrer que tout homomorphisme de $k$-algèbres $\varphi : A \to B$ envoie $x$ sur une constante.
\end{exercise}

\begin{exercise}
\label{exercise-regular-over-Z}
Considérons l'anneau $R = \mathbf{Z}[x, y]/(xy - 7)$.
Démontrer que $R$ est régulier.
\end{exercise}

\noindent
Étant donné un anneau local noethérien $(R, \mathfrak m, \kappa)$,
pour $n \geq 0$, posons
$\varphi_R(n) = \dim_\kappa(\mathfrak m^n/\mathfrak m^{n + 1})$.

\begin{exercise}
\label{exercise-hilbert-function-allowed}
Existe-t-il un anneau local noethérien $R$ tel que
$\varphi_R(n) = n + 1$ pour tout $n \geq 0$ ?
\end{exercise}

\begin{exercise}
\label{exercise-hilbert-function-prohibited}
Soit $R$ un anneau local noethérien.
Supposons que $\varphi_R(0) = 1$, $\varphi_R(1) = 3$,
$\varphi_R(2) = 5$. Montrer que $\varphi_R(3) \leq 7$.
\end{exercise}





\section{Schémas, examen final, printemps 2017}
\label{section-final-exam-spring-2017}

\noindent
Ces questions ont été posées lors de l'examen final d'un cours sur les schémas,
donné au printemps 2017 à l'Université Columbia.

\begin{exercise}[Définitions]
\label{exercise-definitions-spring-2017}
Soit $f : X \to Y$ un morphisme de schémas.
Donner une brève définition des notions en italique.
\begin{enumerate}
\item la {\it fibre schématique} de $f$ en $y \in Y$,
\item $f$ est un {\it morphisme fini},
\item un {\it faisceau quasi-cohérent} de $\mathcal{O}_X$-modules,
\item $X$ est une {\it variété},
\item $f$ est un {\it morphisme lisse},
\item $f$ est un {\it morphisme propre}.
\end{enumerate}
\end{exercise}

\begin{exercise}[Théorèmes]
\label{exercise-results-spring-2017}
Énoncer avec précision, mais brièvement, un résultat non trivial étudié dans le cours
en rapport avec chacun des sujets suivants.
\begin{enumerate}
\item l'image directe des faisceaux quasi-cohérents,
\item la cohomologie des faisceaux cohérents sur les variétés projectives,
\item la dualité de Serre pour un schéma projectif sur un corps, et
\item la formule de Riemann-Hurwitz.
\end{enumerate}
\end{exercise}

\begin{exercise}
\label{exercise-compute-degree}
Soit $k$ un corps algébriquement clos.
Soit $\ell > 100$ un nombre premier différent de la
caractéristique de $k$. Soit $X$ le modèle projectif non singulier
de la courbe affine donnée par l'équation
$$
y^\ell = x(x - 1)^3
$$
dans $\mathbf{A}^2_k$. Répondre aux questions suivantes :
\begin{enumerate}
\item Quel est le genre de $X$ ?
\item Donner une borne supérieure pour la
gonalité\footnote{La gonalité est le plus petit degré d'un morphisme non constant
de $X$ vers $\mathbf{P}^1_k$.} de $X$.
\end{enumerate}
\end{exercise}

\begin{exercise}
\label{exercise-count-components}
Soit $k$ un corps algébriquement clos. Soit $X$ un schéma réduit,
projectif sur $k$, dont toutes les composantes irréductibles
sont de même dimension $1$. Soit $\omega_{X/k}$ le
module dualisant relatif. Montrer que, si
$\dim_k H^1(X, \omega_{X/k}) > 1$, alors $X$ est non connexe.
\end{exercise}

\begin{exercise}
\label{exercise-sections-L-L-inverse}
Donner un exemple d'un schéma $X$ et d'un
$\mathcal{O}_X$-module inversible non trivial $\mathcal{L}$ tels que
$H^0(X, \mathcal{L})$ et $H^0(X, \mathcal{L}^{\otimes -1})$
soient tous deux non nuls.
\end{exercise}

\begin{exercise}
\label{exercise-in-product}
Soit $k$ un corps algébriquement clos. Soit $g \geq 3$.
Soient $X$ et $X'$ des courbes projectives lisses sur $k$, de
genres respectifs $g$ et $g + 1$. Soit $Y \subset X \times X'$
une courbe telle que les projections $Y \to X$ et $Y \to X'$
soient non constantes. Démontrer que le modèle projectif non singulier
de $Y$ est de genre $\geq 2g + 1$.
\end{exercise}

\begin{exercise}
\label{exercise-finitely-many}
Soit $k$ un corps fini. Soit $g > 1$. Esquisser une démonstration de ce qui suit :
il n'existe qu'un nombre fini de classes d'isomorphisme de
courbes projectives lisses sur $k$ et de genre $g$.
(Même une simple tentative de réponse sera prise en compte.)
\end{exercise}







\section{Algèbre commutative, examen final, automne 2017}
\label{section-final-exam-fall-2017}

\noindent
Voici les questions posées à l'examen final d'un cours d'algèbre commutative,
à l'automne 2017, à l'Université Columbia.

\begin{exercise}[Définitions]
\label{exercise-definitions-fall-2017}
Donner une brève définition de chacun des concepts en italique.
\begin{enumerate}
\item l'{\it adjoint à gauche} d'un foncteur $F : \mathcal{A} \to \mathcal{B}$,
\item le {\it degré de transcendance} d'une extension de corps $L/K$,
\item une {\it fonction régulière} sur une variété affine classique
$X \subset k^n$,
\item un {\it faisceau} sur un espace topologique,
\item un {\it anneau local}, et
\item un morphisme de schémas $f : X \to Y$ qui est {\it affine}.
\end{enumerate}
\end{exercise}

\begin{exercise}[Théorèmes]
\label{exercise-results-fall-2017}
Énoncer avec précision mais brièvement un résultat non trivial vu en cours
relatif à chacun des points suivants (s'il en existe plusieurs, en choisir
un seul).
\begin{enumerate}
\item lemme de Yoneda,
\item suite de Mayer-Vietoris,
\item dimension et cohomologie,
\item polynôme de Hilbert, et
\item dualité pour l'espace projectif.
\end{enumerate}
\end{exercise}

\begin{exercise}
\label{exercise-compute-dimension}
Soit $k$ un corps algébriquement clos. Considérons la partie fermée
$X$ de $k^5$ muni de la topologie de Zariski et des coordonnées $x_1, x_2, x_3, x_4, x_5$,
définie par les équations
$$
x_1^2 - x_4 = 0,\quad
x_2^5 - x_5 = 0,\quad
x_3^2 + x_3 + x_4 + x_5 = 0
$$
Quelle est la dimension de $X$, et pourquoi ?
\end{exercise}

\begin{exercise}
\label{exercise-can-there-be}
Soit $k$ un corps. Soit $X = \mathbf{P}^1_k$ l'espace projectif
de dimension $1$ sur $k$. Soit $\mathcal{E}$ un
$\mathcal{O}_X$-module localement libre de rang fini. Pour $d \in \mathbf{Z}$, notons
$\mathcal{E}(d) = \mathcal{E} \otimes_{\mathcal{O}_X} \mathcal{O}_X(d)$
le tordu de Serre d'indice $d$ de $\mathcal{E}$ et
$h^i(X, \mathcal{E}(d)) = \dim_k H^i(X, \mathcal{E}(d))$.
\begin{enumerate}
\item Pourquoi n'existe-t-il pas de $\mathcal{E}$ tel que
$h^0(X, \mathcal{E}) = 5$ et $h^0(X, \mathcal{E}(1)) = 4$ ?
\item Pourquoi n'existe-t-il pas de $\mathcal{E}$ tel que
$h^1(X, \mathcal{E}(1)) = 5$ et $h^1(X, \mathcal{E}) = 4$ ?
\item Pour quels $a \in \mathbf{Z}$ peut-il exister un fibré vectoriel
$\mathcal{E}$ sur $X$ tel que
$$
\begin{matrix}
h^0(X, \mathcal{E})\phantom{(1)} = 1 &
h^1(X, \mathcal{E})\phantom{(1)} = 1 \\
h^0(X, \mathcal{E}(1)) = 2 &
h^1(X, \mathcal{E}(1)) = 0 \\
h^0(X, \mathcal{E}(2)) = 4 &
h^1(X, \mathcal{E}(2)) = a
\end{matrix}
$$
\end{enumerate}
Les réponses partielles sont les bienvenues et vivement encouragées.
\end{exercise}

\begin{exercise}
\label{exercise-banana}
Soit $X$ un espace topologique qui est la réunion
$X = Y \cup Z$ de deux parties fermées
$Y$ et $Z$, dont l'intersection est notée $W = Y \cap Z$.
Notons $i : Y \to X$, $j : Z \to X$ et $k : W \to X$ les applications
d'inclusion.
\begin{enumerate}
\item Montrer qu'il existe une suite exacte courte de faisceaux
$$
0 \to \underline{\mathbf{Z}}_X \to
i_*(\underline{\mathbf{Z}}_Y) \oplus
j_*(\underline{\mathbf{Z}}_Z) \to
k_*(\underline{\mathbf{Z}}_W) \to 0
$$
où $\underline{\mathbf{Z}}_X$ désigne le faisceau constant
de valeur $\mathbf{Z}$ sur $X$, et ainsi de suite.
\item Quelle relation peut-on en déduire entre les
groupes de cohomologie de $X$, $Y$, $Z$ et $W$ à coefficients dans $\mathbf{Z}$ ?
\end{enumerate}
\end{exercise}

\begin{exercise}
\label{exercise-cohomology-infinite-punctured}
Soit $k$ un corps. Soit $A = k[x_1, x_2, x_3, \ldots]$ l'anneau de polynômes
en une infinité d'indéterminées.
Notons $\mathfrak m$ l'idéal maximal de $A$ engendré par
toutes les indéterminées. Soient $X = \Spec(A)$ et $U = X \setminus \{\mathfrak m\}$.
\begin{enumerate}
\item Montrer que $H^1(U, \mathcal{O}_U) = 0$. Indication : calcul de
cohomologie de {\v C}ech.
\item Quelle valeur conjecturez-vous pour $H^i(U, \mathcal{O}_U)$ lorsque $i \geq 1$ ?
\end{enumerate}
\end{exercise}

\begin{exercise}
\label{exercise-principal}
Soit $A$ un anneau local. Soit $a \in A$ un non diviseur de zéro.
Soient $I, J \subset A$ des idéaux tels que $IJ = (a)$. Montrer que
l'idéal $I$ est principal, c'est-à-dire engendré par un seul élément
(qui se révélera être un non diviseur de zéro).
\end{exercise}



\section{Schémas, examen final, printemps 2018}
\label{section-final-exam-spring-2018}

\noindent
Voici les questions posées à l'examen final d'un cours sur les schémas,
au printemps 2018, à l'Université Columbia.

\begin{exercise}[Définitions]
\label{exercise-definitions-spring-2018}
Donner une brève définition de chacun des concepts en italique.
Soit $k$ un corps algébriquement clos.
Soit $X$ une courbe projective sur $k$.
\begin{enumerate}
\item une algèbre {\it lisse} sur $k$,
\item le {\it degré} d'un $\mathcal{O}_X$-module inversible sur $X$,
\item le {\it genre} de $X$,
\item le {\it groupe des classes de diviseurs de Weil} de $X$,
\item $X$ est {\it hyperelliptique}, et
\item le {\it nombre d'intersection}
de deux courbes sur une surface projective lisse sur $k$.
\end{enumerate}
\end{exercise}

\begin{exercise}[Théorèmes]
\label{exercise-results-spring-2018}
Énoncer avec précision mais brièvement un résultat non trivial vu en cours
relatif à chacun des points suivants (s'il en existe plusieurs, en choisir
un seul).
\begin{enumerate}
\item théorème de Riemann-Hurwitz,
\item théorème de Clifford,
\item factorisation des morphismes entre surfaces projectives lisses,
\item théorème de l'indice de Hodge, et
\item hypothèse de Riemann pour les courbes sur les corps finis.
\end{enumerate}
\end{exercise}

\begin{exercise}
\label{exercise-hyperelliptic}
Soit $k$ un corps algébriquement clos. Soit $X \subset \mathbf{P}^3_k$
une courbe lisse de degré $d$ et de genre $\geq 2$. Supposons que $X$
ne soit contenue dans aucun plan et qu'il existe
une droite $\ell$ de $\mathbf{P}^3_k$ rencontrant $X$ en $d - 2$ points.
Montrer que $X$ est hyperelliptique.
\end{exercise}

\begin{exercise}
\label{exercise-singular}
Soit $k$ un corps algébriquement clos. Soit $X$ une courbe projective
ayant pour points singuliers deux à deux distincts $p_1, \ldots, p_n$.
Expliquer pourquoi le genre de la normalisée de $X$ est au plus
$-n + \dim_k H^1(X, \mathcal{O}_X)$.
\end{exercise}

\begin{exercise}
\label{exercise-how-many-blowups}
Soit $k$ un corps. Soit $X = \Spec(k[x, y])$ l'espace affine de dimension $2$.
Soit
$$
I = (x^3, x^2y, xy^2, y^3) \subset k[x, y].
$$
Soit $Y \subset X$ le sous-schéma fermé correspondant à $I$.
Soit $b : X' \to X$ l'éclatement de l'idéal $(x, y)$, c'est-à-dire
l'éclatement de l'espace affine en l'origine.
\begin{enumerate}
\item Montrer que l'image réciproque schématique $b^{-1}Y \subset X'$
est un diviseur de Cartier effectif.
\item Donner un exemple d'idéal $J \subset k[x, y]$
tel que $I \subset J \subset (x, y)$ et que, si $Z \subset X$
est le sous-schéma fermé correspondant à $J$, alors l'image réciproque
schématique $b^{-1}Z$ ne soit pas un diviseur de Cartier effectif.
\end{enumerate}
\end{exercise}

\begin{exercise}
\label{exercise-rational-surface}
Soit $k$ un corps algébriquement clos. Considérons les types
de surfaces suivants :
\begin{enumerate}
\item $S = C_1 \times C_2$ où $C_1$ et $C_2$ sont des courbes projectives lisses,
\item $S = C_1 \times C_2$ où $C_1$ et $C_2$ sont des courbes projectives lisses
et le genre de $C_1$ est $> 0$,
\item $S \subset \mathbf{P}^3_k$ est une hypersurface de degré $4$, et
\item $S \subset \mathbf{P}^3_k$ est une hypersurface lisse de degré $4$.
\end{enumerate}
Pour chaque type, indiquer brièvement si la classe des surfaces de
ce type contient ou non des surfaces rationnelles, et pourquoi.
\end{exercise}

\begin{exercise}
\label{exercise-self-square-line}
Soit $k$ un corps algébriquement clos. Soit $S \subset \mathbf{P}^3_k$
une hypersurface lisse de degré $d$. Supposons que $S$ contienne une
droite $\ell$. Quelle est l'auto-intersection de $\ell$ considérée comme diviseur sur $S$ ?
\end{exercise}






\section{Algèbre commutative, examen final, automne 2019}
\label{section-final-exam-fall-2019}

\noindent
Voici les questions posées à l'examen final d'un cours d'algèbre commutative,
à l'automne 2019, à l'Université Columbia.

\begin{exercise}[Définitions]
\label{exercise-definitions-fall-2019}
Donner une brève définition de chacun des concepts en italique.
\begin{enumerate}
\item une {\it partie constructible} d'un espace topologique noethérien,
\item la {\it localisation} d'un $R$-module $M$ en un idéal premier $\mathfrak p$,
\item la {\it longueur} d'un module sur un anneau local noethérien
$(A, \mathfrak m, \kappa)$,
\item un {\it module projectif} sur un anneau $R$, et
\item un {\it module de Cohen-Macaulay} sur un
anneau local noethérien $(A, \mathfrak m, \kappa)$.
\end{enumerate}
\end{exercise}

\begin{exercise}[Théorèmes]
\label{exercise-results-fall-2019}
Énoncer avec précision mais brièvement un résultat non trivial vu en cours
relatif à chacun des points suivants (s'il en existe plusieurs, en choisir
un seul).
\begin{enumerate}
\item images de parties constructibles,
\item théorème des zéros de Hilbert,
\item dimension des algèbres de type fini sur un corps,
\item lemme de normalisation de Noether, et
\item anneaux locaux réguliers.
\end{enumerate}
\end{exercise}

\noindent
Pour un anneau $R$ et un idéal $I \subset R$, rappelons que $V(I)$
désigne l'ensemble des $\mathfrak p \in \Spec(R)$ tels que $I \subset \mathfrak p$.

\begin{exercise}[Construction d'idéaux premiers]
\label{exercise-infinitely-many-primes}
Construire une infinité d'idéaux premiers distincts
$\mathfrak p \subset \mathbf{C}[x, y]$
tels que $V(\mathfrak p)$ contienne $(x, y)$ et $(x - 1, y - 1)$.
\end{exercise}

\begin{exercise}[Aucun idéal premier]
\label{exercise-no-prime}
Soit $R = \mathbf{C}[x, y, z]/(xy)$. Justifier brièvement qu'il n'existe pas
d'idéal premier $\mathfrak p \subset R$ tel que
$V(\mathfrak p)$ contienne $(x, y - 1, z - 5)$ et $(x - 1, y, z - 7)$.
\end{exercise}

\begin{exercise}[Frobenius]
\label{exercise-frobenius}
Soit $p$ un nombre premier (on pourra supposer $p = 2$ pour simplifier les formules).
Soit $R$ un anneau tel que $p = 0$ dans $R$.
\begin{enumerate}
\item Montrer que l'application $F : R \to R$, $x \mapsto x^p$ est un homomorphisme d'anneaux.
\item Montrer que $\Spec(F) : \Spec(R) \to \Spec(R)$ est l'identité.
\end{enumerate}
\end{exercise}

\noindent
Rappelons qu'une {\it spécialisation} $x \leadsto y$ de points d'un espace topologique
signifie simplement que $y$ appartient à l'adhérence de $x$. Nous disons que $x \leadsto y$ est une
{\it spécialisation immédiate} s'il n'existe aucun $z$ distinct
de $x$ et de $y$ tel que $x \leadsto z$ et $z \leadsto y$.

\begin{exercise}[Dimension]
\label{exercise-dimension}
Supposons que nous ayons un espace topologique sobre $X$ contenant $5$ points distincts
$x, y, z, u, v$ et ayant les spécialisations suivantes :
$$
\xymatrix{
x \ar[d] \ar[r] & u & v \ar[l] \ar[dl] \\
y \ar[r] & z
}
$$
Quelle est la dimension minimale que peut avoir un tel $X$ ? Si $X$ est le spectre
d'une algèbre de type fini sur un corps et si $x \leadsto u$ est une
spécialisation immédiate, que peut-on dire de la
spécialisation $v \leadsto z$ ?
\end{exercise}

\begin{exercise}[Calcul de Tor]
\label{exercise-tor-computation}
Soit $R = \mathbf{C}[x, y, z]$. Soient $M = R/(x, z)$ et $N = R/(y, z)$.
Pour quels $i \in \mathbf{Z}$ le module $\text{Tor}_i^R(M, N)$ est-il non nul ?
\end{exercise}

\begin{exercise}
\label{exercise-depth-goes-up}
Soit $A \to B$ un homomorphisme local plat d'anneaux locaux noethériens.
Montrer que, si $A$ est de profondeur $k$, alors $B$ est de profondeur au moins $k$.
\end{exercise}









\section{Géométrie algébrique, examen final, printemps 2020}
\label{section-final-exam-spring-2020}

\noindent
Voici les questions posées à l'examen final d'un cours de géométrie algébrique,
au printemps 2020, à l'Université Columbia.

\begin{exercise}[Définitions]
\label{exercise-definitions-spring-2020}
Donner une brève définition de chacun des concepts en italique.
\begin{enumerate}
\item un {\it schéma},
\item un {\it morphisme de schémas},
\item un {\it module quasi-cohérent} sur un schéma,
\item une {\it variété} sur un corps $k$,
\item une {\it courbe} sur un corps $k$,
\item un {\it morphisme fini} de schémas,
\item la {\it cohomologie} d'un faisceau de groupes abéliens $\mathcal{F}$
sur un espace topologique $X$,
\item un {\it faisceau dualisant} sur un schéma $X$ de dimension $d$
propre sur un corps $k$, et
\item une {\it application rationnelle} d'une variété $X$ dans une variété $Y$.
\end{enumerate}
\end{exercise}

\begin{exercise}[Théorèmes]
\label{exercise-results-spring-2020}
Énoncer avec précision mais brièvement un résultat non trivial vu en cours
relatif à chacun des points suivants (s'il en existe plusieurs, en choisir
un seul).
\begin{enumerate}
\item cohomologie des faisceaux abéliens sur un espace topologique noethérien
$X$ de dimension $d$,
\item faisceau des différentielles $\Omega^1_{X/k}$ d'une variété lisse
sur un corps $k$,
\item faisceau dualisant $\omega_X$ d'une variété projective lisse $X$
sur le corps $k$,
\item une courbe propre lisse
de genre $0$ sur un corps algébriquement clos $k$, et
\item le genre d'une courbe plane de degré $d$.
\end{enumerate}
\end{exercise}

\begin{exercise}
\label{exercise-coh-P1-infty-doubled}
Soit $k$ un corps. Soit $X$ un schéma sur $k$. Supposons que $X = X_1 \cup X_2$
soit un recouvrement ouvert, $X_1$ et $X_2$ étant tous deux isomorphes à $\mathbf{P}^1_k$,
et $X_1 \cap X_2$ étant isomorphe à $\mathbf{A}^1_k$. (Un tel schéma existe :
par exemple, on peut prendre $\mathbf{P}^1_k$ avec le point $\infty$ dédoublé.)
Montrer que $\dim_k H^1(X, \mathcal{O}_X)$ est infinie.
\end{exercise}

\begin{exercise}
\label{exercise-no-morphism}
Soit $k$ un corps algébriquement clos. Soit $Y$ une courbe projective lisse
de genre $10$. Trouver une bonne borne inférieure pour le
genre d'une courbe projective lisse
$X$ telle qu'il existe un morphisme non constant $f : X \to Y$
qui ne soit pas un isomorphisme.
\end{exercise}

\begin{exercise}
\label{exercise-element-order-n-in-pic}
Soit $k$ un corps algébriquement clos de caractéristique $0$.
Soit
$$
X : T_0^d + T_1^d - T_2^d = 0 \subset \mathbf{P}^2_k
$$
la courbe de Fermat de degré $d \geq 3$. Considérons les points fermés
$p = [1 : 0 : 1]$ et $q = [0 : 1 : 1]$ de $X$. Posons $D = [p] - [q]$.
\begin{enumerate}
\item Montrer que $D$ est non trivial dans le groupe des classes de diviseurs de Weil.
\item Montrer que $d D$ est trivial dans le groupe des classes de diviseurs de Weil.
(Indication : essayer de montrer que $d[p]$ et $d[q]$ sont tous deux l'intersection
de $X$ avec une droite du plan.)
\end{enumerate}
\end{exercise}

\begin{exercise}
\label{exercise-number-equations}
Soit $k$ un corps algébriquement clos. Considérons le plongement $2$-uple
$$
\varphi : \mathbf{P}^2 \longrightarrow \mathbf{P}^5
$$
Dans la terminologie et les notations du cours, il s'agit du morphisme
$$
\varphi = \varphi_{\mathcal{O}_{\mathbf{P}^2}(2)} :
\mathbf{P}^2
\longrightarrow
\mathbf{P}(\Gamma(\mathbf{P}^2, \mathcal{O}_{\mathbf{P}^2}(2)))
$$
En coordonnées homogènes, il est donné par
$$
[a_0 : a_1 : a_2] \longmapsto
[a_0^2 : a_0a_1 : a_0a_2 : a_1^2 : a_1a_2 : a_2^2]
$$
C'est une immersion fermée (on pourra simplement l'admettre).
Soit $I \subset k[T_0, \ldots, T_5]$ l'idéal homogène
de $\varphi(\mathbf{P}^2)$, c'est-à-dire que les éléments de la
composante homogène $I_d$ sont les
polynômes homogènes $F(T_0, \ldots, T_5)$ de degré $d$ dont la
restriction au sous-schéma fermé $\varphi(\mathbf{P}^2)$ est nulle.
Calculer $\dim_k I_d$ en fonction de $d$.
\end{exercise}

\begin{exercise}
\label{exercise-dualizing-sheaf-rank-1}
Soit $k$ un corps algébriquement clos. Soit $X$ un schéma propre de
dimension $d$ sur $k$, muni d'un module dualisant $\omega_X$.
On dispose des informations suivantes :
\begin{enumerate}
\item $\text{Ext}^i_X(\mathcal{F}, \omega_X) \times
H^{d - i}(X, \mathcal{F}) \to H^d(X, \omega_X) \xrightarrow{t} k$
est non dégénéré pour tout $i$ et tout $\mathcal{O}_X$-module cohérent
$\mathcal{F}$, et
\item $\omega_X$ est localement libre de rang fini, disons $r$.
\end{enumerate}
Montrer que $r = 1$. (Indication : examiner ce qui se passe en prenant pour $\mathcal{F}$
un module convenable dont le support est un point fermé.)
\end{exercise}







\section{Algèbre commutative, examen final, automne 2021}
\label{section-final-exam-fall-2021}

\noindent
Les questions suivantes ont été posées lors de l'examen final d'un cours d'algèbre commutative,
donné à l'automne 2021 à l'Université Columbia.

\begin{exercise}[Définitions]
\label{exercise-definitions-fall-2021}
Donner des définitions succinctes des notions en italique.
\begin{enumerate}
\item une {\it partie multiplicative} d'un anneau $A$,
\item un {\it anneau artinien} $A$,
\item le {\it spectre d'un anneau} $A$ en tant qu'espace topologique,
\item un {\it homomorphisme d'anneaux plat} $A \to B$,
\item la {\it hauteur} d'un idéal premier $\mathfrak p$ de $A$, et
\item les foncteurs {\it $\text{Tor}^A_i(-, -)$} sur un anneau $A$.
\end{enumerate}
\end{exercise}

\begin{exercise}[Théorèmes]
\label{exercise-results-fall-2021}
Énoncer avec précision mais brièvement un résultat non trivial vu en cours
se rapportant à chacun des points suivants (s'il y en a plusieurs, n'en choisir
qu'un seul).
\begin{enumerate}
\item les anneaux artiniens,
\item la platitude et les idéaux premiers,
\item les longueurs des $A/\mathfrak m^n$ lorsque $(A, \mathfrak m)$ est un anneau local noethérien,
\item la formule de dimension pour les anneaux noethériens universellement caténaires,
\item la complétion d'un anneau local noethérien, et
\item la dualité de Matlis pour les anneaux locaux artiniens.
\end{enumerate}
\end{exercise}

\begin{exercise}[Unités]
\label{exercise-simple-units}
Quelle est la structure du groupe des unités de $\mathbf{Z}[x, 1/x]$
en tant que groupe abélien ? Aucune justification n'est nécessaire.
\end{exercise}

\begin{exercise}[Idéaux]
\label{exercise-ideals}
Soit $A = \mathbf{F}_2[x, y]/(x^2, xy, y^2)$ et notons
$\overline{x}$ et $\overline{y}$ les images de $x$ et $y$ dans $A$.
Énumérer les idéaux de $A$. Aucune justification n'est nécessaire.
\end{exercise}

\begin{exercise}[Tor et Ext]
\label{exercise-compute-tor-ext}
Soit $(A, \mathfrak m, \kappa)$ un anneau local noethérien.
Posons $\varphi(n) = \dim_\kappa \mathfrak m^n/\mathfrak m^{n + 1}$.
\begin{enumerate}
\item Montrer que $\text{Tor}_1^A(A/\mathfrak m^n, \kappa)$
est de dimension $\varphi(n)$ comme espace vectoriel sur $\kappa$.
\item Montrer que $\text{Ext}^1_A(A/\mathfrak m^n, \kappa)$
est de dimension $\varphi(n)$ comme espace vectoriel sur $\kappa$.
\end{enumerate}
\end{exercise}

\begin{exercise}[Deux vecteurs]
\label{exercise-two-vectors}
Soit $A = \mathbf{Z}[a_1, a_2, a_3, b_1, b_2, b3]$. Posons
$a = (a_1, a_2, a_3)$ et $b = (b_1, b_2, b_3)$ dans $A^{\oplus 3}$.
Considérons l'ensemble
$$
Z = \{\mathfrak p \in \Spec(A) \mid a, b
\text{ ont des images linéairement dépendantes dans }
\kappa(\mathfrak p)^{\oplus 3}\}
$$
\begin{enumerate}
\item Prouver que $Z$ est une partie fermée de $\Spec(A)$.
\item Quelle est la dimension $\dim(Z)$ de $Z$ ?
\item Que deviendrait $\dim(Z)$ si l'on remplaçait $\mathbf{Z}$ par un corps ?
\end{enumerate}
\end{exercise}

\begin{exercise}[Injectifs]
\label{exercise-injective-artinian-local}
Soit $(A, \mathfrak m, \kappa)$ un anneau local artinien.
Supposons que $A$ soit injectif comme $A$-module. Montrer que
$\Hom_A(\kappa, A)$ est de dimension $1$ comme
espace vectoriel sur $\kappa$.
\end{exercise}



\section{Géométrie algébrique, examen final, printemps 2022}
\label{section-final-exam-spring-2022}

\noindent
Les questions suivantes ont été posées lors de l'examen final d'un cours de géométrie algébrique,
donné au printemps 2022 à l'Université Columbia.

\begin{exercise}[Définitions]
\label{exercise-definitions-spring-2022}
Donner des définitions succinctes des notions en italique.
\begin{enumerate}
\item un {\it schéma},
\item un {\it module quasi-cohérent} sur un schéma $X$,
\item un morphisme {\it plat} de schémas $X \to Y$,
\item un morphisme {\it fini} de schémas $X \to Y$,
\item un {\it schéma en groupes} $G$ au-dessus d'un schéma de base $S$,
\item une {\it famille de variétés} au-dessus d'un schéma de base $S$,
\item le {\it degré} d'un point fermé $x$ d'une variété $X$
sur le corps $k$,
\item la {\it hauteur logarithmique} usuelle d'un point
$p = (a_0 : \ldots : a_n)$ de $\mathbf{P}^n(\mathbf{Q})$, et
\item un {\it corps $C_i$}.
\end{enumerate}
\end{exercise}

\begin{exercise}[Théorèmes]
\label{exercise-results-spring-2022}
Énoncer avec précision mais brièvement un résultat non trivial vu en cours
se rapportant à chacun des points suivants (s'il y en a plusieurs, n'en choisir
qu'un seul).
\begin{enumerate}
\item les morphismes d'un schéma $X$ vers le schéma affine $\Spec(A)$,
\item la cohomologie d'un module quasi-cohérent $\mathcal{F}$ sur
un schéma affine $X$,
\item le groupe de Picard de $\mathbf{P}^1_k$, où $k$ est un corps,
\item les dimensions des fibres d'un morphisme propre et plat
$X \to S$ lorsque $S$ est noethérien,
\item les modules $\mathbf{G}_m$-équivariants sur un schéma $S$, et
\item le théorème de Bézout sur les intersections (on pourra se limiter à
un cas particulier).
\end{enumerate}
\end{exercise}

\begin{exercise}[Hypersurfaces cubiques]
\label{exercise-cubics}
Soit $F \in \mathbf{C}[T_0, \ldots, T_n]$ homogène de degré $3$.
Étant donnés $3$ vecteurs $x, y, z \in \mathbf{C}^{n + 1}$, considérons la condition
$$
(*)\quad
F(\lambda x + \mu y + \nu z) = 0 \text{ dans } \mathbf{C}[\lambda, \mu, \nu]
$$
\begin{enumerate}
\item Quelle est la dimension de l'espace de tous les choix de $x, y, z$ ?
\item Combien d'équations portant sur les coordonnées de $x$, $y$ et $z$
la condition (*) représente-t-elle ?
\item Quelle est la dimension attendue de l'espace de tous les triplets
$x, y, z$ vérifiant (*) ?
\item Quelle est la dimension de l'espace de tous les triplets tels que
$x, y, z$ soient linéairement dépendants ?
\item En déduire que, sur une hypersurface de degré $3$ dans $\mathbf{P}^n$,
on s'attend à trouver un sous-espace linéaire de dimension $2$ pourvu que
$n \geq a$, où il vous appartient de déterminer $a$.
\end{enumerate}
\end{exercise}

\begin{exercise}[Hauteurs]
\label{exercise-heights}
Soit $K$ un corps. Soit $h_n : \mathbf{P}^n(K) \to \mathbf{R}$, $n \geq 0$,
une famille de fonctions satisfaisant les $2$ axiomes étudiés
en cours. Soit $X$ une variété projective
sur $K$. Soit $\mathcal{L}$ un $\mathcal{O}_X$-module inversible
et rappelons que nous lui avons associé en cours une
fonction hauteur $h_\mathcal{L} : X(K) \to \mathbf{R}$.
Soit $\alpha : X \to X$ un automorphisme de $X$ au-dessus de $K$.
\begin{enumerate}
\item Montrer que l'application $P \mapsto h_\mathcal{L}(\alpha(P))$
diffère de la fonction $h_{\alpha^*\mathcal{L}}$
par une fonction bornée. (Indication : rappeler que s'il existe un morphisme
$\varphi : X \to \mathbf{P}^n$ tel que
$\mathcal{L} = \varphi^*\mathcal{O}_{\mathbf{P}^n}(1)$, alors, par construction,
$h_\mathcal{L}(P) = h_n(\varphi(P))$ et jouer avec cette identité.
Dans le cas général, écrire $\mathcal{L}$ comme différence de deux modules de ce type.)
\item Supposons que $h_\mathcal{L}(P) - h_\mathcal{L}(\alpha(P))$
ne soit pas bornée sur $X(K)$. Montrer que
$h_\mathcal{N}$, où
$\mathcal{N} = \mathcal{L} \otimes \alpha^*\mathcal{L}^{\otimes -1}$,
n'est pas bornée sur $X(K)$.
\item Supposons que $X$ soit une courbe elliptique et que
$\mathcal{L}$ soit un module inversible ample et symétrique sur $X$
tel que $h_\mathcal{L}$ ne soit pas bornée sur $X(K)$. Montrer
qu'il existe un module inversible $\mathcal{N}$ de degré $0$
tel que $h_\mathcal{N}$ ne soit pas bornée.
(Indications : rappeler que $X$ est une variété abélienne de dimension $1$.
Ainsi, $h_\mathcal{L}$ est quadratique à une constante près, d'après les résultats
vus en cours. Choisir un point convenable $P_0 \in X(K)$.
Soit $\alpha : X \to X$ la translation par $P_0$. Considérer
$P \mapsto h_\mathcal{L}(P) - h_\mathcal{L}(P + P_0)$.
Appliquer les résultats démontrés ci-dessus.)
\end{enumerate}
\end{exercise}

\begin{exercise}[Monomorphismes]
\label{exercise-monomorphisms}
Soit $f : X \to Y$ un monomorphisme dans la catégorie des schémas :
pour toute paire de morphismes $a, b : T \to X$ de schémas, si
$f \circ a = f \circ b$, alors $a = b$. Montrer que $f$ est injectif
sur les points. Votre argument montre-t-il autre chose ?
\end{exercise}

\begin{exercise}[Points fixes]
\label{exercise-fix-points}
Soit $k$ un corps algébriquement clos.
\begin{enumerate}
\item Si $G = \mathbf{G}_{m, k}$, montrer que, si $G$ agit sur une
variété projective $X$ sur $k$, alors l'action admet un point fixe, c'est-à-dire
montrer qu'il existe un point $x \in X(k)$ tel que $a(g, x) = x$ pour
tout $g \in G(k)$.
\item Même question pour $G = (\mathbf{G}_{m, k})^n$, produit de
$n \geq 1$ exemplaires du groupe multiplicatif.
\item Donner un exemple d'une action d'un schéma en groupes connexe $G$
sur une variété projective lisse $X$ qui n'admet aucun point fixe.
\end{enumerate}
\end{exercise}



\section{Géométrie algébrique, examen final, printemps 2025}
\label{section-final-exam-spring-2025}

\noindent
Les questions suivantes ont été posées lors de l'examen final d'un cours de géométrie algébrique,
donné au printemps 2025 à l'Université Columbia.

\begin{exercise}[Définitions]
\label{exercise-definitions-spring-2025}
Donner des définitions succinctes des notions en italique.
\begin{enumerate}
\item un {\it schéma} $X$,
\item un {\it module quasi-cohérent} sur un schéma $X$,
\item le {\it groupe de Picard} d'un schéma $X$,
\item étant donnés un morphisme $f : X \to Y$ de schémas et un
$\mathcal{O}_X$-module $\mathcal{F}$, son {\it image directe} $f_*\mathcal{F}$,
\item une {\it immersion fermée} de schémas, et
\item une {\it variété} sur un corps donné $k$.
\end{enumerate}
\end{exercise}

\begin{exercise}[Théorèmes]
\label{exercise-results-spring-2025}
Énoncer avec précision et concision un résultat non trivial vu en cours
se rapportant à chacun des points suivants (s'il y en a plusieurs, n'en choisir qu'un seul).
\begin{enumerate}
\item l'espace topologique d'un schéma $X$,
\item un critère de représentabilité pour les foncteurs sur la catégorie des schémas,
\item le foncteur de Picard d'une courbe projective lisse $X$ sur un
corps algébriquement clos $k$,
\item la torsion dans le groupe de Picard d'une courbe projective lisse $X$ sur un
corps algébriquement clos $k$,
\item un théorème sur le groupe de Picard d'un produit $X \times Y \times Z$
de variétés projectives lisses $X$, $Y$, $Z$ sur un corps algébriquement
clos $k$.
\end{enumerate}
\end{exercise}

\begin{exercise}
\label{exercise-affine-line-infinite-nr-points}
Soit $k$ un corps. Expliquer pourquoi $\Spec(k[x])$ a une infinité de points.
\end{exercise}

\begin{exercise}
\label{exercise-hypersurface-variety}
Soit $k$ un corps et soient $x_1, \ldots, x_n$ des indéterminées.
Soit $f \in k[x_1, \ldots, x_n]$ un polynôme non constant.
Soit $X = \Spec(k[x_1, \ldots, x_n]/(f))$. Quelle propriété
$f$ doit-il posséder pour que $X$ soit une variété sur $k$ ?
\end{exercise}

\begin{exercise}
\label{exercise-find-singular-point}
Trouver un point singulier (il n'est pas demandé de les trouver tous)
du spectre de $\mathbf{C}[x, y]/(x^5 - 5xy^4 + 20y - 16)$
considéré comme une variété sur $\mathbf{C}$.
\end{exercise}

\begin{exercise}
\label{exercise-kernel-restriction-pic-curve}
Soit $X$ une courbe projective lisse sur un corps algébriquement clos
$k$. Soit $x \in X(k)$ un point fermé de $X$. Soit
$U = X \setminus \{x\}$. Que peut-on dire du noyau de
l'homomorphisme de restriction $\Pic(X) \to \Pic(U)$ ?
\end{exercise}

\begin{exercise}
\label{exercise-sections-degree-2g-minus-4}
Soit $X$ une courbe projective lisse sur un corps algébriquement clos
$k$ de genre $g \geq 100$. Quelles valeurs
$h^0(\mathcal{L}) = \dim_k H^0(X, \mathcal{L})$ peut-il prendre lorsque
$\mathcal{L}$ est un $\mathcal{O}_X$-module inversible de degré $2g - 4$
sur $X$ ? Répondre de manière aussi complète que possible.
\end{exercise}









\begin{multicols}{2}[\section{Autres chapitres}]
\noindent
Préliminaires
\begin{enumerate}
\item \hyperref[introduction-section-phantom]{Introduction}
\item \hyperref[conventions-section-phantom]{Conventions}
\item \hyperref[sets-section-phantom]{Théorie des ensembles}
\item \hyperref[categories-section-phantom]{Catégories}
\item \hyperref[topology-section-phantom]{Topologie}
\item \hyperref[sheaves-section-phantom]{Faisceaux sur les espaces}
\item \hyperref[sites-section-phantom]{Sites et faisceaux}
\item \hyperref[stacks-section-phantom]{Champs}
\item \hyperref[fields-section-phantom]{Corps}
\item \hyperref[algebra-section-phantom]{Algèbre commutative}
\item \hyperref[brauer-section-phantom]{Groupes de Brauer}
\item \hyperref[homology-section-phantom]{Algèbre homologique}
\item \hyperref[derived-section-phantom]{Catégories dérivées}
\item \hyperref[simplicial-section-phantom]{Méthodes simpliciales}
\item \hyperref[more-algebra-section-phantom]{Compléments d'algèbre}
\item \hyperref[smoothing-section-phantom]{Lissage des morphismes d'anneaux}
\item \hyperref[modules-section-phantom]{Faisceaux de modules}
\item \hyperref[sites-modules-section-phantom]{Modules sur les sites}
\item \hyperref[injectives-section-phantom]{Objets injectifs}
\item \hyperref[cohomology-section-phantom]{Cohomologie des faisceaux}
\item \hyperref[sites-cohomology-section-phantom]{Cohomologie sur les sites}
\item \hyperref[dga-section-phantom]{Algèbre différentielle graduée}
\item \hyperref[dpa-section-phantom]{Algèbre à puissances divisées}
\item \hyperref[sdga-section-phantom]{Faisceaux différentiels gradués}
\item \hyperref[hypercovering-section-phantom]{Hyperrecouvrements}
\end{enumerate}
Schémas
\begin{enumerate}
\setcounter{enumi}{25}
\item \hyperref[schemes-section-phantom]{Schémas}
\item \hyperref[constructions-section-phantom]{Constructions de schémas}
\item \hyperref[properties-section-phantom]{Propriétés des schémas}
\item \hyperref[morphisms-section-phantom]{Morphismes de schémas}
\item \hyperref[coherent-section-phantom]{Cohomologie des schémas}
\item \hyperref[divisors-section-phantom]{Diviseurs}
\item \hyperref[limits-section-phantom]{Limites de schémas}
\item \hyperref[varieties-section-phantom]{Variétés}
\item \hyperref[topologies-section-phantom]{Topologies sur les schémas}
\item \hyperref[descent-section-phantom]{Descente}
\item \hyperref[perfect-section-phantom]{Catégories dérivées de schémas}
\item \hyperref[more-morphisms-section-phantom]{Compléments sur les morphismes}
\item \hyperref[flat-section-phantom]{Compléments sur la platitude}
\item \hyperref[groupoids-section-phantom]{Schémas en groupoïdes}
\item \hyperref[more-groupoids-section-phantom]{Compléments sur les schémas en groupoïdes}
\item \hyperref[etale-section-phantom]{Morphismes étales de schémas}
\end{enumerate}
Thèmes de théorie des schémas
\begin{enumerate}
\setcounter{enumi}{41}
\item \hyperref[chow-section-phantom]{Homologie de Chow}
\item \hyperref[intersection-section-phantom]{Théorie de l'intersection}
\item \hyperref[pic-section-phantom]{Schémas de Picard des courbes}
\item \hyperref[weil-section-phantom]{Théories cohomologiques de Weil}
\item \hyperref[adequate-section-phantom]{Modules adéquats}
\item \hyperref[dualizing-section-phantom]{Complexes dualisants}
\item \hyperref[duality-section-phantom]{Dualité pour les schémas}
\item \hyperref[discriminant-section-phantom]{Discriminants et différentes}
\item \hyperref[derham-section-phantom]{Cohomologie de de Rham}
\item \hyperref[local-cohomology-section-phantom]{Cohomologie locale}
\item \hyperref[algebraization-section-phantom]{Géométrie algébrique et formelle}
\item \hyperref[curves-section-phantom]{Courbes algébriques}
\item \hyperref[resolve-section-phantom]{Résolution des surfaces}
\item \hyperref[models-section-phantom]{Réduction semi-stable}
\item \hyperref[functors-section-phantom]{Foncteurs et morphismes}
\item \hyperref[equiv-section-phantom]{Catégories dérivées de variétés}
\item \hyperref[pione-section-phantom]{Groupes fondamentaux de schémas}
\item \hyperref[etale-cohomology-section-phantom]{Cohomologie étale}
\item \hyperref[crystalline-section-phantom]{Cohomologie cristalline}
\item \hyperref[proetale-section-phantom]{Cohomologie pro-étale}
\item \hyperref[relative-cycles-section-phantom]{Cycles relatifs}
\item \hyperref[more-etale-section-phantom]{Compléments de cohomologie étale}
\item \hyperref[trace-section-phantom]{La formule des traces}
\end{enumerate}
Espaces algébriques
\begin{enumerate}
\setcounter{enumi}{64}
\item \hyperref[spaces-section-phantom]{Espaces algébriques}
\item \hyperref[spaces-properties-section-phantom]{Propriétés des espaces algébriques}
\item \hyperref[spaces-morphisms-section-phantom]{Morphismes d'espaces algébriques}
\item \hyperref[decent-spaces-section-phantom]{Espaces algébriques décents}
\item \hyperref[spaces-cohomology-section-phantom]{Cohomologie des espaces algébriques}
\item \hyperref[spaces-limits-section-phantom]{Limites d'espaces algébriques}
\item \hyperref[spaces-divisors-section-phantom]{Diviseurs sur les espaces algébriques}
\item \hyperref[spaces-over-fields-section-phantom]{Espaces algébriques sur des corps}
\item \hyperref[spaces-topologies-section-phantom]{Topologies sur les espaces algébriques}
\item \hyperref[spaces-descent-section-phantom]{Descente et espaces algébriques}
\item \hyperref[spaces-perfect-section-phantom]{Catégories dérivées d'espaces}
\item \hyperref[spaces-more-morphisms-section-phantom]{Compléments sur les morphismes d'espaces}
\item \hyperref[spaces-flat-section-phantom]{Platitude sur les espaces algébriques}
\item \hyperref[spaces-groupoids-section-phantom]{Groupoïdes dans les espaces algébriques}
\item \hyperref[spaces-more-groupoids-section-phantom]{Compléments sur les groupoïdes dans les espaces}
\item \hyperref[bootstrap-section-phantom]{Amorçage}
\item \hyperref[spaces-pushouts-section-phantom]{Sommes amalgamées d'espaces algébriques}
\end{enumerate}
Thèmes de géométrie
\begin{enumerate}
\setcounter{enumi}{81}
\item \hyperref[spaces-chow-section-phantom]{Groupes de Chow des espaces}
\item \hyperref[groupoids-quotients-section-phantom]{Quotients de groupoïdes}
\item \hyperref[spaces-more-cohomology-section-phantom]{Compléments sur la cohomologie des espaces}
\item \hyperref[spaces-simplicial-section-phantom]{Espaces simpliciaux}
\item \hyperref[spaces-duality-section-phantom]{Dualité pour les espaces}
\item \hyperref[formal-spaces-section-phantom]{Espaces algébriques formels}
\item \hyperref[restricted-section-phantom]{Algébrisation des espaces formels}
\item \hyperref[spaces-resolve-section-phantom]{Résolution des surfaces revisitée}
\end{enumerate}
Théorie des déformations
\begin{enumerate}
\setcounter{enumi}{89}
\item \hyperref[formal-defos-section-phantom]{Théorie formelle des déformations}
\item \hyperref[defos-section-phantom]{Théorie des déformations}
\item \hyperref[cotangent-section-phantom]{Le complexe cotangent}
\item \hyperref[examples-defos-section-phantom]{Problèmes de déformation}
\end{enumerate}
Champs algébriques
\begin{enumerate}
\setcounter{enumi}{93}
\item \hyperref[algebraic-section-phantom]{Champs algébriques}
\item \hyperref[examples-stacks-section-phantom]{Exemples de champs}
\item \hyperref[stacks-sheaves-section-phantom]{Faisceaux sur les champs algébriques}
\item \hyperref[criteria-section-phantom]{Critères de représentabilité}
\item \hyperref[artin-section-phantom]{Axiomes d'Artin}
\item \hyperref[quot-section-phantom]{Espaces de Quot et de Hilbert}
\item \hyperref[stacks-properties-section-phantom]{Propriétés des champs algébriques}
\item \hyperref[stacks-morphisms-section-phantom]{Morphismes de champs algébriques}
\item \hyperref[stacks-limits-section-phantom]{Limites de champs algébriques}
\item \hyperref[stacks-cohomology-section-phantom]{Cohomologie des champs algébriques}
\item \hyperref[stacks-perfect-section-phantom]{Catégories dérivées de champs}
\item \hyperref[stacks-introduction-section-phantom]{Introduction aux champs algébriques}
\item \hyperref[stacks-more-morphisms-section-phantom]{Compléments sur les morphismes de champs}
\item \hyperref[stacks-geometry-section-phantom]{La géométrie des champs}
\end{enumerate}
Thèmes de théorie des espaces de modules
\begin{enumerate}
\setcounter{enumi}{107}
\item \hyperref[moduli-section-phantom]{Champs de modules}
\item \hyperref[moduli-curves-section-phantom]{Espaces de modules de courbes}
\end{enumerate}
Divers
\begin{enumerate}
\setcounter{enumi}{109}
\item \hyperref[examples-section-phantom]{Exemples}
\item \hyperref[exercises-section-phantom]{Exercices}
\item \hyperref[guide-section-phantom]{Guide bibliographique}
\item \hyperref[desirables-section-phantom]{Desiderata}
\item \hyperref[coding-section-phantom]{Style de programmation}
\item \hyperref[obsolete-section-phantom]{Obsolète}
\item \hyperref[fdl-section-phantom]{Licence de documentation libre GNU}
\item \hyperref[index-section-phantom]{Index généré automatiquement}
\end{enumerate}
\end{multicols}


\bibliography{my}
\bibliographystyle{amsalpha}

\end{document}

